\documentclass{article}
\usepackage[letterpaper,margin=1.0in]{geometry}
\usepackage{hyperref}
\usepackage{enumitem}
%\usepackage{natbib}
\usepackage[utf8]{inputenc} % allow utf-8 input
\usepackage[T1]{fontenc}    % use 8-bit T1 fonts
\usepackage{hyperref}       % hyperlinks
\usepackage{url}            % simple URL typesetting
\usepackage{booktabs}       % professional-quality tables
\usepackage{amsfonts}       % blackboard math symbols
\usepackage{nicefrac}       % compact symbols for 1/2, etc.
\usepackage{microtype}      % microtypography
%\usepackage{lipsum}
%\usepackage{fancyhdr}       % header
\usepackage{graphicx}       % graphics
\usepackage[numbers,sort]{natbib}
\usepackage{multirow}
\usepackage{bbding}
\graphicspath{{media/}}     % organize your images and
\usepackage{color}
\usepackage{algorithmic,algorithm}
\usepackage{nicefrac}
\usepackage{cancel}
\usepackage{graphicx} 
\usepackage{booktabs}
\usepackage{makecell}
\usepackage{multirow}
\usepackage{amssymb,amsmath,amsthm}
\usepackage[table]{xcolor}

%\pagestyle{fancy}
%\thispagestyle{empty}
%\rhead{ \textit{ }} 

\input{math}
\usepackage{thmtools}
\usepackage{thm-restate}

\title{On the Condition Number Dependency in
Bilevel Optimization}


\author{Lesi Chen \qquad 
\qquad Jingzhao Zhang  \\
\vspace{-2mm} \\
\normalsize{IIIS, Tsinghua University
} \\
\normalsize{ \texttt{ chenlc23@mails.tsinghua.edu.cn,  jingzhaoz@mail.tsinghua.edu.cn}} 
}

% \author{Name \thanks{Address; Email}
% }
%\date{}
      
\begin{document}
\maketitle
\begingroup
\begin{NoHyper}
% \renewcommand\thefootnote{*}
% \footnotetext{Equal contributions.}
\end{NoHyper}
\endgroup
\begin{abstract}
Bilevel optimization minimizes an objective function, defined by an upper-level problem whose feasible region is the solution of a lower-level problem. We study the oracle complexity of finding an $\epsilon$-stationary point with first-order methods when the upper-level problem is nonconvex and the lower-level problem is strongly convex. 
Recent works (Ji et al., ICML 2021; Arbel and Mairal, ICLR 2022; Chen el al., JMLR 2025)  achieve a $\tilde \gO(\kappa^4 \epsilon^{-2})$ upper bound that is near-optimal in $\epsilon$. However, the optimal dependency on the condition number $\kappa$ is unknown.
In this work, we establish a new $\Omega(\kappa^2 \epsilon^{-2})$ lower bound and $\tilde \gO(\kappa^{7/2} \epsilon^{-2})$ upper bound for this problem, establishing the first provable gap between bilevel problems and minimax problems in this setup. 
Our lower bounds can be extended to various settings, including high-order smooth functions, stochastic oracles, and convex hyper-objectives: (1) For second-order and arbitrarily smooth problems, we show $\Omega(\kappa_y^{13/4} \epsilon^{-12/7})$ and $\Omega(\kappa^{17/10} \epsilon^{-8/5})$ lower bounds, respectively. 
(2) For convex-strongly-convex problems, we improve the previously best lower bound (Ji and Liang, JMLR 2022) from $\Omega(\kappa /\sqrt{\epsilon})$ to $\Omega(\kappa^{5/4} / \sqrt{\epsilon})$. 
(3) For smooth stochastic problems, we show an $\Omega(\kappa^4 \epsilon^{-4})$ lower bound.  
\end{abstract}

\section{Introduction}
We study the first-order oracle complexity for solving the bilevel optimization
\begin{align} \label{prob:main-BLO}
    \min_{\vx \in \sR^{d_x}} F(\vx) = f(\vx,\vy^*(\vx)) , \quad \vy^*(\vx) = \arg \min_{\vy \in \sR^{d_y}} g(\vx,\vy).
\end{align}
The goal is to minimize the hyper-objective $F(\vx)$, which is implicitly defined by two explicit objectives $f(\vx,\vy)$ and $g(\vx,\vy)$ and the $\arg \min$ operator in $g(\vx,\vy)$ with respect to~$\vy$. 
This formulation originates from the two-player general-sum Stackelberg game \citep{stackelberg1934marktform}, and reflects the sequential nature of the decision process of two players $\vx$ and $\vy$. Bilevel optimization of this form got extensive attention in the machine learning community due to its wide applications, such as meta-learning \citep{rajeswaran2019meta}, hyper-parameter tuning \citep{bao2021stability,franceschi2018bilevel,mackay2019self},  generative adversarial networks \citep{goodfellow2020generative}, and reinforcement learning \citep{konda1999actor,zeng2024two,hong2023two}.

% Since even when the lower-level function $g(\vx,\vy)$ is convex in $\vy$, it is impossible to derive dimension-free upper bounds for any first-order methods \citep[Theorem 3.2]{chen2024finding},

To ensure that Problem (\ref{prob:main-BLO}) is well-defined, we follow \citep{ghadimi2018approximation,ji2021bilevel,hong2023two} to study the standard nonconvex-strongly-convex (NC-SC) setting, where the lower-level function $g(\vx,\vy)$ is strongly convex in $\vy$ while the upper-level function $f(\vx,\vy)$ can be nonconvex.
%This  setting is the most standard setting in the literature 
% This is possibly the largest 
In NC-SC bilevel problems, a first-order method applied to the 
 hyper-objective $F(\vx)$ acts as a second-order method on $g$
because 
\begin{align} \label{eq:hyper-grad}
    \nabla F(\vx) = \nabla_x f(\vx,\vy^*(\vx)) - \nabla_{xy}^2 g(\vx,\vy^*(\vx)) [ \nabla_{yy}^2 g(\vx,\vy^*(\vx))]^{-1} \nabla_y f(\vx,\vy^*(\vx)).
\end{align}
In large-scale problems, directly inverting the Hessian matrix $\nabla_{yy}^2 g(\vx,\vy)$ is very expensive. There are mainly two lines of work to address this issue. 
The first line studies algorithms based on Hessian-vector product (HVP) oracles. Among them, the best-known 
theoretical guarantees \citep{ji2021bilevel,arbel2022amortized} show that they
can find an $\epsilon$-stationary point of $F(\vx)$ with $\gO(\bar \kappa_y^{4} \epsilon^{-2})$ and $\gO(\bar \kappa_y^9 \epsilon^{-4})$  first-order oracle complexity for smooth deterministic and stochastic problems, respectively, 
where $\bar \kappa_y$ is the global condition number for the problem (see Definition~\ref{dfn:cond-number}). In addition,
\citet{yang2023accelerating} showed an improved upper bound of $\tilde \gO(\bar \kappa_y^{13/4} \epsilon^{-7/4})$ for second-order smooth deterministic problems.
% \citet{ji2021bilevel,arbel2022amortized} also proposed methods with  upper bounds on stochastic HVP calls for smooth stochastic problems.

% Although these methods originate from the second-order method that uses the hyper-gradient in Eq. (\ref{eq:hyper-grad}), each iteration of these methods requires only Hessian-vector product oracles, which are almost as cheap as 
% first-order oracles \citep{yang2023achieving,allen2018neon2,agarwal2017finding}. In this paper, we use the terminology ``second-order-ish'' methods to call them since they have the features of both second-order and first-order methods.

Since HVP oracles are more expensive than gradient oracles,
another line of work
\citep{liu2020generic,liu2021value,liu2021towards,liu2022bome, kwon2023fully,shen2023penalty,lu2024first,pan2024scalebio}
directly uses gradients to optimize $F(x)$ by solving the penalty function:
\begin{align} \label{eq:Lag}
   \min_{\vx \in \sR^{d_x}, \vy \in \sR^{d_y}}  f(\vx,\vy) + \lambda (g(\vx,\vy) - g(\vx,\vy^*(\vx)),
\end{align}
where $\lambda = \Omega(\epsilon^{-1})$ is a large penalty coefficient.
%Thanks to Danskin's theorem, the gradient of the penalty function only involves first-order information. The behavior of penalty methods has been extensively studied in many recent works . % \citet{liu2022bome,lu2024first} proved non-asymptotic convergence to the KKT points of problem (\ref{eq:Lag}). \citet{shen2023penalty,shen2025penalty} studied the relationship of local solutions to the original problem and the penalty problem.
Among them, 
\citet{kwon2023fully} 
developed the first fully first-order method that can find an $\epsilon$-stationary point of $F(\vx)$ with $\tilde \gO({\rm poly}(\bar \kappa_y) \epsilon^{-3})$ first-order oracles for smooth deterministic problems.
In a follow-up work, \citet{chen2023near} improved the complexity to $\tilde \gO(\bar \kappa_y^{4} \epsilon^{-2})$, which achieves the near-optimal dependency on $\epsilon$. This result was also extended to $\tilde \gO(\bar \kappa_y^{15/4} \epsilon^{-7/4})$ for second-order smooth deterministic problems \citep{chen2023near}, and to $\tilde \gO(\bar \kappa_y^{11} \epsilon^{-6})$ for smooth stochastic problems \citep{kwon2024complexity,chen2023near}.

Although both HVP-based and gradient-based methods can achieve (near)-optimal $\epsilon^{-2}$ rates for NC-SC bilevel problems, their upper bounds have high~$\bar \kappa_y$~dependencies compared with 
NC-SC minimax problems, which are special bilevel problems with $f = -g$. For instance, 
\citet{lin2020near} proposed an algorithm with a $\tilde \gO(\sqrt{\kappa_y} \epsilon^{-2})$ complexity for smooth NC-SC 
deterministic minimax problems, where $\kappa_y \le \bar \kappa_y$ is the condition number of the lower-level problem (see Definition \ref{dfn:cond-number}). Under the same setting,
\citet{li2021complexity,zhang2021complexity} proved a nearly matching lower bound of $\Omega(\sqrt{\kappa_y} \epsilon^{-2})$.
However, the state-of-art methods 
for deterministic smooth NC-SC bilevel problems requires an $\gO(\bar \kappa_y^4 \epsilon^{-2})$ complexity.
Given this, we wonder  \textit{whether NC-SC bilevel problems are provably more challenging than NS-SC minimax problems measured by the condition number dependency.}
% For second-order smooth problems, \citet{chen2023near} proposed the AccF${}^2$BA method 
% with 
% For stochastic problems, \citet{kwon2024penalty,chen2023near} showed the double-loop F${}^2$SA achieves the  stochastic gradient complexity.
%\citet{pan2024scalebio} showed that penalty methods, combined with memory-efficient training techniques, can scale to a 32B-sized large language model (LLM) data cleaning problem, which can not be done by HVP-based methods.
% Motivated by the restarted accelerated gradient descent method \citep{li2023restarted}, \citet{chen2023near} also proposed an accelerated method AccF${}^2$BA 
% under the additional smoothness assumptions that guarantee $\nabla^2 F(\vx)$ is Lipschitz continuous \citep{huang2025efficiently}.
% Since these methods only rely on gradient oracles, they rapidly became the dominant approach for large-scale bilevel problems \citep{liu2022bome,pan2024scalebio}. 
%Very recently, \citet{pan2024scalebio} showed that these penalty methods can scale to 32B-sized by combining .
% %For both HVP- and gradient-based methods for NC-SC bilevel problems, we all observe . 
% \begin{center}
%     % \textit{Whether the high condition number dependency is necessary due to the intrinsic difficulty of the problem, or is avoidable using more intricate algorithm designs?}
   
% \end{center}

% However, we find no specific lower bounds for general bilevel optimization, and the best applicable results were given by minimax optimization where $g$ is always equivalent to $-f$: In the context of NC-SC minimax optimization, \citet{li2021complexity,zhang2021complexity} established the tight $\Omega(\sqrt{\kappa_y} \epsilon^{-2})$ lower bound, where $\kappa_y$ is the \textit{lower-level condition number} (see Definition \ref{dfn:cond-number}). Moreover,
% their results can be extended to 
% $\Omega(\sqrt{\kappa_y} \epsilon^{-12/7})$ for second-order smooth problems \citep{wang2024efficient}, and to 
%  $\Omega(\kappa_y^{1/3} \epsilon^{-4})$  for stochastic problems \citep{li2021complexity}.
% However, all these lower bounds 
% have a huge gap in condition number compared to the corresponding upper bounds in bilevel optimization.


\begin{table*}[t] 
    \centering
     \caption{We present the oracle complexities of different \textbf{deterministic} methods for finding an $\epsilon$-stationary point for \textbf{smooth} NC-SC bilevel problems.
     }
    \label{tab:det}
    \begin{tabular}{c c c}
    \hline 
    Oracle   &  Method & Complexity \\ 
    \hline \addlinespace 
    & BA \citep{ghadimi2018approximation} & $\gO( \bar \kappa_y^5 \epsilon^{-5/2})$ 
    \\ \addlinespace
      & AID/ITD  \citep{ji2021bilevel} & $\gO(\bar \kappa_y^4 \epsilon^{-2})$ \\ 
    \vspace{-2.5mm} 
    HVP \\
      &  AmIGO \citep{arbel2022amortized} & $\gO(\bar \kappa_y^4 \epsilon^{-2})$ \\ \addlinespace
      & \cellcolor{gray!20} Lower Bound (Theorem \ref{thm:NC-SC}) 
      & 
      \cellcolor{gray!20} 
      $\Omega(\kappa_y^2 \epsilon^{-2})$  \\
    % & ITD \citep{ji2021bilevel} & $\tilde \gO(\kappa_y^4 \epsilon^{-2})$ \\
    \hline \addlinespace 
    & F${}^2$SA \citep{kwon2023fully} & $\tilde \gO(\bar \kappa_y^7 \epsilon^{-3})$ \\
    \addlinespace
    & F${}^2$BA \citep{chen2023near} & $\tilde \gO(\bar \kappa_y^4 \epsilon^{-2})$ \\
    \addlinespace 
   Gradient & \cellcolor{gray!20} F${}^2$BA${}^+$ (Theorem \ref{thm:ub-p1}) & \cellcolor{gray!20} $\tilde \gO(\bar \kappa_y^{7/2} \epsilon^{-2})$ \\ 
    \addlinespace 
    % & F${}^2$BA${}^+_{\rm llq}$ (Theorem \ref{thm:ub-p1}) & $\tilde \gO(\kappa_y^{5/2} \epsilon^{-2})$ \\ \addlinespace 
    & Lower Bound$_{f=-g}$ \citep{li2021complexity} %~\textsuperscript{{\color{blue}*}} 
    & $\Omega( \sqrt{\kappa_y} \epsilon^{-2})$ \\
    \addlinespace
    & \cellcolor{gray!20} Lower Bound (Theorem \ref{thm:NC-SC}) & 
    \cellcolor{gray!20}
    $\Omega(\kappa_y^2 \epsilon^{-2})$ \\ 
    \hline 
    \end{tabular}
\end{table*}


%  extended their result to second-order smooth problems and showed an 
%  lower bound. \citet{li2021complexity} extended the $\Omega(\sqrt{\kappa_y} \epsilon^{-2})$ deterministic lower bound to the stochastic setting and showed an 
% stochastic lower bound.
% But these lower bounds 
% in bilevel optimization 
% lower bounds 
% from minimax optimization 
% look not tight if  since they have

To answer this question, we first provide lower bounds for NC-SC bilevel problems. Almost all previously applicable lower bounds are attained in minimax problems where~$f = -g$. In this work, by leveraging the asymmetry of the objectives $f$ and $g$, we modify the chain-like hard instance for NC-SC minimax problems \citep{li2021complexity} to create a longer effective chain in $\vx$.
Our construction leads to larger $\kappa_y$ dependencies in various setups:
\vspace{-1mm}
\begin{enumerate}
    \item For deterministic NC-SC bilevel problems, we show lower bounds of 
$\Omega(\kappa_y^2 \epsilon^{-2})$ and  $\Omega(\kappa_y^{25/14}\epsilon^{-12/7})$ under first- and second-order smoothness assumptions, respectively, both improving the trivial lower bounds 
of $\Omega(\sqrt{\kappa_y} \epsilon^{-2})$ \citep{li2021complexity} and 
 $\Omega(\sqrt{\kappa_y} \epsilon^{-12/7})$ 
\citep{wang2024efficient} from minimax problems.
\vspace{-1mm}
\item For stochastic NC-SC bilevel problems, we prove a lower bound of 
$\Omega(\kappa_y^4 \epsilon^{-4})$, improving the lower bound of $\Omega(\kappa^{1/3} \epsilon^{-4})$ 
from minimax problems \citep{li2021complexity}.
\vspace{-1mm}
\item For convex-strongly-convex (C-SC) and strongly-convex-strongly-convex (SC-SC) bilevel problems, we also show lower bounds of 
$\Omega(\kappa_y^{5/4} / \sqrt{\epsilon})$ and $\tilde \Omega(\kappa_y^{5/4} \sqrt{\kappa_x})$, respectively, improving 
prior results \citep{ji2023lower}
of $\tilde \Omega(\min \{ \kappa_y, \epsilon^{-3/2} \} / \sqrt{\epsilon} )$ and~$\tilde \Omega(\kappa_y \sqrt{\kappa_x})$.
\end{enumerate}
% resulting in an $\Omega(\kappa_y^2 \epsilon^{-2})$ lower bound with a $\kappa_y^{3/2}$ improvement.  
%achieving a stronger $\kappa_y$ dependency in the lower bounds. 
% Answering this question requires studying both upper and lower complexity bounds for bilevel optimization. 
% In this work, we enhance the construction in \citep{li2021complexity,wang2024efficient}  for bilevel problems to establish a much larger dimension dependency in the lower bounds. Consequently, we prove the lower bound of $\Omega(\kappa_y^2 \epsilon^{-2})$, $\Omega(\kappa_y^{25/14} \epsilon^{-12/7})$, $\Omega(\kappa_y^{17/10} \epsilon^{-8/5})$, and $\Omega(\kappa_y^4 \epsilon^{-4})$ for smooth, second-order smooth, arbitrarily smooth, and stochastic NC-SC bilevel problems, respectively. 
% Moreover, although our construction is mainly designed for NC-SC problems, it also implies better results % Specifically, we improve  their lower bound 

Then, we provide improved upper bounds to reduce the gap between existing ones and our lower bounds. For smooth NC-SC problems, we show an improved upper bound of $\tilde \gO(\bar \kappa_y^{7/2} \epsilon^{-2})$ by solving the lower-level problems
in the F${}^2$BA method \citep{chen2023near} using accelerated gradient descent (AGD) \citep{nesterov1983method}. Our result improves the best-known $\gO(\bar \kappa_y^4 \epsilon^{-2})$ results \citep{ji2021bilevel,arbel2022amortized,chen2023near} by a factor of $\sqrt{\kappa_y}$. Moreover, if we further applied nonconvex AGD \citep{li2023restarted} in the upper level, the upper bounds can be improved to 
$\tilde \gO(\bar \kappa_y^{13/4} \epsilon^{-7/4})$ for second-order smooth NC-SC problems, which matches the state-of-the-art guarantee \citep{yang2023accelerating} without directly using any HVP oracles or their finite difference approximations \citep{yang2023achieving}.

% By solving the lower-level problems in fully first-order methods \citep{kwon2023fully,chen2023near} ,
% we obtain the improved upper bound of and  for first-order and second-order smooth problems, respectively. Our results improve the previous method \citep{chen2023near} by a factor of $\sqrt{\kappa_y}$.
To summarize, we compare our lower and upper bounds with existing results for smooth deterministic, second-order smooth deterministic, and smooth stochastic bilevel problems in Table \ref{tab:det}, \ref{tab:det-second-smooth}, and~\ref{tab:stoc},
respectively. 
In addition, since we prove our lower bounds on lower-level quadratic problems, all our lower bounds hold for both gradient- and HVP-based methods (see Lemma \ref{lem:zr-grad-hvp}). 
In this simpler setting, where $g(\vx,\vy)$ is quadratic, our upper bounds can also be refined to 
$\tilde \gO(\kappa_y^{5/2} \epsilon^{-2})$ and $\tilde \gO(\bar \kappa_y^{9/4} \epsilon^{-7/4})$ for first- and second-order smooth deterministic problems, respectively, indicating the gaps of $\sqrt{\kappa_y}$ and $\kappa_y^{13/28} \epsilon^{-1/28}$ still exist for quadratic lower-level problems.

% When the lower-level problem is quadratic, we show
% the above complexities can also be refined to 


%Since our corresponding $\Omega(\kappa_y^2 \epsilon^{-2})$ and $\Omega(\kappa_y^{25/14} \epsilon^{-7/4})$ lower bounds are also proved with lower-level quadratics, it suggests that a gap still exists even under this simpler setting. 

% Our results suggest there is a significant gap 
% For second-order smooth NC-SC minimax problems, \citet{wang2024efficient} proved a $\tilde \gO(\sqrt{\kappa_y} \epsilon^{-7/4})$ upper bound and an ; For stochastic smooth NC-SC minimax problems, the current best upper and lower bounds are $\tilde \gO(\kappa_y \epsilon^{-4})$ by \citet{zhang2022sapd+} and by \citet{li2021complexity}, respectively.
% We first study the lower complexity bounds for this problem. Before our work,  were given by the case $f = -g$, which is also known as minimax optimization, and there is no 
% which are much worse than the results for minimax optimization, a special case of bilevel optimization when $f= -g$:
% This answer to this question is (almost) clear for the special case when $f=-g$, which is only known as minimax optimization in the literature:
% For smooth minimax problems,
% \citet{lin2020near} proposed a method with an $\tilde \gO(\sqrt{\kappa_y} \epsilon^{-2})$ upper complexity bound, and \citet{li2021complexity} proved a nearly matched lower bound of $\Omega(\sqrt{\kappa_y} \epsilon^{-2})$. 
% For Hessian smooth minimax problems, \citet{wang2024efficient} proved a $\tilde \gO(\sqrt{\kappa_y} \epsilon^{-7/4})$ upper bound and an $\Omega(\sqrt{\kappa_y} \epsilon^{-12/7})$ lower bound.
% However, for the general bilevel optimization ($f \ne -g)$, the current upper bounds are much worse than the complexity of minimax optimization.
% Given the above observation, it is natural to ask:
% \textit{What is the complexity of first-order bilevel optimization, especially the condition number dependency? Can we prove lower bounds to show that the high $\kappa_y$ dependency is necessary for bilevel optimization? Can we prove improved upper bounds to reduce the $\kappa_y$ dependency in existing results?
% }

% \subsection{Our Contributions}
% This paper answers the above question from two sides, by both proving novel lower bounds and improved upper bounds for various setups in bilevel optimization. 
% \begin{enumerate}
%     \item For smooth NC-SC problems, we prove an $\Omega(\kappa_y^2 \epsilon^{-2})$ lower bound and a $\tilde \gO(\kappa_y^{7/2} \epsilon^{-2})$ upper bounds.
%     \item For second-order smooth NC-SC problems, we prove an $\Omega(\kappa_y^{25/14} \epsilon^{-12/7})$ lower bound and a $\tilde \gO(\kappa_y^{13/4} \epsilon^{-7/4})$ upper bound.
%     \item For stochastic smooth NC-SC problems, we prove an $\Omega(\kappa_y^{4} \epsilon^{-4}  )$ lower bound, where the best-known upper bound is $\tilde \gO(\kappa_y^{9} \epsilon^{-4})$ in \citep{ji2021bilevel}.
% \end{enumerate}
% We compare our results with prior results in Table \ref{tab:det}, \ref{tab:det-second-smooth} and~\ref{tab:stoc}, respectively, for the above three settings.
% The improvements in our upper bounds come from using accelerated gradient descent \citep{nesterov1983method} to solve the lower-level problem in fully first-order methods \citep{kwon2023fully,chen2023near}. The more interesting part is our lower bounds, which relies on a generic framework (see Section \ref{subsec:framework}) to construct zero-chains that can simultaneously reflect the two-source $\kappa_y$ dependency in bilevel optimization. Under the same framework, we also show novel results for other setups, including:
% \begin{enumerate}
%     \item[4.] An $\Omega(\kappa_y^{17/10} \epsilon^{-8/5})$ lower bound for arbitrarily smooth problems.  \vspace{1mm}
%     \item[5.] An $\Omega(\kappa_y^{5/4}/\sqrt{\epsilon})$ lower bound for convex-strongly-convex problems and an $\Omega(\kappa_y^{5/4} \sqrt{\kappa_x})$ lower bound for strongly-convex-strongly-convex problems.
% \end{enumerate}
% We highlight that our lower bounds are the first result for settings 1-4.
% Setting 5 is the only one that has been studied by prior work \citep{ji2023lower}, but their results are worse than ours by a factor of $\kappa_y^{1/4}$, which is also a strong evidence of the advantage of our novel lower bound framework.




% We answer (Q1) by proving new lower bounds under various setups.
% \begin{enumerate}
%     \item[(L1)] We prove lower bounds for deterministic zero-respecting algorithms applied on a hard bilevel problem where $f$ has Lipschitz continuous derivatives up to $p$th-order and $g$ is a smooth strongly convex quadratic, showing any algorithms require at least $\Omega(\kappa_y^2 \epsilon^{-2})$, $\Omega(\kappa_y^{25/14} \epsilon^{-12/7} )$, or $\Omega( \kappa_y^{17/10} \epsilon^{-8/5})$ DFO calls to find an $\epsilon$-stationary point when $p=1$, $p=2$, or $p \ge 3$, respectively.
%     \item[(L2)] We extend our results to the case the hyper-objective $F(\vx)$ is additionally convex or strongly convex, showing the lower bounds of $\Omega( \kappa_y^{5/4} \epsilon^{-1/2})$ or $\Omega(\kappa_y^{5/4} \sqrt{\kappa_x})$.
%     \item[(L3)] We also extend our results to the stochastic setting,  showing an $\Omega(\kappa_y^{7/3} \epsilon^{-4})$ lower bound of the stochastic first-order oracle (SFO) calls.
% \end{enumerate}
% Our result suggests a significant gap between our lower bounds and existing upper bounds, which motivates us to turn to (Q2).  
% We obtain the following upper bounds by solving the lower-level problem in fully first-order methods \citep{kwon2023fully,chen2023near} with accelerated gradient descent (AGD) \citep{nesterov1983method}:
% \begin{enumerate}
%     \item[(U1)] For nonconvex-strongly-convex bilevel problems where $f$ and $g$ has Lipschitz continuous derivatives up to $p$th and $(p+1)$th-order derivatives, respectively, our methods can provably find an $\epsilon$-stationary point with $\tilde \gO(\kappa_y^{7/2} \epsilon^{-2})$ or $\tilde \gO( \kappa_y^{13/4} \epsilon^{-7/4})$ DFO calls when $p=1$ or $p=2$, respectively. When the lower-level problem $g$ is quadratic, the upper bounds can be tighten to $\tilde \gO(\kappa_y^{5/2} \epsilon^{-2})$ or $\tilde \gO( \kappa_y^{9/4} \epsilon^{-7/4})$, which only have small gaps of orders $\sqrt{\kappa_y}$ or $\kappa_y^{13/28} \epsilon^{1/28}$ between our lower bounds in (L1).
%     % \item[(U2)] For bilevel problems satisfying the same assumptions above with the hyper-objective $F(\vx)$ additionally being convex or strongly convex, our methods can provably find an $\epsilon$-stationary point with  $\tilde \gO(\kappa_y^2 \epsilon^{-1/2})$ or  $\tilde \gO(\kappa_y^2 \sqrt{\kappa_x})$ DFO calls. 
%     %  When the lower-level problem $g$ is quadratic, the upper bounds can be tighten to $\tilde \gO(\kappa_y^{3/2} \epsilon^{-1/2})$ and $\tilde \gO( \kappa_y^{3/2} \sqrt{\kappa_x})$, which only have a $\kappa_y^{1/4}$ gap between our lower bounds in (L2).
%      \item[(U3)] AC-SA in the lower-level. 
% \end{enumerate}

\begin{table*}[htbp] 
    \centering
     \caption{We present the oracle complexities of different \textbf{deterministic} methods for finding an $\epsilon$-stationary point for \textbf{second-order smooth} NC-SC bilevel problems.}
    \label{tab:det-second-smooth}
    \begin{tabular}{c c c}
    \hline 
    Oracle   &  Method & Complexity \\ 
    \hline \addlinespace 
    & AIPUN \citep{wang2024efficient} &  $\tilde \gO(\bar \kappa_y^{19/4} \epsilon^{-7/4})$ \\ \addlinespace
    HVP    & RAHGD \citep{yang2023accelerating} & $\tilde \gO(\bar \kappa_y^{13/4} \epsilon^{-7/4})$ \\  \addlinespace 
    & \cellcolor{gray!20} Lower Bound (Theorem \ref{thm:NC-SC}) & \cellcolor{gray!20} $\Omega(\kappa_y^{25/14} \epsilon^{-12/7})$ \\ 
    % & ITD \citep{ji2021bilevel} & $\tilde \gO(\kappa_y^4 \epsilon^{-2})$ \\
    \hline \addlinespace 
    & AccF${}^2$BA \citep{chen2023near} & $\tilde \gO(\bar \kappa_y^{15/4} \epsilon^{-7/4})$ \\
    \addlinespace 
    & \cellcolor{gray!20} AccF${}^2$BA${}^+$ (Theorem \ref{thm:ub-p2}) & 
    \cellcolor{gray!20}
    $\gO(\bar \kappa_y^{13/4} \epsilon^{-7/4})$  \\ 
    \vspace{-2.5mm} 
    Gradient \\
    % Gradient & AccF${}^2$BA${}^+_{\rm llq}$ (Theorem \ref{thm:ub-p2}) & $\gO(\bar \kappa_y^{9/4} \epsilon^{-7/4})$ \\ \addlinespace
    & Lower Bound$_{~f=-g}$ \citep{wang2024efficient} & $\Omega( \sqrt{\kappa_y} \epsilon^{-12/7})$ \\
    \addlinespace
    &
    \cellcolor{gray!20}
    Lower Bound (Theorem \ref{thm:NC-SC}) &
    \cellcolor{gray!20}
    $\Omega(\kappa_y^{25/14} \epsilon^{-12/7})$ \\ 
    \hline
    \end{tabular}
  % \begin{tablenotes} {\scriptsize   
  % \item {\color{blue} *} {The lower bound from minimax optimization. }
  %  }
  %   \end{tablenotes}
  %    \begin{tablenotes} {\scriptsize   
  % \item {\color{blue} (a)} {The lower bound from minimax optimization. }
  %  }
  %   \end{tablenotes}
\end{table*}


\begin{table*}[htbp] 
    \centering
     \caption{We present the oracle complexities of different \textbf{stochastic} methods for finding an $\epsilon$-first-order stationary point for \textbf{smooth} NC-SC bilevel problems.}
    \label{tab:stoc}
    \begin{tabular}{c c c}
    \hline 
    Oracle   &  Method & Complexity \\ 
    \hline \addlinespace 
      & BSA \citep{ghadimi2018approximation} & ~~$\gO({\rm poly}(\bar \kappa_y) \epsilon^{-6})$ \\ \addlinespace 
     & TTSA \citep{hong2023two} & ~~$\gO({\rm poly}(\bar \kappa_y)\epsilon^{-5})$ \\ \addlinespace
    HVP & stocBiO \citep{ji2021bilevel} & $\gO(\bar \kappa_y^{9} \epsilon^{-4})$
    \\ \addlinespace
    & 
    AmIGO \citep{arbel2022amortized} 
    & $\gO(\bar \kappa_y^{9} \epsilon^{-4})$
    \\ \addlinespace
    %  & ALSET \citep{chen2021closing} & $\tilde \gO( {\rm poly} (\bar \kappa_y) \epsilon^{-4} )$\\ \addlinespace
    % &
    % MA-SOBA \citep{chen2024optimal} & $\gO( {\rm poly} (\bar \kappa_y) \epsilon^{-4} )$ \\ \addlinespace
    &  \cellcolor{gray!20} Lower Bound (Theorem \ref{thm:NC-SC-stoc}) & 
     \cellcolor{gray!20}
    $\Omega(\kappa_y^{4} \epsilon^{-4})$ \\ 
    %\addlinespace 
    % & ITD \citep{ji2021bilevel} & $\tilde \gO(\kappa_y^4 \epsilon^{-2})$ \\
    \hline \addlinespace 
    & F${}^2$SA \citep{kwon2023fully} & ~~~$\tilde \gO({\rm poly}(\bar \kappa_y) \epsilon^{-7})$ \\
    %~\textsuperscript{{\color{blue}(b)}}  \\
    \addlinespace 
    & F${}^2$SA \citep{kwon2024complexity} & $\tilde \gO({\rm poly}(\bar \kappa_y) \epsilon^{-6} )$ 
   \\ \addlinespace
    Gradient &  F${}^2$BSA \citep{chen2023near} & $\tilde \gO(\bar \kappa_y^{11} \epsilon^{-6})$ \\ 
    \addlinespace 
   % & F${}^2$SA (to post on arXiv) & $\tilde \gO(\bar \kappa_y^{11} \epsilon^{-6})$  \\ 
   %  \addlinespace 
    & Lower Bound$_{f=-g}$ \citep{li2021complexity} %
    %~\textsuperscript{{\color{blue}*}} 
    & $\Omega( \kappa_y^{1/3} \epsilon^{-4})$ \\
    \addlinespace
    &  \cellcolor{gray!20} Lower Bound (Theorem \ref{thm:NC-SC-stoc}) & 
     \cellcolor{gray!20}
    $\Omega(\kappa_y^{4} \epsilon^{-4})$ \\ 
    \hline
    \end{tabular}
%      \begin{tablenotes} {\scriptsize   {%\color{blue} (a)} Stochastic HVP methods additionally the existence of an estimator to $\nabla^2 g(\vx,\vy)$. 
% %   \item  {\color{blue} (b)} {We use $\kappa_y^{p_1}$, $\kappa_y^{p_2}$, and $\kappa_y^{p_3}$ to denote the not given polynomial dependency in \citep{ghadimi2018approximation,hong2023two,kwon2023fully}.} 
% \item {\color{blue} *} {The lower bound from minimax optimization. }
% }}
%     \end{tablenotes}
\end{table*}

\subsection{Additional Related Work}
In addition to the works discussed above, we further review the closely related work on lower bounds, including lower bounds for minimization, minimax, and bilevel problems.
%not only other lower bounds for bilevel optimization, but also the results for minimization and minimax problems that we refer to in constructing our lower bounds.
% . 
% % In this subsection, we review 
% \vspace{4mm}

\paragraph{Lower bounds for minimization problems} In the textbooks, \citet{nemirovskij1983problem,nesterov2018lectures} established tight lower bounds for convex optimization using chain-like quadratic functions. Recently, \citet{carmon2020lower,carmon2021lower} further abstracted the above idea into the concept of zero-chain as a generic approach for proving lower bounds. By carefully constructing zero-chains, \citet{carmon2020lower} proved the tight $\Omega(\epsilon^{-p/(p+1)})$  lower bound for 
$p$th-order methods to find $\epsilon$-stationary points of $p$th-order smooth nonconvex functions. In the same series of work, \citet{carmon2021lower}
proved $\Omega(\epsilon^{-12/7})$ and  $\Omega(\epsilon^{-8/5})$  lower bounds for 
first-order methods to optimize second-order smooth and arbitrarily smooth nonconvex functions, respectively.  
% proved lower bounds for  when optimizing high-order smooth nonconvex functions, and showed that are the lower bounds for minimizing
Many advancements were made following these seminal works. \citet{arjevani2023lower} proved the tight $\Omega(\epsilon^{-4})$ and $\Omega(\epsilon^{-3})$ lower bounds for nonconvex stochastic first-order optimization under the general expectation and the mean-squared smooth setting, respectively. 
\citet{fang2018spider,li2021page} proved the tight $\Omega(n + \sqrt{n} \epsilon^{-2})$ lower bound for nonconvex finite-sum first-order optimization; \citet{arjevani2020second} proved 
the tight $\Omega(\epsilon^{-7/2})$ and $\Omega(\epsilon^{-3})$ lower bounds for first-order stochastic nonconvex optimization on second-order smooth objectives, and
high-order stochastic nonconvex optimization on arbitrarily smooth objectives, respectively.

\paragraph{Lower bounds for minimax problems} Inspired by the lower bounds in minimization problems, researchers also proved tight lower bounds for (first-order) minimax optimization under many setups. 
\citet{nemirovski2004prox} gave tight $\Omega(\epsilon^{-1})$ lower bound for convex-concave minimax problems. \citet{ouyang2021lower,zhang2022lower} proved the tight 
 $\Omega(\sqrt{\kappa_y/ \epsilon})$ and
$\tilde \Omega(\sqrt{\kappa_x \kappa_y})$
 lower bounds for C-SC and SC-SC problems, respectively.
\citet{han2024lower} extended those results to the finite-sum setting.
\citet{li2021complexity,zhang2021complexity} proved the tight $\Omega(\sqrt{\kappa_y} \epsilon^{-2})$ lower bound for nonconvex-strongly-convex (NC-SC) problems.
Their results have also been extended to
$\Omega(\sqrt{\kappa_y} \epsilon^{-12/7})$, $\Omega(n + \sqrt{n \kappa_y} \epsilon^{-2})$, and  
$\Omega(\kappa_y^{1/3} \epsilon^{-4})$ for 
second-order smooth problems \citep{wang2024efficient},
finite-sum problems \citep{zhang2021complexity}, and 
general stochastic problems \citep{li2021complexity}, respectively.
% general expectation setting,  
% with extensions of  to the general expectation setting \citep{li2021complexity}, and 
%  to the finite-sum setting \citep{zhang2021complexity}. \citet{wang2024efficient} proved an  lower bound for second-order smooth NC-SC minimax problems. 

\paragraph{Lower bounds for bilevel problems} Although the lower bounds for both minimization and minimax problems have been extensively studied, we find there are only a few works on the lower complexity bounds of bilevel optimization, except for the following ones. In their novel work,
\citet{ji2023lower} proved lower bounds for (S)C-SC bilevel problems instead of the more common NC-SC problems. 
The established the lower bound of $\tilde \Omega(\min \{ \kappa_y, \epsilon^{-3/2} \} / \sqrt{\epsilon}  )$ and 
$\tilde \Omega (\kappa_y \sqrt{\kappa_x})$
for C-SC and SC-SC setups, respectively.
% For the C-SC setups, they established a lower bound of , where $\kappa_y$ is the condition number of lower-level problems. For the SC-SC setups, they established a lower bound 
% However, to the best of our knowledge, there are no results for the general nonconvex-strongly-convex (NC-SC) bilevel problems before our work, which is also the most important application of bilevel optimization. 
\citet{kwon2024complexity} established tight upper and lower bounds for $r$-$\vy^*(\vx)$-aware optimization, 
where the oracle only gives  $r$-estimates of $\vy^*(\vx)$ along with $r$-locally reliable gradients. This setting recovers stochastic bilevel optimization if $r = +\infty$, but requires $r = 
\gO(\epsilon)$, which is different from the setting of bilevel optimization with standard oracles.
\citet{chen2024finding} considered bilevel problems without lower-level strong convexity, and constructed a hard instance such that any first-order algorithm will get stuck at $\vx_0$ within finite iterations. \citet{dagreou2024lower} proved a lower bound for finite-sum NC-SC problems, but the $\Omega(n+ \sqrt{n} \epsilon^{-2})$ lower bound misses the condition number dependency and is subsumed by the $\Omega(n + \sqrt{n \kappa_y} \epsilon^{-2})$ lower bound from NC-SC minimax optimization \citep{zhang2021complexity}.

 \paragraph{Notations} For a vector $\vx \in \sR^d$, 
 we let $x_i $ be its $i$th coordinate; define its support by ${\rm supp}(\vx) = \{ i \mid \vx_i \ne 0 \}$ and its progress by ${\rm prog} = \max\{ i \ge 0 \mid \vx_i \ne 0 \}$. We use $\Vert \,\cdot \Vert$ to denote both the Euclidean norm of a vector and the operator norm of a matrix/tensor. 
 Given an event $A$, we let $\vone_A$ denote its indicator function. We let $\ve_i$ be the vector with its $i$th element being $1$ and all the other elements being $0$.
 Given functions $f,g : \gX \rightarrow [0,+\infty)$, we use $f = \gO(g)$ to denote there exists a numerical constant $C< +\infty$
such that $f(x) \le C g(x)$ for all $x \in \gX$,  $f = \Omega(g)$ to denote $g = \gO(f)$,  and $f = \tilde \gO(g)$ to denote $f = \gO(g \max\{1, \log g \})$. We also denote $f \asymp g$ when $f = \gO(g)$ and $f = \Omega(g)$ both holds.

% \newpage

\section{Problem Setup}

In this section, we introduce the bilevel function class considered in this paper, the algorithm class to which our lower bound applies, and the complexity measurement of the algorithms.

\subsection{Function Class}

Our main focus is the general high-order smooth nonconvex-strongly-convex (NC-SC) bilevel problems as follows, where the case $p=1$ is introduced by \citet{ghadimi2018approximation}, and the case $p=2$ is recently introduced by \citet{huang2025efficiently}.

\begin{dfn}[$p$th-order smooth NC-SC problem] \label{dfn:NC-SC-func}
Given $p \in \sN_+$, $L_0, \cdots,L_{p+1} >0$, $\Delta>0$, and $\mu_y \in (0,L_1]$, 
we use $\gF^{\text{nc-sc}}(L_0,\cdots,L_{p+1},\mu_y, \Delta)$ to denote the set of all bilevel problems $(f,g)$, where the upper-level function $f:\sR^{d_x} \rightarrow \sR$ and lower-level function $g: \sR^{d_y} \rightarrow \sR$ satisfy the following assumptions:
\begin{enumerate}
    \item $f$ have $L_q$-Lipschitz continuous $q$th-order derivative for all $q = 1,\cdots,p$; and $g$ have $L_q$-Lipschitz continuous $q$th-order derivative for all $q = 1,\cdots,p+1$, that is, for every $\vx_1, \vx_2 \in \sR^{d_x}$ and $\vy_1,\vy_2 \in \sR^{d_y}$, we have
    \begin{align*}
        \Vert \nabla^q f(\vx_1,\vy_1) - \nabla^q f(\vx_2,\vy_2) \Vert &\le L_q \Vert (\vx_1,\vy_1) - (\vx_2,\vy_2) \Vert, \quad q=0,\cdots,p~; \\
        \Vert \nabla^q g(\vx_1,\vy_1) - \nabla^q g(\vx_2,\vy_2) \Vert &\le L_q \Vert (\vx_1,\vy_1) - (\vx_2,\vy_2) \Vert, \quad q=1,\cdots,p+1.
    \end{align*}
    % \item $g$ have $L_q$-Lipschitz continuous $q$th-order derivative for $q = 1,\cdots,p+1$, that is, for every $\vx_1, \vx_2 \in \sR^{d_x}$ and $\vy_1,\vy_2 \in \sR^{d_y}$,
    % \begin{align*}
    %     \Vert \nabla^q f(\vx_1,\vy_1) - \nabla^q f(\vx_2,\vy_2) \Vert &\le L_q \Vert (\vx_1,\vy_1) - (\vx_2,\vy_2) \Vert, \\
    %     \Vert \nabla^q g(\vx_1,\vy_1) - \nabla^q g(\vx_2,\vy_2) \Vert &\le L_q \Vert (\vx_1,\vy_1) - (\vx_2,\vy_2) \Vert.
    % \end{align*}
    \item For any fixed $\vx \in \sR^{d_x}$, $g$ is $\mu_y$-strongly convex in $\vy$, that is, for any $\vy_1,\vy_2 \in \sR^{d_y}$,
    \begin{align*}
        g(\vx,\vy_2) \ge g(\vx,\vy_1) + \langle \nabla_y g(\vx,\vy_1), \vy_2 - \vy_1 \rangle + \frac{\mu_y}{2} \Vert \vy_1 - \vy_2 \Vert^2.
    \end{align*}
    \item The hyper-objective $F(\vx) = f(\vx,\vy^*(\vx))$ in Eq. (\ref{prob:main-BLO}) is lower bounded and 
    \begin{align*}
        F(\vzero) - \inf_{\vx \in \sR^{d_x}} F(\vx) \le \Delta.
    \end{align*}
\end{enumerate}
\end{dfn}

For these problems, we define the lower-level and global condition numbers $\kappa_y$ and $\bar \kappa_y$
as follows. By definition, we have $\kappa_y \le \bar \kappa_y$.
% We define the condition numbers of these problems as follows.

\begin{dfn} \label{dfn:cond-number}
    For the $p$th-order smooth NC-SC function class in Definition \ref{dfn:NC-SC-func}, we define $\kappa_y: =L_1/\mu_y$ as the condition number of $g$  and $\bar \kappa_y: = \bar L / \mu_y$ as the global condition number of the bilevel problem, where $\bar L = \max_{q = 0,\cdots,p+1} L_q$.
\end{dfn}

In our lower bound analyses, we construct bilevel problems with quadratic lower-level functions, which are formally defined in Definition \ref{dfn:qg-sub-class} below. 
Note that for this subclass, we do not make assumptions on the boundedness of $L_0$ since both the lower and upper bounds do not depend on it.

%But our upper bounds always hold for the general bilevel class.

% \begin{dfn}[C-SC function class]\label{dfn:C-SC-func}
% Given $p \in \sN_+$, $L \ge \mu >0$, $D>0$ and and $L_2,\cdots,L_p >0$, we use $\gF^{\text{c-sc}}(L_1,\cdots,L_p,\mu, D)$ to denote the set of all bilevel problems, where the upper-level function $f:\sR^{d_x} \rightarrow \sR$ and lower-level function $g: \sR^{d_y} \rightarrow \sR$ satisfy item 1 and item 2 in Definition \ref{dfn:NC-SC-func}, and item 3 is replaced by:
% \begin{enumerate}
%     \item[3.] $F(\vx)$ is convex in $\vx$.
%     \item[4.] $\vx^* = \arg \min_{\vx \in \sR^{d_x}} F(\vx)$ exists and $\Vert \vx^* \Vert \le D$.
% \end{enumerate}
% \end{dfn}
%, and can be improved for the quadratic $g$ subclass.

\begin{dfn}[Quadratic lower-level problem] \label{dfn:qg-sub-class}
For $p \in \sN_+$,~$L_1,\cdots,L_p>0$, $\mu_y \in (0,L_1]$, and $\Delta>0$, we denote
$\gF^{\text{nc-scq}}(L_1,\cdots,L_{p},\mu_y, \Delta) \subseteq \gF^{\text{nc-sc}}(+\infty, L_1,\cdots,L_p,0,\mu_y, \Delta)$
as the subset of NC-SC problems, where the lower-level function takes the form of
\begin{align} \label{eq:quadratic-g}
    g(\vx,\vy) = \frac{1}{2} \vy^\top \mH \vy+ \vx^\top \mJ \vy +  \vb^\top \vy,
\end{align}
where $\mu_y \mI_{d_y} \preceq \mH \preceq L_1 \mI_{d_y}$ and $-L_1 \mI_{d_y} \preceq \mJ \preceq L_1 \mI_{d_y}$.
% \begin{align*}
%     %  \gF^{\text{sc-scq}}(L_1,\mu_x,\mu_y, D) &\subseteq \gF^{\text{sc-sc}}(+\infty, L_1, 0 ,\mu_x,\mu_y, D), \\
%     % \gF^{\text{c-scq}}(L_1,\mu_y, D) &\subseteq \gF^{\text{c-sc}}(+\infty, L_1,0, \mu_y, D)
% \end{align*}
\end{dfn}

From a complexity theory perspective, it is natural to consider quadratic functions as optimizing convex quadratics is usually as hard as general convex functions \citep[Theorem 2.1.13]{nesterov2018lectures}. Therefore, 
understanding the complexity of quadratic problems is essential to tackle more complex setups.
From a practical perspective, this problem class already covers many useful applications in reinforcement learning, meta-learning, robust optimization, and graph neural network training, as detailed in Appendix \ref{apx:example-Q}.


% \begin{dfn}[Averaged smooth NC-SC class]
%     Given $L_1 \ge \mu_y >0$, $\Delta>0$ and $L_2>0$, we use $\gF^{\text{nc-sc-as}}(L_1,L_2,\mu_y,\Delta) \subset \gF^{\text{nc-sc}}(L_1,L_2,\mu_y,\Delta) $ to denote the set of all bilevel problems $(f,g)$ which lie in the smooth NC-SC function class and additional  (a) have the finite-sum structure such that $f=  \sum_{j=1}^n f_j / n$ and $g = \sum_{j=1}^m g_j / m$; (b) 
%     is $L_1$-averaged smooth, that is, for every $\vx_1,\vx_2 \in \sR^{d_x}$ and $\vy_1,\vy_2 \in \sR^{d_y}$, we have 
%     \begin{align*}
%         \frac{1}{n} \sum_{j=1}^n \Vert \nabla f_j(\vx_1,\vy_1) - \nabla f_j(\vx_2,\vy_2) \Vert^2 \le L_1^2 \Vert (\vx_1,\vy_1)  - (\vx_2,\vy_2) \Vert^2; \\ 
%         \frac{1}{m} \sum_{j=1}^m \Vert \nabla g_j(\vx_1,\vy_1) - \nabla g_j(\vx_2,\vy_2) \Vert^2 \le L_1^2 \Vert (\vx_1,\vy_1)  - (\vx_2,\vy_2) \Vert^2. 
%     \end{align*}
% We also let $\gF^{\text{nc-scq-as}}(L_1,\mu_y, \Delta) \subset  \gF^{\text{nc-sc-as}}(L_1, 0, \mu_y, \Delta) \cap \gF^{\text{nc-scq}}(L_1,\mu_y, \Delta)$ to denote the subclass such that the lower-level function $g$ is quadratic.
% \end{dfn}
\subsection{Algorithm Class}

Following the seminal framework proposed by \citet{nemirovskij1983problem}, we model the algorithmic complexity by the number of oracle accesses, where first-order oracles reveal the gradient information at the query points, defined as follows.

\begin{dfn}[First-order oracle] 
% For a single differentiable objectives $h: \sR^d \rightarrow \sR$, we call
% $\sO_h^{{\rm fo}}: \sR^d \rightarrow \sR$ a (deterministic) first-order oracle for $h$ if $\sO_h^{{\rm fo}}(\vz) = \nabla h(\vz)$. 
Let $f$ and $g$ be two continuous differentiable functions. We define the first-order oracle 
for the bilevel problem $(f,g)$
as the mapping $\sO^{\rm fo}(\vx,\vy) = ( \nabla_x f(\vx,\vy), \nabla_y f(\vx,\vy) ) $.
\end{dfn}

% Besides the above deterministic oracles, we also study the complexity of algorithms that use the following stochastic oracles.
% \begin{dfn}[Incremental first-order oracle] 
% \label{dfn:IFO}
% Let the differentiable function $h: \gZ \rightarrow \sR$ satisfy the finite-sum structure such that $h = (1/n) \sum_{i=1}^n h_i$.  The incremental first-order oracle (IFO) of $\{ h\}_{i=1}^n$ maps any point $\vz \in \gZ$ and an index $i \in \{1,\cdots,n\}$ to $\nabla h_i(\vz)$. 
% % \end{dfn}

% We define first-order oracles for bilevel problems as the tuple of oracles for both upper- and lower-level functions, and then define the first-order algorithm class as follows.
% follow
%  to define the algorithm class as follows.

% Given the above definitions, the oracle for bilevel problems is then naturally defined as the tuple 

% \begin{dfn}[First-order oracle for bilevel problem] \label{dfn:fo-blo}
%  For a bilevel problem $(f,g)$, we define its first-order oracle $\sO = (\sO_f,\sO_g)$,  The definition applies to both deterministic and stochastic oracles. 
% \end{dfn}
%As in \, we prove lower bounds for the algorithm classes defined as follows.

We then define the (deterministic) first-order algorithms as follows, which 
generalizes the algorithm classes defined 
in \citep{agarwal2018lower,arjevani2020second,carmon2020lower,arjevani2023lower}
from a single objective to two objectives.
Our definition  
encompasses almost all existing first-order algorithms for bilevel optimization \citep{liu2022bome,shen2023penalty,kwon2023fully,lu2024first}.

% {\color{red} Update 10/12: extend the lower bound to all deterministic algorithms.}
\begin{dfn}[First-order algorithm] \label{dfn:rand-alg}
% Under both the deterministic and stochastic settings, 
% Let $\sO^{{\rm fo}} = (\sO_f^{{\rm fo}}, \sO_g^{{\rm fo}})$, where 
% where $\sO_f^{{\rm fo}}$ and $\sO_g^{{\rm fo}}$ are the first-order oracles of $f$ and $g$, respectively. 
A first-order algorithm $\texttt{A}$ consists of a sequence of measurable mappings $\{ \texttt{A}^{t}\}_{t \in \sN}$ such that $ \texttt{A}^{t}$ takes in the first $t-1$ oracle responses to product the $t$-th query $\vw^t = (\vx^t,\vy^t)$, such that $\vw^0 = (\vzero,\vzero)$ and  
\begin{align}  \label{eq:rand-seq}
    \vw^t = \texttt{A}^t \left(\sO^{\rm fo}(\vw^0), \cdots, \sO^{\rm fo}(\vw^{t-1}) \right), \quad \forall t \in \sN_+.
\end{align}                                                                             
We denote $\gA^{{\rm fo}}$ as the set of all deterministic algorithms that follow protocol (\ref{eq:rand-seq}).
% A first-order algorithm $\texttt{A}$ consists of a distribution $\sP_r$ over a measurable set $\gR$ and a sequence of measurable mappings $\{ \texttt{A}^{t}\}_{t \in \sN}$ such that $ \texttt{A}^{t}$ takes in the first $t-1$ oracle responses and the random seed $r \in \gR$ to product the $t$-th query $\vz^t$, such that $\vz^0 = \vzero$ and  
% \begin{align}  \label{eq:rand-seq}
%     \vz^t = \texttt{A}^t \left(r, \sO^{\rm fo}(\vz^0), \cdots, \sO^{\rm fo}(\vz^{t-1}) \right), \quad \forall t \in \sN_+.
% \end{align}                                                                            
% We denote $\gA^{\text{rand-fo}}(\sO^{{\rm fo}})$ as the set of all randomized algorithms that follow protocol (\ref{eq:rand-seq}) and $\gA^{\text{det-fo}}(\sO^{\rm fo}) \subseteq \gA^{\text{rand-fo}}(\sO^{\rm fo})$ as the set of all deterministic algorithms that follow protocol~(\ref{eq:rand-seq}) with a fixed $r$. 
\end{dfn}
% \begin{dfn}[First-order zero-respecting algorithm] \label{dfn:zr-alg}
% For given functions $f,g: \sR^{d_x} \times \sR^{d_y} \rightarrow \sR$, we let $\sO_f$ and $\sO_g$ be the first-order oracle for $f$ and $g$, respectively.
% % Let $f,g: \sR^{d_x} \times \sR^{d_y} \rightarrow \sR$ and $\sO= (\sO_f,\sO_g)$, where $\sO_f$ and $\sO_g$ are the first-order oracle for $f$ and $g$, respectively.
% We denote $\gA^{\text{zr}}(\sO_f,\sO_g) \in \gA^{\text{rand}}(\sO_f,\sO_g)$ as the set of all (first-order) zero-respecting algorithms that generate a sequence $\{\vz^t\}_{t=1}^{\infty}$ with $\vz^t = (\vx^t,\vy^t)$ such that $\vz^0 = \vzero$ and 
% \begin{align} \label{eq:zr-seq}
%     {\rm supp}(\vz^{t}) \subseteq  \bigcup_{s < t} \left( {\rm supp}(\sO_f(\vz^s)) 
%     \cup {\rm supp}(\sO_g(\vz^s)) 
%     \right), \quad \forall t \in \sN_+.
% \end{align}
% % \end{dfn}
% Definition \ref{dfn:rand-alg} 

In addition to first-order algorithms, all of our lower bounds also apply to the 
following algorithm class in Definition \ref{dfn:rand-alg-hvp}, which includes many standard HVP-based algorithms \citep{ghadimi2018approximation,ji2021bilevel,hong2023two}. 
These algorithms introduce an additional variable $\vv$ that tracks $-(\nabla_{yy}^2 g(\vx,\vy))^{-1} \nabla_y g(\vx,\vy)$ via gradient updates on the quadratic function $\frac{1}{2} \vv^\top \nabla_{yy}^2 g(\vx,\vy) \vv - \nabla_y g(\vx,\vy)^\top \vv$ and then use the approximate hyper-gradient $\nabla_x f(\vx,\vy) + \nabla_{xy}^2 g(\vx,\vy) \vv$ to update $\vx$. Formally, we define the deterministic HVP oracle and HVP-based methods as follows.

% as formally defined as follows.
 % Furthermore, when restricted to the quadratic $g$ subclass as per Definition \ref{dfn:qg-sub-class}, we can show in Lemma \ref{lem:zr-grad-hvp} that Definition \ref{dfn:zr-alg} also contains the class of  HVP-based zero-respecting algorithms \citep{ghadimi2018approximation,ji2021bilevel,dagreou2022framework,hong2023two,chen2024optimal,arbel2022amortized} as defined below.

% , and the randomized algorithm class is even broader.

% {\color{red} HVP methods.}


\begin{dfn}[HVP oracle] 
% For a single twice differentiable objective $h: \sR^d \rightarrow \sR$, we 
% we call $\sO_h^{\rm hvp}: \sR^d \times \sR^d \rightarrow$ a (deterministic) HVP oracle for $h$ if $\sO_{h}^{\rm hvp}(\vz,\vv) = \nabla^2 h(\vz) \vv$. 
Let $f$ be differentiable and $g$ be twice differentiable. We define the HVP oracle for the bilevel problem $(f,g)$ as the mapping
\begin{align*}
 \sO^{\rm hvp}(\vx,\vy,\vv) = (\nabla_x f(\vx,\vy), \nabla_y f(\vx,\vy), \nabla_{xy}^2 g(\vx,\vy) \vv, \nabla_{yy}^2 g(\vx,\vy) \vv ).
\end{align*}
% \begin{align*}
%     \sO_{f,g}^{\rm hvp}(\vx,\vy,\vv) = ( \sO_f^{\rm fo}(\vx,\vy), \sO_g^{\rm fo}(\vx,\vy),  )
% \end{align*}
% $\sO_{f,g}^{\rm hvp}(\vx,\vy,\vv) = $
% We call $\sO_{h}^{\rm hvp}: \sR^{d} \times \sR^{d_y} \times \sR^{d_y}$ a (deterministic) HVP oracle of a twice differentiable function $g: \sR^{d_x} \times  \sR^{d_y} \rightarrow \sR $ if its output satisfies $\sO_g^{\rm hvp}(\vx,\vy,\vv)= ()$.
% maps any point $\vz \in \sR^d$ to its gradient $\nabla h(\vz)$.
\end{dfn}


\begin{dfn}[HVP-based algorithm] \label{dfn:rand-alg-hvp}
% Given functions $f,g: \sR^{d_x} \times \sR^{d_y} \rightarrow \sR$, we 
%Under the deterministic setting, 
An HVP-based algorithm $\texttt{A}$ consists of a sequence of measurable mappings $\{ \texttt{A}^{t}\}_{t \in \sN}$ such that $ \texttt{A}^{t}$ takes in the first $t-1$ oracle responses to product the $t$-th query $\vw^t = (\vx^t,\vy^t, \vv^t)$, such that $\vw^0 = (\vzero,\vzero, \vzero)$ and  
\begin{align}  \label{eq:rand-seq-hvp}
    \vw^t = \texttt{A}^t \left(\sO^{\rm hvp}(\vw^0), \cdots, \sO^{\rm hvp}(\vw^{t-1}) \right), \quad \forall t \in \sN_+.
\end{align}                                                                             
We denote $\gA^{{\rm hvp}}$ as the set of all deterministic algorithms that follow protocol (\ref{eq:rand-seq}).

% An HVP-based randomized algorithm $\texttt{A}$ consists of a distribution $\sP_r$ over a measurable set $\gR$ and a sequence of measurable mappings $\{ \texttt{A}^{t}\}_{t \in \sN}$ such that $ \texttt{A}^{t}$ product the $t$-th query $(\vx^t,\vy^t,\vv^t)$, such that $\vx^0 = \vy^0 = \vv^0 = \vzero$ and  
% \begin{align}  \label{eq:rand-seq-hvp} 
% \begin{split}
%     (\vx^t,\vy^t,\vv^t) &= \texttt{A}^t \left(r, \sO_{{\rm hvp}}(\vx^0,\vv^0,\vy^0),\cdots,  \sO_{{\rm hvp}}(\vx^{t-1},\vv^{t-1},\vy^{t-1}) \right), \quad \forall t \in \sN_+.
% \end{split}
% \end{align}
% We denote $\gA^{\text{rand-hvp}}(\sO)$ to be the set of all algorithms that follow the follow protocol (\ref{eq:rand-seq-hvp}), and denote $\gA^{\text{det-hvp}}(\sO) \subseteq \gA^{\text{rand-hvp}}(\sO)$ to denote the set of all deterministic algorithms that follow protocol (\ref{eq:rand-seq-hvp}) with a fixed $r$. 
\end{dfn}

% \begin{dfn}[HVP-based zero-respecting algorithms] \label{dfn:zr-alg-hvp}
% % Given functions $f,g: \sR^{d_x} \times \sR^{d_y} \rightarrow \sR$, we
% Let $\sO_f$ and $\sO_g$ be the deterministic first-order oracle for $f$ and $g$, respectively.
% We denote $\gA^{\text{zr-hvp}}(f, g) \in \gA^{\text{rand}}(\sO_f, \sO_g)$ as the set of all HVP-based zero-respecting algorithms that generate a sequence $\{ \vx^t,\vv^t,\vy^t \}_{t=1}^{\infty}$ such that $\vx^0 = \vy^0 = \vv^0 = \vzero$ and 
% \begin{align} \label{eq:zr-seq-hvp}
% \begin{split}
%     \begin{cases}
%     {\rm supp}(\vy^{t})
%     \subseteq
%      \bigcup_{s < t} {\rm supp}\left( \nabla_y g(\vx^s,\vy^s )\right); \\
%      {\rm supp}(\vv^t) \subseteq
%      \bigcup_{s < t} \left( {\rm supp}
%      \left( \nabla_{yy}^2 g(\vx^s,\vy^s) \vv^s \right) \bigcup {\rm supp} \left( \nabla_y g(\vx^s,\vy^s) \right) \right); \\
%     {\rm supp}(\vx^t) \subseteq
%      \bigcup_{s < t} \left( {\rm supp} \left(  \nabla_{xy}^2 g(\vx^s,\vy^s) \vv^s
%     \right) \bigcup {\rm supp} \left( 
%     \nabla_x f(\vx^s,\vy^s) \right) \right).
% \end{cases}
% \quad \forall t \in \sN_+.
% \end{split}
% \end{align}
% \end{dfn}

Definition \ref{dfn:rand-alg-hvp} subsumes the hyper-gradient-based algorithm class previously defined in \citep[Definition 2]{ji2023lower}.
Since the HVP and gradient oracles for quadratic functions are equivalent, proving lower bounds for first-order algorithms on lower-level quadratic problems implies that the same lower bound applies to HVP-based algorithms. Formally, we have the following statement.

%The formal statement is as follows.
% As all the lower bounds in this paper will be proved by constructing hard instances with quadratic lower-level functions, our lower bounds can also apply to HVP-based algorithms because  
% Formally,  when the lower-level function is quadratic, we show below that the
% HVP-based algorithm class in Definition \ref{dfn:rand-alg-hvp} is a subclass of the first-order algorithm class in Definition \ref{dfn:rand-alg}.

%For this subclass, we can show that the HVP-based algorithms in Definition \ref{dfn:rand-alg-hvp} also lies in the first-order algorithm classes in Definition \ref{dfn:rand-alg}.

\begin{lem} \label{lem:zr-grad-hvp}
% Let $\sO_f$ and $\sO_g$ be the deterministic first-order oracle for $f$ and $g$, respectively.
If $g$ is a quadratic function of the form (\ref{eq:quadratic-g}), then
$  \gA^{{\rm hvp}} = \gA^{\rm fo}$.
% the relationships
% \begin{align*}
%     \in \gA^{\text{rand}}(\sO_f,\sO_g) \quad {\rm and} \quad
%     \gA^{\text{det-hvp}}(f,g) \in \gA^{\text{det}}(\sO_f,\sO_g).
% \end{align*}
\end{lem}

\begin{proof} 
If $g$ is a quadratic function of the form (\ref{eq:quadratic-g}), then we have that $\nabla_y g(\vzero,\vzero) = \vb$, $\nabla_{yy}^2 g(\vx,\vy) =  \mH$, and $\nabla_{xy}^2 g(\vx,\vy) = \mJ$. Therefore, 
it holds that
\begin{align} \label{eq:rela-gq}
\begin{split}
\begin{cases}
    \nabla_{yy}^2 g(\vx,\vy) \vv  = \mH \vv = \nabla_y g(\vx,\vv) + \nabla_y g(\vzero,\vzero), \\
    \nabla_{xy}^2 g(\vx,\vy) \vv = \mJ \vv = \nabla_x g(\vx,\vv). 
\end{cases}    
\end{split}
\end{align}
Then,
for any 
sequence $\{(\vx^0,\vy^0,\vv^0),\cdots,(\vx^t,\vy^t,\vv^t) \}$ generated by an
HVP-based algorithm that satisfies Eq. (\ref{eq:rand-seq-hvp}), we know from the above equation that the expanded sequence $\{ (\vx^0,\vy^0), (\vx^0,\vv^0),\cdots, (\vx^t,\vy^t), (\vx^t,\vv^t)  \}$ satisfies Eq. (\ref{eq:rand-seq}), which means that ~$\gA^{\rm hvp} \subseteq \gA^{\rm fo} $.
\textit{Conversely}, for any sequence $\{  (\vx^0,\vy^0), \cdots, (\vx^t,\vy^t)\}$ generated by a first-order algorithm that satisfies Eq. (\ref{eq:rand-seq}), we also
know from Eq. (\ref{eq:rela-gq}) that the expanded sequence $\{ (\vx^0,\vy^0,\vy^0), \cdots, (\vx^t,\vy^t,\vv^t) \}$ satisfies Eq.~(\ref{eq:rand-seq-hvp}), which means  that $\gA^{\rm fo} \subseteq \gA^{\rm hvp}$.

% {\color{red} to prove the converse relationship.}
% from Eq. (\ref{eq:eq-hvp-zo}) we have that
% \begin{align} \label{eq:eq-hvp-zo-gq}
% \begin{split}
%      \begin{cases}
%     {\rm supp}(\vy^{t})
%     \subseteq
%      \bigcup_{s < t} {\rm supp}\left( \nabla_y g(\vx^s,\vy^s )\right); \\
%      {\rm supp}(\vv^t) \subseteq
%      \bigcup_{s < t} {\rm supp}
%      \left( \nabla_y g(\vx^s,\vv^s) + \nabla_y g(\vx^0,\vy^0) + \nabla_y g(\vx^s,\vy^s) \right); \\
%     {\rm supp}(\vx^t) \subseteq
%      \bigcup_{s < t} {\rm supp} \left( 
%     \nabla_x f(\vx^s,\vy^s) +  \nabla_{x} g(\vx^s,\vv^s)
%     \right).
%     \end{cases}
% \end{split}
% \end{align}
% Since ${\rm supp}(\vc+\vd) \subseteq {\rm supp}(\vc) \cup {\rm supp}(\vd)$ holds for any vectors $\vc,\vd \in \sR^d$, we can conclude from Eq. (\ref{eq:eq-hvp-zo-gq}) that 
% % \begin{align*}
% %     \begin{cases}
% %          , & \forall s\in \sN ; \\
% %         , & \forall s\in \sN ; \\
        
% %     \end{cases}
% % \end{align*}
\end{proof}

% Note that the above relationship between HVP-based and first-order algorithms only hold under the deterministic setting. In contrast,
% the complexity of HVP-based and first-order algorithms under the stochastic setting may differ a lot {\color{red} variance} \citep{arjevani2020second,kwon2024complexity}. We leave the lower bounds on the condition number dependency in stochastic HVP-based methods for future study.

% Although our lower bounds in Section \ref{subsec:lb-det} do not apply to any randomized algorithms, 
% The above definition generalized the same definition in \citep[Section 2.2]{carmon2020lower} from a single objective to two objectives $f$ and $g$. 

\subsection{Complexity Measure}
The ultimate goal of bilevel optimization is to optimize the hyper-objective $F(\vx)$ implicitly defined in Eq. (\ref{prob:main-BLO}). Given that $F(\vx)$ is nonconvex in general, finding its approximate minima requires exponential complexity in the worst case \citep{nemirovskij1983problem}. Hence, we follow \citep{carmon2020lower,carmon2021lower,arjevani2023lower} to measure the complexity of an algorithm by the number of oracle calls to find an $\epsilon$-stationary point.

\begin{dfn} \label{dfn:fo-sta}
We call 
a point $\vx \in \sR^d$ an $\epsilon$-stationary point of $F: \sR^{d_x} \rightarrow \sR$ if 
\begin{align*}
    \Vert \nabla F(\vx) \Vert \le \epsilon.
\end{align*}
\end{dfn}

\begin{dfn}
Given a bilevel function class $\gF$  and the associated first-order oracle $\sO$, we define the worst-case complexity of an algorithm class $\gA$ to find an $\epsilon$-stationary point as
\begin{align*}
    {\rm Compl}_{\epsilon}(\gA, \gF, \sO) = \sup_{(f,g) \in \gF} \inf_{\texttt{A} \in \gA} \inf \{t \in \sN \mid \E \Vert \nabla F(\vx^t) \Vert \le \epsilon \},
\end{align*}
where $\{\vz^t\}_{t \in \sN}$ with $\vz^t = (\vx^t,\vy^t)$ is the sequence generated by \texttt{A} running on $(f,g)$.
\end{dfn}


%We aim to lower bound the number of necessary first-order oracle calls to $f(\vx,\vy)$ and $g(\vx,\vy)$ for any zero-respecting algorithms to find an $\epsilon$-stationary point of $F(\vx)$.


% \section{Building Blocks of Our Lower Bound}

%  We also review some known zero-chains for minimization problems, which will be the ``gadges'' for our bilevel lower bound.
% \begin{dfn}[Orthogonal matrix sequence] \label{dfn:orth-matrix}
% We say the matrix sequence $\{ \mQ_i \}_{i=1}^n \in {\rm Orth}(a,b,n)$ if for each $i  =1,\cdots, n$, we have $\mQ_i \in \sR^{a \times b}$ and $\mQ_i \mQ_i^\top = \mI_a$; and for any $1 \le i \ne j \le n$, we have $\mQ_i \mQ_j^\top = \mI_a$. 
% \end{dfn}

% \begin{lem}[{\citet[Lemma 5.1]{zhou2019lower}}] \label{lem:fs-smooth}
% Let $f : \sR^d \rightarrow \sR$ be $L$-smooth and $\{ \mQ_i \}_{i=1}^n \in {\rm Orth}(d,nd,n)$. Then $f^{\rm fs} (\vx)= \sum_{i=1}^n f(\mQ_i \vx) / \sqrt{n}$ is $L$-averaged smooth.
% \end{lem}

% \begin{lem}[Implicit used in \citep{zhou2019lower}] \label{lem:fs-ifo-lem}
% Let $f : \sR^T \rightarrow \sR$ be a zero-chain and define $\{ \mQ_j \}_{j=1}^n \in {\rm Orth}(d,nd,n)$. Define $f^{\rm fs}(\vx) = \sum_{j=1}^n f(\mQ_j\vx) / n$ and let $\sO$ be the associated incremental first-order oracle of $\{ f(\mQ_j\vx) \}_{j=1}^n$. Then for any sequence $\{\vx^t \}_{t=1}^{\infty}$ generated by a zero-respecting algorithm $\texttt{A} \in \gA^{\text{zr}}(\sO)$, there exists an index set $\gJ $ with $\vert \gJ \vert \ge n/2$ such that $(\mQ_j \vx^t)_i = 0$ for all $j \in \gJ$, $i \ge t$, $t \le n T$.
% \end{lem}



% \subsection{Smooth Convex Optimization}

% Given the above preliminaries, we now introduce the well-known zero-chains in minimization problems. 

% \citet{nesterov2018lectures} showed the following chain-like quadratic function leads to the tight $\Omega( \sqrt{L/ \mu} \log (1/\epsilon))$ lower bound for smooth strongly convex optimization.

% \begin{dfn}[Strongly convex zero-chain  \citep{nesterov2018lectures}] \label{dfn:Nes-lb}
% Let $f^{\rm sc}: \sR^{\infty} \rightarrow \sR$ be
% \begin{align*} 
%     f^{\rm sc}(\vx) &:=
%     \frac{L- \mu}{8} \left( (x_1 - 1)^2 + \sum_{i=1}^{\infty} (x_i - x_{i+1})^2 \right) + \frac{\mu}{2} \Vert \vx \Vert^2 \\
%     &= \frac{L- \mu}{8} \vx^\top \mA_{\infty} \vx + \frac{\mu}{2} \Vert \vx \Vert^2,
% \end{align*}    
% where $L> \mu>0$.
% \end{dfn}

% \citet{carmon2021lower} defined a finite version of Nesterov's worst-case function. We also recall this (unscaled) hard function and its large gradient property below.


% \citet{li2021complexity} defined the following smooth strongly convex function as a component of their hard NC-SC minimax instance.
% We also recall this function since it will also be useful in proving our bilevel lower bounds. 

% \begin{lem}[{\citet[Lemma 10]{li2021complexity}}] \label{lem:li-y-star}
% Let $\vy^*(x_1,x_2) = \arg \min_{\vy \in \sR^K} h^{\rm sc}(x_1,x_2/2,\vy)$.
% If $R_y \ge 30 R_x$, for every $x_1,x_2$ such that $\vert x_1 \vert, \vert x_2 \vert \le R_x$, we have $ \Vert \vy^*(x_1,x_2) \Vert_{\infty} \le K R_y$.
% \end{lem}

% \subsection{Smooth Nonconvex Optimization}
% % Lemma \ref{lem:Car-Upsion} shows that $f^{\rm nc}$ is a smooth zero-chain with minimizer $\vx^* = \vone$, while Lemma \ref{lem:Car-large-grad} shows that $\vx$ will have a large gradient unless the last two coordinates are activated. Based on these lemmas and by rescaling $f^{\rm nc}$ so that it meets the requirements of the function class, \citet{carmon2021lower} showed lower bounds for first-order methods applied on functions of different orders of smoothness.

% \subsection{Stochastic Nonconvex Optimization}

% In the stochastic setup, \citet{arjevani2023lower} obtained a tight $\Omega(\epsilon^{-4})$ lower bound for any zero-respecting algorithm to find an $\epsilon$-stationary point of a smooth nonconvex function using the following hard instance and stochastic first-order oracle (SFO).


% \begin{dfn}[Second nonconvex zero-chain \citep{li2021complexity}] \label{dfn:li-NC-chain}
% Let {$\bar f^{\text{nc-2}}: \sR^{2T+1} \rightarrow \sR$} be
% \begin{align*}
%     \bar f^{\text{nc-2}}(\vx)= - \Psi(1) \Phi(x_1)
%     + \sum_{i=1}^{T}
%     H(x_{2i},x_{2i+1}) + 6 \sum_{i=1}^{T} \left( x_{2i-1} - \frac{1}{2} x_{2i} \right)^2,
% \end{align*}
% where $H(x_1,x_2) = \Psi(-x_1) \Phi(-x_2) - \Psi(x_1) \Phi(x_2) $ with the component functions $\Psi, \Phi: \sR \rightarrow \sR$ being  $\Psi(x) =  \vone_{x>1/2} \exp ( 1 -{1}/{(2x-1)^2} ) $ and $\Phi(x) = \sqrt{\rm e} \int_{-\infty}^x \exp\left(- \frac{t^2}{2} \right) {\rm d}t$.
% \end{dfn}

% \begin{lem}[{\citet[Lemma 1] {carmon2020lower}}] \label{lem:Car-Psi-Phi}
% The functions $\Psi, \Phi : \sR \rightarrow \sR$ satisfies 
% \begin{enumerate}
%     \item For all $x \le 1/2$ and $q \in \sN$, $\Psi^{(q)} (x)= 0$.
%     \item For all $x \ge 1$ and $\vert y \vert <1$, $\Psi(x) \Phi'(y) > 1$.
%     \item Both $\Psi$ and $\Phi$ are infinitely differentiable. For all $q \in \sN$, we have
%     \begin{align*}
%         \sup_x \vert \Psi^{(q)}(x) \vert \le \exp (2.5q \log (4q)) \quad {\text{and}} \quad \sup_x \vert \Phi^{(q)}(x) \vert \le \exp ( 1.5q  \log (1.5q)).  
%     \end{align*}
%     \item The function values and derivatives $\Psi,\Psi',\Phi,\Phi'$ are non-negative and bounded, with
%     \begin{align*}
%         0 < \Psi < {\rm e}, ~~ 0< \Psi' < \sqrt{54/}{\rm e}, ~~ 0 < \Phi < \sqrt{2 \pi {\rm e}}, ~~ \text{and} ~~ 0 < \Phi' < \sqrt{\rm e}.
%     \end{align*}
% \end{enumerate}
% \end{lem}


% \begin{lem}[{\citet[Lemma 12]{li2021complexity}}] \label{lem:li-large-grad}
% If $\vert x_{2i} \vert <1$ for some $i \le T$, then $\vx$ is not a $1/3$-stationary point of $\bar f^{\text{nc-2}}$ on the domain $\gC^{2T+1}_{2} := \{\vx \in \sR^{2T+1} \mid \Vert \vx \Vert_{\infty} \le 2 \}$.
% \end{lem}

\section{Lower Bounds} \label{sec:lb}

In this section, we establish new lower bounds for bilevel optimization. We begin in Section~\ref{subsec:framework} by introducing a generic framework for deriving these lower bounds. We then apply this framework to obtain lower bounds across various setups. Specifically, Section~\ref{subsec:lb-det} presents lower bounds for deterministic nonconvex-strongly-convex (NC-SC) and (S)C-SC problems under different orders of smoothness. Section~\ref{subsec:lb-stoc} 
uses a slightly different construction to prove lower bounds for 
randomized and stochastic algorithms in NC-SC problems.

% In this section, we establish new lower bounds for bilevel optimization. 
% % In Section~\ref{sec:zero-chain}, we review the zero-chain as
% % a useful tool for proving lower bounds.
% In Section~\ref{subsec:framework}, we first introduce a generic framework for proving lower bounds in bilevel optimization. Then we apply the framework to obtain lower bounds under various setups, where Section~\ref{subsec:lb-det} gives lower bounds for deterministic NC-SC and 
% (S)C-SC problems under different orders of smoothness, and Section \ref{subsec:lb-stoc} 
% extends the scope of our first-order NC-SC lower bound to randomized and stochastic algorithms.

% \subsection{The Concept of Zero-Chain} \label{sec:zero-chain}



%to randomized and stochastic algorithms.
%give results under the stochastic setting, with extension to randomized algorithms.



\subsection{Construction Overview} \label{subsec:framework}

% Because minimax optimization is a special case of bilevel optimization if $f = -g$, the existing lower bound for NC-SC minimax optimization \citep{zhang2021complexity,li2021complexity,wang2024efficient} also implies the same lower bound for NC-SC bilevel optimization. However, 
% the $\kappa_y$ dependency in both the $\Omega(\sqrt{\kappa_y} \epsilon^{-2})$ lower bound \citep{zhang2021complexity,li2021complexity} and $\Omega(\sqrt{\kappa_y} \epsilon^{12/7})$ lower bound \citep{wang2024efficient} for first- and second-order smooth NC-SC minimax optimization is much smaller than those in the upper bounds for bilevel optimization \citep{chen2023near}.
% In our work, we amplify the $\kappa_y$ dependency 
% in 

We construct hard instances with chain-like structures to prove optimization lower bounds.
Such functions are known as zero-chains, which are 
formally introduced by~\citet{carmon2020lower} as an abstraction of the idea in ~\citep{nesterov2018lectures}. 

% To prove lower bounds, we will . 

%We recall the formal definition below. 


% \begin{lem}[{\citet[Observation 1]{carmon2020lower}}] \label{lem:zero-chain}
% Let $f: \sR^T \rightarrow \sR$ be a zero-chain associated with a first-order oracle $\sO$, and $\{\vx^t \}_{t \in \sN}$ be the sequence generated by a zero-respecting algorithm $\texttt{A} \in \gA^{\text{zr}}(\sO)$. Then $\vx^t_i = 0$ for all $i \ge t$ and $t \le T$. 
% \end{lem}
% For stochastic problems, \citet{arjevani2023lower} introduced the following concept of probability zero-chain to define the function such that its associated stochastic first-order oracle (SFO) only has probability $p$ to activate one coordinate.

% \begin{dfn}[Probability-$p$ zero-chain] A function $f:\sR^T \rightarrow \sR$ with a stochastic first-order oracle $\sO: \vx \mapsto \hat \nabla f(\vx)$ is a probability-$p$ zero-chain if
% \begin{align*}
%     {\rm supp}(\vx) \subseteq \{ 1,\cdots,t-1\} 
%     \Longrightarrow ~~
%     \begin{cases}
%         \sP({\rm supp}(\hat \nabla f(\vx))) \nsubseteq \{1,\cdots,t-1 \}) \le p, \\
%         \sP({\rm supp}(\hat \nabla f(\vx))) \subseteq \{1,\cdots,t\}) =1,
%     \end{cases}
%     \quad \forall \vx \in \sR^T.
% \end{align*}
% \end{dfn}

% \citet{arjevani2023lower} also showed that a length-$T$ probability-$p$ zero-chain imply an $\Omega(T/p)$ lower bound to activate the last coordinate. We restate their result as follows.

% \begin{lem}[{\citet[Lemma 1]{arjevani2023lower}}] \label{lem:pro-zero-chain}
% Let $f: \sR^T \rightarrow \sR$ be a probability-$p$ zero-chain associated with a stochastic first-order oracle $\sO$, and $\{\vx^t \}_{t \in \sN}$ be the sequence generated by a zero-respecting algorithm $\texttt{A} \in \gA^{\text{zr}}(\sO)$. Then with probability at least $1-
% \delta$ we have $\vx^t_i = 0$ for all $i \ge t$ and $t \le (T - \log (1/\delta))/(2p)$. 
% \end{lem}

% Another important observation in \citep{arjevani2023lower} is that: given a zero-chain, there is a generic way to construct probability zero-chains under the stochastic setting. 

% % The next question is: how to construct a probability zero-chain? \citet{arjevani2023lower} gave a clear answer to this question, where they showed how to construct an SFO for a given zero-chain and make them together become a probability zero-chain.

% \begin{lem}[{\citet[Lemma 3]{arjevani2023lower}}] \label{lem:bernoulli-sg}
% Let $f: \sR^T \rightarrow \sR$ be a first-order zero-chain. For $p \in (0,1]$, we define 
% \begin{align*}
%      [\hat \nabla f(\vx)]_i =  
%     [\nabla_z f(\vx) ]_i \left( 1 +
%     \vone_{i= {\rm prog}(\vx) +1} \left( \xi/ p  -1\right) \right),
% \end{align*}
% where $\xi \sim {\rm Bernoulli}(p)$. Then $\sO:\vx \mapsto \hat \nabla f(\vx)$ is a stochastic first-order oracle with variance
% $\sigma^2 \le \Vert \nabla f(\vx) \Vert_{\infty}^2 (1-p)/p $ such that $f$ equipped with $\sO$ is a probability-$p$ zero chain.
% \end{lem}


\begin{dfn}[Zero-chain]
A function $f: \sR^T\rightarrow \sR$ is a (first-order) zero-chain if 
\begin{align*}
    {\rm supp}(\vx) \subseteq \{1,\cdots,t-1 \} ~~ \Longrightarrow ~~ {\rm supp}(\nabla f(\vx)) \subseteq \{1,\cdots,t \}, \quad \forall \vx \in \sR^T.
\end{align*}
\end{dfn}

We say a coordinate is activated if it becomes non-zero from zero. As shown in \citep{arjevani2019oracle,carmon2020lower}, if $f(\vx)$ is a zero-chain, then for any deterministic algorithm~$\texttt{A}$, one can adversarially rotate the $f(\vx)$ such that $\texttt{A}$ running on the rotated function can only activate its coordinates one by one.
This leads to a generic recipe for proving lower bounds: Construct a function $f: \sR^T \rightarrow \sR$ 
that meets two criteria: (1) it is a zero-chain, and (2) a point $\vx \in \sR^T$ cannot be an $\epsilon$-solution unless its $T$-th coordinate is activated.

We apply this framework to prove lower bounds for bilevel problems, with the intention of maximizing the chain length $T$
while ensuring that the constructed instance remains 
in the function class of interest.
To achieve this goal, we 
make novel changes to the constructions in (near)-optimal NC-SC minimax lower bounds 
\citep{li2021complexity,wang2024efficient}, amplifying their 
$\kappa_y$ dependency by leveraging the bilevel structure.

% In the following, we introduce a construction framework that leads to a long chain length in bilevel problems.
% To achieve this, we  
% In the next section, we will introduce a construction framework tailored to bilevel problems, which can induce a long chain length.
% % The zero-chains naturally restrict the progress of any zero-respecting algorithms (Definition \ref{dfn:zr-alg}) since they must ,  as shown in the following lemma.

\begin{figure}[t]
\centering\includegraphics[width=0.8 \linewidth]{fig/BLO_chain.pdf}
    \caption{ (a) The construction for NC-SC minimax lower bound \citep{li2021complexity,wang2024efficient}. (b) The construction for our NC-SC bilevel lower bound. An orange node or arrow represents a coordinate or connection in the upper-level, while the black represent the ones in the lower-level.  A bold node indicates that a nonconvex regularization is added on the corresponding coordinate. 
    }
    \label{fig:BLO-chain}
\end{figure}

First, we recap the constructions for NC-SC minimax problems in \citep{li2021complexity,wang2024efficient}.
Let $\bar f_{\nu, r}^{\rm nc}: \sR^T \rightarrow \sR$ be the NC zero-chain \citep{carmon2021lower} for minimization problems (see Definition \ref{dfn:nc-zo-chain}).
\citet{li2021complexity} set
$\vx,\vz$ as the minimization (upper-level) variables and $\vy$ as the maximization (lower-level) variable, and constructed functions such that 
(1) every progress from $z_i$ to $x_i$ 
requires traversing a lower-level strongly-convex sub-chain with length $K=\Omega(\sqrt{\kappa_y})$; (2) 
the hyper-objective $F(\vx,\vz)= \max_{\vy} f(\vx,\vz,\vy)$ has similar
properties to $\bar f^{\rm nc}_{\nu, r}(\vx): \sR^T \rightarrow \sR$. 
Such a construction leads to a zero-chain with a total length~$KT$.
The analysis in \citep{carmon2021lower} indicates that the largest possible value of~$T$ that ensures $\bar f_{\nu,r}^{\rm nc}(\vzero) - \inf_{\vx \in \sR^T} f(\vx) \le \Delta$ is 
$T = \Omega(\epsilon^{-2})$ and $T = \Omega(\epsilon^{-12/7})$ under first- and second-order smoothness, which leads to the lower bounds of
$\Omega(\sqrt{\kappa_y} \epsilon^{-2})$ \citep{li2021complexity} and $\Omega(\sqrt{\kappa_y} \epsilon^{-12/7})$ \citep{wang2024efficient} for first- and second-order smooth NC-SC minimax problems, respectively.

% % and coupled $\bar f_{\nu, r}^{\rm nc}$ with $T$ strongly convex sub-chains in Eq. (\ref{eq:sc-sub-chain}), such that  $\bar f_{\nu, r}^{\rm nc}$ 

% Then, setting 
%    to be the largest possible length for 
%    -order and second-order smooth NC functions 
%    \citep{carmon2021lower} 
% In addition, setting for second-order smooth NC functions \citep{carmon2021lower} leads to the 
%  lower bound for second-order smooth NC-SC minimax problems \citep{wang2024efficient}.
% Let $K = $ be the length of strongly convex sub-chain in Definition \ref{dfn:li-SC-subchain}.
% \citet{li2021complexity,wang2024efficient} , such that the total chain length is $\Omega(KT)$ and in Definition~\ref{dfn:nc-zo-chain}. Then setting
% $T = \Omega(\epsilon^{-2})$ and $T= \Omega(\epsilon^{-12/7})$ to be the largest possible length for nonconvex objectives \citep{carmon2021lower} leads to the  and  lower bound for first- and second-order smooth NC-SC minimax optimization,
% respectively proven by \citet{li2021complexity} and \citet{wang2024efficient}.

In this work, we modify the previous construction in NC-SC minimax optimization as follows (also illustrated in Figure \ref{fig:BLO-chain}):
\begin{enumerate}
    \item We move $\vz$ from the upper to the lower level, and let the lower-level connection in $\vx$ and $\vz$ lead to $\vz^*(\vx) = K^2 \vx$. This step is possible in bilevel problems because even with the upper-level function $f(\vx,\vz,\vy)$ fixed, we can
   still have the freedom to manipulate the lower-level function $g(\vx,\vz,\vy)$.
    \item We carefully design the connections of $\vx,\vz,\vy$ (the triangles in Figure \ref{fig:BLO-chain} (b)) and add nonconvex regularization on $\vz$ instead of $\vx$, such that the final hyper-objective $F(\vx) = f(\vx, \vz^*(\vx), \vy^*(\vx))$ is equivalent to $\bar f^{\rm nc}_{\nu,r} (K^{3/2} \vx) $ rather than $\bar f^{\rm nc}_{\nu,r} (\vx) $. This step combines a technical lemma from \citep[Lemma 7]{li2021complexity} with a careful modification in both $f$ and $g$ to boost the smoothness of $F(\vx)$ by a factor of $K^3$.
\end{enumerate}
{Our hard instance characterizes two sources of condition number dependencies in the complexity of bilevel problems:
First, we use the variable $
\vy$ to couple sub-chains into the zero-chain in $\vx$, which reflects the $\kappa_y$ dependency from the computation of $\vy^*(\vx)$ in Lemma~\ref{lem:AGD}. Second, we use the variable $\vz$ to rescale $\vx$ to $\kappa_y \vx$, which reflects the $\kappa_y$ dependency on the smoothness constant of the hyper-objective $F(\vx)$ in Lemma \ref{lem:smooth-phi}. }

Using this hard instance, we can prove a set of lower bounds for NC-SC bilevel problems.
If we let $\epsilon^{-\gamma}$ be the length of nonconvex zero-chain under different orders of smoothness \citep{carmon2021lower}, where $\gamma=2$ , $\gamma = 12/7$, and $\gamma = 8/5$ for first-order, second-order, or arbitrarily smooth functions, respectively, then the scaling factor in our hyper-objective means that our construction allows for a larger length $T = \Omega( (K^{3/2}/ \epsilon)^{\gamma})$ in the upper level. Together with $K = \Omega(\sqrt{\kappa_y})$ in the lower level, this implies a lower bound of $\Omega(KT) = \Omega\left(\kappa_y^{3 \gamma /4 + 1/2} \epsilon^{-\gamma} \right)$ for NC-SC bilevel problems. 

% To prove lower bounds for C-SC problems, we simply remove the nonconvex regularizers in the upper-level function of our hard NC-SC instance. It leads to $F(\vx) = \bar f^{c}(K^{3/2}\vx)$, where $\bar f^{\rm c}$ is Nesterov's worst-case convex zero-chain (see Definition \ref{dfn:finite-Nes-func}). The scaling factor in the hyper-objective allows us to set $T = \Omega(K^{3/2} / \sqrt{\epsilon})$ while still 
% ensuring the regular conditions of the objectives. Then, along with $K = \Omega(\sqrt{\kappa_y})$, we obtain the lower bounds of $\Omega(TK) = \Omega(\kappa_y^{5/4} / \sqrt{\epsilon} ) $ for C-SC problems. As for SC-SC problems, we can obtain the lower bound of $\tilde \Omega(\kappa_y^{5/4}  \sqrt{\kappa_x} )$ by adding a regularizer of $\mu_y \Vert \vz \Vert^2/2$ in $F(\vx)$. 

% Finally, we compare a different (but simpler) construction that abandons variable $\vy$ in our previous construction for NC-SC problems. 
% Now, we set $f(\vx,\vz) = \bar f^{\rm nc}(\vz)$, $g(\vx,\vz) = \mu_y \Vert \vz \Vert^2/2 - L_1 \langle \vx, \vz \rangle$, which induces $\vz^*(\vx) = \kappa_y \vx$ and $F(\vx)= \bar f^{\rm nc}(\kappa_y x)$. This second construction now gives a lower bound of $\Omega( (\kappa_y / \epsilon)^\gamma)$ for NC-SC bilevel problems, which matches that of our first one when $\gamma = 2$, but is worse when $\gamma \le 2$.  Although the second construction does not lead to better results for deterministic cases where $\gamma \le 2$, we use it to prove lower bounds for stochastic NC-SC problems where $\gamma=4$ \citep{arjevani2023lower}.

% note that our lower bound for first-order NC-SC problems can be obtained using 


% To prove lower bounds for
% stochastic first-order smooth NC-SC problems, we use   since it induced extra variance in the stochastic setting \citep{li2021complexity}.
%  Then 

% Essentially, we note that the results in  can be obtained via modifying our second construction to the convex setting, which gives $F(\vx) = \bar f^{\rm c}(\kappa_y \vx)$ and leads to .
% to obtain the lower bound of $\Omega()$
% Moreover, we can also apply a similar (but simpler) construction in (S)C-SC problems and prove lower bounds under the corresponding setups.

% to the larger condition number dependency.

% Because nonconvex-strongly-convex bilevel optimization subsumes nonconvex optimization, the lower bound for nonconvex single-level optimization \citep{carmon2020lower,carmon2021lower} already gives the lower bound in terms of the $\epsilon$ dependency. The main focus of this paper is to amplify these results o reveal the $\kappa_y$ dependency of bilevel problems, which comes from the following two sources: 
% First, the smoothness of  hyper-objective $F(\vx) = f(\vx,\vy^*(\vx))$ depends on the smoothness of $\vy^*(\vx)$, which further depends on $\kappa_y$ in general \citep[Lemma 2.2]{ghadimi2018approximation}.
% Second, the expression of hyper-gradient $\nabla F(\vx)$ in Eq. (\ref{eq:hyper-grad}) involves the optimal solution to the lower-level function $\vy^*(\vx)$. This means that each time an iterative algorithm uses $\nabla F(\vx)$, it needs to solve a strongly convex lower-level problem, which requires at least $\Omega(\sqrt{\kappa_y})$ complexity in the worst case \citep[Section 2.1.4]{nesterov2018lectures}.

% According to the above analysis, we design a generic lower bound framework that can recover the same $\epsilon$ dependency in single-level optimization while amplifying the 
% $\kappa_y$ dependency that comes from the above two sources. We let $\vx$ be the upper-level variable, $(\vy,\vz)$ be the lower-level variable, and construct the upper- and lower-level functions $f$ and $g$ such that they jointly form the zero-chain structure illustrated in Figure \ref{fig:BLO-chain}. In the figure, each node represents a coordinate of the variable.  and the arrows indicate the order in which the coordinates are activated, where the solid line and dotted line represent the connection in $f$ and $g$, respectively. The \textbf{black} arrows in the upper level represent the same connections from the nonconvex zero-chain in Definition \ref{dfn:nc-zo-chain}, which forms the backbone with length $2T+1$ and contributes to the $\epsilon$ dependency in our lower bounds. The green and blue arrows in the lower-level contribute to the two different sources of the $\kappa_y$ dependency: the \textbf{\color{darkgreen}green} arrows rescale $\vy_j^{(i)*}(\vx) = \kappa_y^{3/4} \vx_j$ for $i = 1,2$ to boost the smoothness constant of the hyper-objective~$F(\vx)$ and make $T$ larger than in single-level optimization; the \textbf{\color{darkblue}blue} arrows form $T$ sub-chains from Definition \ref{dfn:li-SC-subchain}, each with length $K = \Theta(\sqrt{\kappa_y})$ to reflect the complexity of obtaining $\vy^*(\vx)$ to make progress in the upper level and lead to a total chain length of $\Omega(KT)$.
% The directions of the arrows while their colors indicate the different types of connections and their roles:  arrows are the connections same as \cite{carmon2021lower},  arrows are used to rescale the smoothness constant of $F(\vx)$, and  are used for coupling a convex zero-chain.

% According to the above analysis, we use two corresponding techniques to show our lower bounds. First, we use the lower-level function to rescale $\vy^*(\vx) = \kappa_y^{3/4} \vx$, which boosts the smoothness constant of $F(\vx)$.
% Second, we couple the convex zero-chain and nonconvex zero-chain to ensure that, in order to activate a new coordinate in the upper-level, any algorithm must go through a whole strongly convex lower-level zero-chain with length $K= \Omega(\sqrt{\kappa_y})$ for finding $\vy^*(\vx)$. 

\subsection{Lower Bounds for Deterministic Problems} \label{subsec:lb-det}

In this subsection, we formally prove our lower bounds for deterministic high-order smooth NC-SC 
and (S)C-SC bilevel problems.

\subsubsection{Lower Bounds for Deterministic NC-SC Problems} \label{subsec:lb-det-nc}

% In the nonconvex setup, \citet{carmon2021lower}  as follows, and showed that the resulting function can lead to  

To establish lower bounds for NC-SC bilevel problems, we first review the NC zero-chain introduced by \citet{carmon2021lower}, which 
argumented the convex zero-chain \citep{nesterov2018lectures}
with a separable nonconvex regularizer to obtain lower bounds for finding approximate stationary points of high-order smooth NC functions.

\begin{dfn}[NC zero-chain \citep{carmon2021lower}] \label{dfn:nc-zo-chain}
Let $\bar f_{\nu,r}^{\rm nc}: \sR^{T}\rightarrow \sR$ be
\begin{align*}
    \bar f_{\nu, r}^{\rm nc}(\vx) = \frac{\sqrt{\nu}}{2} (x_1 - 1)^2 + \frac{1}{2} \sum_{i=2}^T (x_{i}-  x_{i-1})^2
    + \nu \sum_{i=1}^{T-1} \Upsilon_r(x_i),
\end{align*}
where the nonconvex regularizer $\Upsilon_r(x) = 120 \int_1^x \frac{t^2 (t-1)}{1+(t/r)^2} {\rm d}t$.
\end{dfn}

We summarize the key properties of this function in the following lemma.

%The first one is about the properties of the regularizer $\Upsilon_r$, and the second on is the large gradient property of the whole function $\bar f_{\nu, r}^{\rm nc}(\vx)$.

\begin{lem}[{\citet[Lemma 2 and 3]{carmon2021lower}}] \label{lem:Car-Upsion}
 The functions $\Upsilon_r$ and $\bar f_{\nu,r}^{\rm nc}$ 
 satisfy the following:
 \begin{enumerate}
 \item We have $\Upsilon_r'(0) = \Upsilon_r'(1) = 0$.
 \item For all $x \in \sR$, we have $\Upsilon_r(x) > \Upsilon_r(1) = 0$, and for all $r$, $\Upsilon_r(0) \le 10$.
    \item For every $r \ge 1$ and every $p \ge 1$, the $p$th-order derivatives of $\Upsilon_r$ are $r^{3-p}\ell_p$-Lipschitz continuous, where $\ell_p \le \exp(3p \log(p)/2 + c_0p)$ for a numerical constant $c_0>0$.
\item  Let $r \ge 1$ and $\nu \le 1$, then for any $\vx \in \sR^T$ such that $x_{T-1} = x_{T}= 0$, we have 
\begin{align*}
    \Vert \nabla \bar f^{\rm nc}_{\nu, r}(\vx) \Vert > \nu^{3/4} / 4.
\end{align*}
 \end{enumerate}
\end{lem}

To extend the above NC hard instance for minimization to an NC-SC hard instance for NC-SC bilevel problems, we set $\vx \in \sR^T$, $(\vz,\vy) \in  \sR^T \times \sR^{K(T-1)} $ be the upper- and lower-level variables, respectively, partition $\vy \in \sR^{(T-1)K}$ into $\vy = (\vy^{(1)},\cdots, \vy^{(T-1)})$ so that $\vy^{(i)} \in \sR^K$, 
% $\vy^{(i)} \in \sR^{2T+1}$ for $i = 1,2$, $\vz = $ with $\vz^{(i)} \in \sR^K$ for all $i = 1,\cdots,T$, 
and define the (unscaled and unrotated) upper- and lower-level functions $\bar f^{\text{nc-sc}}, \bar g^{\rm sc}: \sR^{T} \times \sR^{T} \times \sR^{K(T-1)} \rightarrow \sR $ as
\begin{align} \label{eq:our-f}
\begin{split}
     \bar f^{\text{nc-sc}}(\vx,\vz,\vy)&:= \frac{\sqrt{\nu}}{2} \left( \frac{z_1}{\sqrt{K}} - 1 \right)^2
     + \nu \sum_{i=1}^{T} \Upsilon_r \left( \frac{z_i}{\sqrt{K}} \right) + {\color{blue}h(\vz,\vy)}, \\
\bar g^{\rm sc}(\vx,\vz,\vy) &:= 
{\color{blue} \frac{1}{2K^2} \Vert \vz \Vert^2 - \langle \vx, \vz \rangle} +  \sum_{i=1}^{T-1} h^{\rm sc}(x_{i}, x_{i+1}, \vy^{(i)}),
\end{split}
\end{align}
where $\Upsilon_r: \sR \rightarrow \sR$  follows from Definition \ref{dfn:nc-zo-chain}. For the function $h: \sR^{T} \times \sR^{K(T-1)}: \rightarrow \sR$ in the upper level, we design it as
\begin{align} \label{eq:our-hard-h}
    h(\vz,\vy) := \sum_{i=1}^{T-1} \frac{a_K}{2} (z_i^2+ z_{i+1}^2) +  \frac{b_K}{2} (z_i y_1^{(i)} - z_{i+1} y_K^{(i)}),
\end{align}
where the constants $a_K, b_K \in \sR$ depend on $K$, with specific values provided in Lemma \ref{lem:hyperfunc}.
For the function $h^{\rm sc}: \sR \times \sR \times \sR^K \rightarrow \sR$ in the lower level, we set it as 
\begin{align} \label{eq:sc-sub-chain}
\begin{split}
h^{\rm sc}(x_1,x_2,\vy) &:= \frac{1}{2} \sum_{i=1}^{K-1} (y_i - y_{i+1})^2 + \frac{1}{2K^2} \Vert \vy \Vert^2 - \left( x_1 y_1 - x_2 y_K \right) \\
     &=\frac{1}{2} \vy^\top \left( \frac{1}{K^2} \mI_K + \mA_K \right)
 \vy -  \langle x_1 \ve_1 - x_2 \ve_K, \vy \rangle,
\end{split}
\end{align}
where the matrix $\mA_K \in \sR^{K \times K} $ is
\begin{align} \label{eq:finite-A}
   \mA_K
    =
    {\small
    \begin{bmatrix}
        1 & -1 \\
        -1 & 2 & -1 \\
        & -1 & \ddots & \ddots  \\
        & & \ddots & 2 & -1 \\
        & &  & -1 & 1\\
    \end{bmatrix}}.
\end{align}


% \begin{lem}[{\citet[Lemma 3]{carmon2021lower}}] \label{lem:Car-large-grad}

% \end{lem}

% For convenience, we define t as


% \begin{dfn}[Strongly convex sub-chain \citep{li2021complexity}]
% \label{dfn:li-SC-subchain}
% Let $h^{\rm sc} :\sR \times \sR \times \sR^K \rightarrow \sR$ be
% \begin{align*}
     
% \end{align*}
%where $\vb_{x_1,x_2} = $.
%\end{dfn}


% Specifically, the bilevel problem is 
% \begin{align*}
%     \min_{\vx \in \sR^{2T+1}, \vy \in \sR^{4T+2}, \vz \in \sR^{K T}} \bar f(\vx, \vy,\vz), \quad {\rm s.t.} \quad (\vy,\vz) = \arg \min_{\vy \in \sR^{2T+1}, \vz \in \sR^{K T}} \bar g(\vx,\vy,\vz),
% \end{align*}
% where 
% To rescale the smoothness constant, we let the upper-level function be basically the same as the nonconvex zero-chain $\bar f^{\rm nc}_{\nu, r}$ in Definition \ref{dfn:nc-zo-chain}, but defined on the variable $\vy$ instead of $\vx$. Hence, after plugging in $\vy^*(\vx) = \kappa_y^{3/4} \vx$, the hyper-objective becomes $F(\vx) = \bar f^{\rm nc}_{\nu, r}(\kappa_y^{3/4} \vx)$, which has the larger smoothness constant than the original nonconvex zero-chain. Although the best possible rescaling is $\vy^*(\vx) = \kappa_y \vx$, here we can only have $\vy^*(\vx) = \kappa_y^{3/4} \vx$ in order not to interfere the role of variable~$\vz$.
% To couple the strongly convex zero-chain, we break the connections between two adjacent upper-level coordinates $x_{2i}$ and $x_{2i+1}$ and insert $T$ lower-level strongly-convex sub-chains $h^{\rm sc}(x_{2i},x_{2i+1},\vz^{(i)})$ in Definition \ref{dfn:li-SC-subchain}. Because $x_{2i}$ and $x_{2i+1}$ are only connected through the lower-level sub-chain $h^{\rm sc}(x_{2i},x_{2i+1},\vz^{(i)})$, for any algorithm to make progress in the upper level, it must go through an entire new strongly convex zero-chain in the lower level each step.
% \paragraph{Our construction.} As discussed in the previous section, our novel construction used both the ``inserting'' and ``rescaling'' techniques in a single instance. To achieve this effect, we participate the lower-level variable into two parts and denote it by $(\vy,\vz)$, where the role of $\vy$ is for rescaling and $\vz$ is for ``inserting''. Specifically, 
% we construct our hard instance as follows:
% Now, we present our lower bounds for deterministic bilevel problems. 
% We start from the NC-SC problems, and later show that our construction also implies improved lower bounds for (S)C-SC problems considered in prior work \citep{ji2023lower}.

%Formally, 
% $h^{\rm chain}: \sR^{2T+1} \rightarrow \sR$, $h^{\rm reg}:\sR \times \sR \rightarrow \sR$, and $h^{\rm scale}: \sR^{2T+1} \times \sR^{2T+1} \rightarrow \sR$ are
% \begin{align} \label{eq:our-components-h}
% \begin{split}
%       h^{\rm chain}(\vy) &:=\frac{\sqrt{\nu}}{2} (y_1 - 1)^2 + \frac{1}{2} \sum_{i=1}^T  (y_{2i-1}- y_{2i})^2; \\
%        h^{\rm scale}(\vx,\vy) &:= \frac{1}{2 K^{3/2}} \Vert \vy \Vert^2 -  \langle \vx, \vy \rangle; \\
%      h^{\rm reg}(y_1,y_2, \vz)&: = c_1 y_1^2 + c_1 y_2^2 + c_2 \Vert \vz \Vert^2
% \end{split}
% \end{align}
The strongly-convex sub-chain $h^{\rm sc}$ aligns with those used in \citep{wang2024efficient} except that the affine term is scaled by $\sqrt{K}$. This sub-chain ensures that every progress from $x_1$ to $x_{2}$ needs to pass the chain $ y_1 \rightarrow y_2 \rightarrow \cdots y_K$.
In Eq. (\ref{eq:our-f}), we mark the main changes we made to the hard NC-SC minimax instance \citep{li2021complexity,wang2024efficient} in blue, which are crucial to obtain our improved lower bound on condition number dependency.
In the following lemma, we present the basic properties of our hard instance in Eq. (\ref{eq:our-f}).
%For any given $\kappa_y >0 $, we need to restrict $K = \gO(\sqrt{\kappa_y})$ to ensure the lower-level function is $\mu_y$-strongly convex. In our constructions,  need to depend on $K$, and we . % $h^{\rm sc}: \sR \times \sR \times \sR^{K} \rightarrow \sR$ is the convex zero-chain $f^{\rm sc}$ defined with $x_1$ and $x_2$ coupled at the beginning and the end of the chain, defined as
% \begin{align} \label{eq:our-hsc}
%     h^{\rm sc}(x_1,x_2,\vz) &:= \frac{1}{2} \sum_{i=1}^{K-1} (z_i - z_{i+1})^2 + \frac{1}{2K^2} \Vert \vz \Vert^2 - \left( x_1 z_1 - x_2 z_K \right),
%  \end{align}
% \paragraph{Rationale behind our construction.}
% It is motivated by the nonconvex zero-chain \citep{carmon2021lower} defined in Eq. (\ref{eq:Car-lb-2}) on the variable $\vy^{(1)}$, except we only maintain the connections of $\vy_{2i-1}^{(1)} \rightarrow \vy_{2i}^{(1)}$ for $i =1,\cdots,T$, and we break the original connections of $\vy_{2i}^{(1)} \rightarrow \vy_{2i+1}^{(1)}$ for $i =1,\cdots,T$ is to insert a new zero-chain given by the lower-level function.

\begin{restatable}{lem}{lempro}
\label{lem:property}
Let $\nu \in (0,1]$, $r \ge 1$ and $K \ge 1$.
The hard instance $(\bar f^{\text{nc-sc}}, \bar g^{\rm sc})$ satisfies:
\begin{enumerate}
    \item 
    Let  $\ell_p$ be the constant defined in Lemma \ref{lem:Car-Upsion}.
    For every $p \ge 1$, the $p$th-order derivatives of $\bar f^{\text{nc-sc}}(\vx,\vz,\vy)$ are $\bar \ell_p$-Lipchitz continuous, where 
    \begin{align*}
      \bar \ell_p = 
      \begin{cases}
          1 +|a_K|  + |b_K| + \nu r^2 \ell_1, & p =1 \\
          \nu r^{3-p}  \ell_p, & p \ge 2.
      \end{cases}
    \end{align*}
    And $\bar g^{\rm sc}(\vx,\vz, \vy)$ has $5$-Lipschitz continuous gradients and $0$ higher-order derivatives.    
    % $p$th-order derivatives of  are  $5$-Lipschitz for $p=1$ and $0$
    \item $\bar g^{\rm sc}(\vx,\vz,\vy)$ is  $(1/K^2)$-strongly convex jointly in $(\vz,\vy)$.
    % \item Let $\vx^0 = \vz^0 = \vy^0 = \vzero$. For any $ t \le K(T-1)$, the iterate $(\vx^t,\vz^t,\vy^t)$ generated by any zero-respecting algorithm $\texttt{A} \in \gA^{\text{zr}}(\sO)$ satisfies $x^t_{T-1} = x^t_{T} = 0$, where $\sO$ is the deterministic first-order oracle $(\vx,\vz,\vy) \mapsto (\nabla \bar f^{\text{nc-sc}}(\vx,\vz,\vy), \nabla \bar g^{\rm sc}(\vx,\vz,\vy))$.
\end{enumerate}
\end{restatable}

\begin{proof}
In our construction, both the upper- and lower-level functions are quadratic except for the nonconvex regularizer~$\Upsilon_r$. Therefore, we can separately bound the Lipschitz constant of the quadratic and non-quadratic parts. For the quadratic part, the Lipschitz constant of the gradient can be bounded by the spectral norm of the Hessian, and all of the higher-order derivatives are zero. For the non-quadratic part, 
we know from Lemma~\ref{lem:Car-Upsion} that its $p$th-order derivative of 
$\Upsilon_r$ is $r^{3-p} \ell_p$-Lipschitz continuous.
This proves items 1 and~2.
Finally, item~3 immediately follows from the form of Eq. (\ref{eq:sc-sub-chain}).
%Finally, item 4 follows from Lemma \ref{lem:zero-chain} and the observation that any zero-respecting algorithm 
% only activate the coordinate following the order indicated by the arrows in Figure \ref{fig:BLO-chain} (b).
% Finally, we prove item 4 by observing the following facts:\begin{enumerate} \item At the first iteration, only $y_1^{(1)}$ is activated by connection $(y_1^{(1)} - 1)^2$ in $f$.\item For all $i = 1,\cdots,2T+1$, $y_i^{(2)}$ is activated after $x_i$ by connection  $x_i y_i^{(2)}$ in $g_1$.
% \item For all $i = 2,\cdots,2T+1$, $y_i^{(1)}$  is activated after $y_{i-1}^{(1)}$ by  connection $(y_i^{(1)} -y_{i-1}^{(1)})^2$ in $f$.
% \item For all $i = 1, \cdots,T$, $x_{2i}$ is activated after $y_{2i}^{(1)}$ by connection $ x_{2i} y_{2i}^{(1)} $ in $g_1$.
% \item For all $i = 1,\cdots, T$, $z^{(i)}_1$ is activated after $x_{2i}$ by connection $x_{2i} z^{(i)}_1$ in $g_2$.
% \item For all $i=1,\cdots, T$, $j=2,\cdots,K$, $\vz^{(i)}_j$ is activated after $\vz^{(i)}_{j-1}$ by $(\vz^{(i)}_j - \vz^{(i)}_{j-1})^2$ in $g_2$.
% \item For all $i = 1,\cdots, T$, $x_{2i+1}$ is activated after $z^{(i)}_n$  by connection $x_{2i+1} z^{(i)}_n$ in $g_2$.
% \item  For all $i = 1,\cdots, T$, $y_{2i+1}^{(1)}$ is activated after $x_{2i+1}$ by connection $x_{2i+1} y_{2i+1}^{(1)}$ in $g_1$.
% \end{enumerate}
% The activation of coordinates should follow the order described above, as we visualize in Figure \ref{fig:BLO-chain}.
\end{proof}

Let $(\vz^*(\vx), \vy^*(\vx)) = \arg \min_{\vz \in \sR^{T}, \vy \in \sR^{K (T-1)}} \bar g^{\rm sc}(\vx,\vz,\vy)$ be the lower-level solution. We carefully choose the coefficients $\alpha_K$ and $\beta_K$ in the function $h(\vz,\vy)$ such that $H(\vx): = h(\vz^*(\vx), \vy^*(\vx))$ admits a simple quadratic form up to scaling.

% as we will soon show in Lemma \ref{lem:hyperfunc}. 
% Note that our upper-level function only has half of the connections in the nonconvex zero-chain $\bar f^{\rm nc}_{\nu,r}$ in Definition \ref{dfn:nc-zo-chain}. The other half of the connections are recovered 
% by plugging in the optimal lower-level solution set $(\vy^*(\vx), \vz^*(\vx))$ to the regularizer $h^{\rm reg}$. We formally show this in the following lemma.

% We can use Lemma \ref{lem:li} to prove the following Lemma \ref{lem:hyperfunc}.
% Now, we present Lemma \ref{lem:hyperfunc}.

\begin{restatable}{lem}{lemhyperfunc}  \label{lem:hyperfunc}
If $K \ge 10$, then there exists some values $a_K, b_K$ satisfying $\vert a_K \vert \le 20$ and $\vert b_K \vert \le 10$ such that 
$H(\vx) = h(\vz^*(\vx), \vy^*(\vx))$ is equivalent to
\begin{align} \label{eq:Hx-goal}
    H(\vx) = \frac{K^3}{2} \sum_{i=1}^{T-1} (x_i - x_{i+1})^2.
\end{align}
\end{restatable}

\begin{proof}
    See Appendix \ref{apx:proof-lem-hyperfunc} for the complete proof.
\end{proof}
% \begin{cor} \label{cor:hyperfunc}
% The hyper-objective $\bar F^{\text{nc-sc}}(\vx) := \bar f^{\text{nc-sc}}(\vx, \vz^*(\vx), \vy^*(\vx))$ is 
% \begin{align*}
%     \bar F^{\text{nc-sc}}(\vx)  = \frac{\sqrt{\nu}}{2} ( \hat x_1 - 1)^2 + \frac{1}{2} \sum_{i=2}^{2T+1} (\hat x_{i} - \hat x_{i-1})^2 +  \nu \sum_{i=1}^{2T} \Upsilon_r (\hat x_i)=  \bar f^{\rm nc}_{\nu, r}(\hat \vx),
% \end{align*}
% where $\hat \vx = K^{3/2} \vx$ and  $\bar f_{\nu, r}^{\rm nc}:\sR^{2T+1} \rightarrow \sR$ the nonconvex zero-chain in Definition \ref{dfn:nc-zo-chain}.
% \end{cor}

% Plugging in the above expression, the resulting hyper-objective becomes exactly the nonconvex zero-chain in Definition \ref{dfn:nc-zo-chain} up to scaling.
Using Lemma \ref{lem:hyperfunc} in Eq. (\ref{eq:our-f}), we can observe that the hyper-objective is
\begin{align*}
\bar F^{\rm nc}(\vx) := \bar f^{\text{nc-sc}}(\vx,\vz^*(\vx), \vy^*(\vx))= \bar f^{\rm nc}_{\nu, r}(K^{3/2} \vx),
\end{align*}
where $\bar f^{\rm nc}_{\nu, r}: \sR^T \rightarrow \sR$ is the NC zero-chain in Definition \ref{dfn:nc-zo-chain}. 
Now, we can rescale and rotate the hard instance as in Eq. (\ref{eq:scaled-fg}) and (\ref{eq:dfn-matrix-P}), then apply a similar analysis to that in~\citep[Theorem 2]{carmon2021lower}
to prove lower bounds for NC-SC bilevel problems under different orders of smoothness.

%Then, we can prove  by rescaling and rotating the hard instance in Eq. (\ref{eq:scaled-fg}) and  by choosing different $\nu$ and $r$.

%[TODO] xxx 
%By taking , we also obtain the corresponding results 
%We first state them together in the following theorem whose proofs will be provided soon.
%with a proper scaling, such that , we know it suffices to show the following lemma.

\begin{thm}[NC-SC lower bound] \label{thm:NC-SC}
There exists a numerical constants $a_0 \in (0,1)$ such that, 
for any $p \in \sN_+$, $L_1,\cdots,L_p>0$, $\Delta >0$ and $\mu_y \in (0, a_0 L_1]$, there exists a lower-level-quadratic bilevel problem $(f, g) \in \gF^{\text{nc-scq}}(L_1,\cdots,L_p, \mu_y, \Delta)$ such that, to find an $\epsilon$-stationary point of the  hyper-objective  $F(\vx)$, any deterministic first-order algorithm $\texttt{A} \in \gA^{\rm fo}$ or HVP-based algorithm $\texttt{A} \in \gA^{\rm hvp}$ requires 
    \begin{align*}
        \begin{cases}
        \vspace{2mm} 
             \Omega \left( \kappa_y^{2} 
             L_1 \Delta \epsilon^{-2} \right), &p=1; \\
             \vspace{2mm} 
             \Omega \left( \kappa_y^{25/14} (L_1)^{3/7} ( L_2)^{2/7} \Delta \epsilon^{-12/7}    \right), & p =2; \\ 
             \Omega \left( \kappa_y^{17/10} (L_1)^{3/5} L_*^{1/5} \Delta \epsilon^{-8/5} \right), & p \ge 3,
        \end{cases}
    \end{align*}
    calls of first-order oracles when $\epsilon \rightarrow 0$,
    where $\kappa_y = L_1/\mu_y$ is the condition number of the lower-level problem, $L_* = \min_{q=2, \cdots, p}  \left( {L_q}/{L_1 } \right)^{2/(q-1)}$.
\end{thm}

\begin{proof}
    See Appendix \ref{apx:proof-thm-SCNC} for the complete proof.
\end{proof}

% Now, let us compare our results with known lower bounds in NC-SC minimax problems. \citet{li2021complexity,zhang2021complexity} established an $\Omega(\sqrt{\kappa_y} \epsilon^{-2})$ lower bound for first-order smooth NC-SC minimax problems, and \citet{wang2024efficient} established an $\Omega(\sqrt{\kappa_y} \epsilon^{-12/7})$ for second-order NC-SC minimax problems. For arbitrarily smooth NC-SC minimax problems $(p \ge 3)$, there are no known lower bounds, but we expect that the analysis for $p=2$ \citep{wang2024efficient} can be extended to such a setting and gives an $\Omega(\sqrt{\kappa_y} \epsilon^{-8/5})$ lower bound.

Compared with known lower bounds for NC-SC minimax problems whose optimal
condition number dependency is both $\sqrt{\kappa_y}$ for $p=1$
\citep{li2021complexity,zhang2021complexity} and $p=2$ \citep{wang2024efficient}, our lower bounds for NC-SC bilevel problems exhibit much larger~$\kappa_y$ dependencies. Our results show that NC-SC bilevel optimization is provably more challenging than NC-SC minimax optimization, with the increased $\kappa_y$ dependency reflecting a fundamental difficulty stemming from the bilevel problem structure.

% Comparing our lower bounds with the known lower bounds for NC-SC minimax problems, 
%{\color{red} gap between minimax problems.}
% In the following, we compare the lower bounds in Theorem \ref{thm:NC-SC} with the upper bounds established in recent work \citep{chen2023near} for the cases $p=1$ and $p=2$. We are not aware of any upper bounds for the case $p=3$, but it may possibly be obtained by applying a similar technique to \citep[Theorem 2]{carmon2017convex}. We leave it to future work.

% \begin{remark}[Known NC-SC upper bounds.] 
% The best-known upper bounds given by fully first-order methods are $\tilde \gO(\bar \kappa_y^4 \epsilon^{-2})$ and $\tilde \gO(\bar \kappa_y^{15/4} \epsilon^{-7/4})$ for $p=1$ and $p=2$, respectively \citep{chen2023near}, where $\bar \kappa_y: = \bar L / \mu_y$ is the global condition number. It suggests that a significant gap in the condition number still exists for the current results. 
% \end{remark}

% Now, we formally prove Theorem \ref{thm:NC-SC}.

% C-SC/SC-SC bilevel problems that apply to first-order zero-respecting algorithms (Definition \ref{dfn:zr-alg}).
% , and lower bounds for stochastic first-order smooth NC-SC bilevel problems that apply to first-order randomized algorithms (Definition \ref{dfn:rand-alg}) .  

%The lower bounds in Theorem \ref{thm:NC-SC} have the same $\epsilon$ dependency as single-level nonconvex optimization \citep{carmon2021lower}, and obtain a much larger $\kappa_y$ dependency compared to all existing lower bounds in NC-SC minimax optimization,w which is $\Omega(\sqrt{\kappa_y} \epsilon^{-2})$ for $p=1$ \citep{li2021complexity,zhang2021complexity} and $\Omega(\sqrt{\kappa_y} \epsilon^{-12/7})$ for $p=2$ \citep{wang2024efficient}.This theorem shows that NC-SC bilevel problems are provably more challenging than NC-SC minimax problems, especially in the $\kappa_y$ dependency. 

\subsubsection{Lower Bounds for Deterministic C-SC and SC-SC Problems} \label{subsec:lb-det-sc}

\citet{ji2023lower} proved lower bounds of $\tilde \Omega(\min \{ \kappa_y, \epsilon^{-3/2} \} / \sqrt{\epsilon}  )$ and
$\tilde \Omega (\kappa_y \sqrt{\kappa_x})$   
for the convex-strongly-convex (C-SC) problems and
strongly-convex-strongly-convex (SC-SC) in Definition \ref{dfn:SC-SC-func} and \ref{dfn:SC-SC-func-quadra}, which demonstrates the first larger $\kappa_y$ dependency than the lower bounds of  $\Omega( \sqrt{\kappa_y/ \epsilon} )$ and $\tilde \Omega(\sqrt{\kappa_y \kappa_x})$ for C-SC and SC-SC minimax problems \citep{ouyang2021lower,zhang2022lower}.
However, the optimality of the lower bounds in \citep{ji2023lower} remains unknown, as they still have a gap between the upper bounds of $\tilde \gO(\bar \kappa_y^2 / \sqrt{\epsilon})$ and $\tilde \gO(\bar \kappa_y^2 \sqrt{\kappa_x})$ established in~\citep{ji2023lower}.

In this subsection, we show that our hard NC-SC instance in the previous section can also be extended to C-SC and SC-SC problems \citep{ghadimi2018approximation,ji2023lower} and achieve improved lower bounds of $\Omega(\kappa_y^{5/4} /\sqrt{\epsilon} )  $ and $\tilde \Omega( \kappa_y^{5/4} \sqrt{\kappa_x} )$, respectively. To begin with, we formally define smooth and quadratic (S)C-SC bilevel problems as follows.

 % upper bounds
% In the prior work, 
% \citet{ji2023lower} did not establish lower bounds for the more common NC-SC bilevel problems, but they instead consider the following  below, where 
% the hyper-objective $F(\vx)$ is assumed to be 
% (strongly) convex. Although the main focus of this paper is NC-SC problems, we also 
% provide lower bounds for (S)C-SC problems for better comparing our construction with \citep{ji2023lower}. 
% instead of (S)C-SC problems since we do not know any sufficient conditions to guarantee the convexity of the hyper-objective, as also pointed out by \citet{hong2023two}. 
% We 
% convexity of 
% We notice that the following more restricted function classes are also considered in the literature \citep{ji2023lower}. We also list them below.

\begin{dfn}[Smooth (S)C-SC problems] \label{dfn:SC-SC-func}
Given
$L_0,L_1, L_2>0$, $\mu_x \in [0,L_1]$, $\mu_y \in (0,L_1]$, and $D>0$, we use $\gF^{\text{sc-sc}}(L_0,L_1,L_2,\mu_x,\mu_y, D)$ to denote the subset of all bilevel problems $(f,g)$ 
that satisfies items 1 and 2 in Definition \ref{dfn:NC-SC-func}, \textit{i.e.}, 
$\gF^{\text{sc-sc}}(L_0,L_1,L_2,\mu_x,\mu_y, D)$ $ \subset $ $\gF^{\text{nc-sc}} (L_0,L_1,L_2, \mu_y, +\infty )$, and the hyper-objective
$F(\vx)$ also satisfies
\begin{enumerate}
    \item $F(\vx)$ is $\mu_x$-strongly convex in $\vx$. 
    \item $\min\{ \Vert \vx \Vert : \vx \in  \arg \min_{\vx \in \sR^{d_x}} F(\vx)  \} \le D$.
\end{enumerate}
If $\mu_x = 0$, we call them C-SC problems: $\gF^{\text{c-sc}}(L_0,L_1,L_2,\mu_y, D) = \gF^{\text{sc-sc}}(L_0,L_1,L_2,0,\mu_y, D)$.
\end{dfn}

{  We note that the (S)C-SC problems are not defined by the convexity of the upper-level function $f$, but instead by the convexity of the hyper-objective function $F$, making the definition potentially difficult to check in practice, as discussed in~\citep{hong2023two}. }

% \begin{dfn}[Smooth C-SC problems]  \label{dfn:C-SC-func}
% Given
% $L_0,L_1, L_2>0$,  $\mu_y \in (0,L_1]$, $D>0$, we denote .
% \end{dfn}

\begin{dfn}[Quadratic (S)C-SC problems]  \label{dfn:SC-SC-func-quadra}
Given $L_1,L_2>0$, $\mu_x,\mu_y \in (0,L_1]$, and $D>0$, we denote
$\gF^{\text{scq-scq}}(L_1,\mu_x,\mu_y, D) \subseteq \gF^{\text{sc-sc}}(+\infty, L_1, 0 ,\mu_x,\mu_y, D)$ as the subset of SC-SC problems such that both the upper- and lower-level functions are quadratic. If $\mu_x=0$, we call them quadratic C-SC problems: $\gF^\text{cq-scq}(L_1,\mu_y,D) = \gF^\text{scq-scq}(L_1,0,\mu_y,D)$.
% We also denote .
% % $\gF^{\text{cq-scq}}(L_1,\mu_y, D) \subseteq \gF^{\text{c-sc}}(+\infty, L_1,0, \mu_y, D)$
% % as the subclasses 
% % \begin{align*}
%     %  , \\
%     % 
% \end{align*}
\end{dfn}

To see how our NC-SC hard instance can be modified to the C-SC settings, we first note that if the NC regularizers in the NC zero-chain $\bar f_{\nu,r}^{\rm nc}$ in Definition \ref{dfn:nc-zo-chain} are removed, the function becomes a finite version of Nesterov's convex zero-chain in Definition \ref{dfn:finite-Nes-func} and has the following property stated in Lemma \ref{lem:Car-large-grad-convex}.

%{\color{red} Add definition of the quadratic subclass.}
% Recall our hard NC-SC bilevel instance in 
% Eq. (\ref{eq:our-f}). To prove lower bounds for C-SC and SC-SC problems, we simply remove the 
% NC regularizer $\nu \sum_{i=1}^{T} \Upsilon_r (z_i / \sqrt{K})$ in the NC upper-level function and obtain a quadratic C-SC hard instance. 
% Next, we recall the properties of the following convex zero-chain \citep{carmon2021lower}.

% We will show that, by
% simply removing the 
% in , the same hard instance also leads to the lower bounds of 
% $\Omega (\kappa_y^{5/4} /\sqrt{\epsilon} ) $ and $\Omega (\kappa_y^{5/4} \sqrt{\kappa_x} ) $ 
% for C-SC and SC-SC problems, respectively.


%$when the NC regularizer in $\bar f_{\nu, r}^{\rm nc}$ is removed.

% Before we give the proof of Theorem \ref{thm:NC-SC}, we present two direct corollaries for SC-SC and C-SC problems, which are obtained by simply replacing the 
% %Although we only state the results for first-order smooth (S)C-SC problems, our lower bound is also valid under higher-order smoothness assumptions sincethe functions we use are fully quadratic and have zero higher-order derivatives. 
% Our results improve prior lower bounds , showcasing the advantage of our novel construction in 

% In this subsection, we prove lower bounds for convex-strongly-convex (C-SC) and strongly-convex-strongly-convex (SC-SC) bilevel problems, where we additionally assume the hyper-objective $F(\vx)$ is convex or strongly convex. The study of this setting dates back to \citet{ghadimi2018approximation}, and is later followed by \citet{hong2023two,ji2023lower}. However, prior works provide very few examples provably satisfying this condition. In Appendix \ref{apx:exmple-C-SC}, we show that this setting covers many useful applications such as meta-learning, robust optimization, and graph neural network training, 
% verifying the rationality of studying this setup.

% \citet{zhang2021complexity,ouyang2021lower} proved the $\Omega(\sqrt{\kappa_y/ \epsilon})$ and $\Omega(\sqrt{\kappa_y \kappa_x})$ lower bound for C-SC and SC-SC minimax optimization, which also applies to bilevel optimization. \citet{ji2023lower} extended the construction of \citet{zhang2021complexity} to the bilevel setup, and showing a larger lower bound of $\tilde \Omega( \min\{\kappa_y, \epsilon^{-3/2} \} \epsilon^{-1/2})$ and $\tilde \Omega(\kappa_y \sqrt{\kappa_x})$ for the corresponding setups. Below, we further improve the lower bound to $\Omega(\kappa_y^{5/4} \epsilon^{-1/2}  )$ and $\Omega(\kappa_y^{1/4} \sqrt{\kappa_x})$ following a similar (but simpler) construction as our NC-SC lower bound.

% For the C-SC setup, we take $\sigma=0$, while for the SC-SC setup we take $\sigma>0$. We can also apply Lemma
% \ref{lem:hyperfunc} to show that the (unscaled) hyper-objective $\bar F^{\text{sc-sc}}(\vx) = \bar f^{\text{sc-sc}}(\vx,\vy^*(\vx), \vz^*(\vx))$ is 
% \begin{align} \label{eq:bar-varphi-scsc}
%     \bar F^{\text{sc-sc}}(\vx) = \bar f^{\rm c}_\nu(\hat \vx) + \frac{\sigma}{2} \Vert \vx \Vert^2,
% \end{align}
% where $\bar f^{\rm c}_\nu$ is the convex zero-chain in Definition \ref{dfn:finite-Nes-func}, and $\hat \vx = K^{3/2} \vx$ is the rescaled variable.
% Finally, we can obtain the following improved lower bounds for 
% SC-SC and C-SC bilevel problems.


\begin{dfn}[Convex zero-chain \citep{carmon2021lower}] \label{dfn:finite-Nes-func}
Let $\bar f^{\rm c}: \sR^T \rightarrow \sR$ be the nonconvex function $\bar f_{1,0}^{\rm nc}$ in Definition \ref{dfn:nc-zo-chain} with the NC regularizer $\Upsilon_r$ removed, \textit{i.e.},
\begin{align*}
    \bar f^{\rm c}(\vx) 
    % &:=\frac{\sqrt{\nu}}{2} (x_1 - 1)^2 + \frac{1}{2} \sum_{i=2}^T (x_{i}-  x_{i-1})^2 + \nu \sum_{i=1}^{T-1} \Upsilon_1(x_i) \\
    &:= 
    \frac{1}{2} \vx^\top \left( \mA_T  +  \ve_1 \ve_1^\top \right) \vx - \langle \ve_1, \vx \rangle + \frac{1}{2}.    % \frac{\sqrt{\nu}}{2} (x_1 - 1)^2 + \frac{1}{2} \sum_{i=2}^T (x_{i}-  x_{i-1})^2.
\end{align*}
\end{dfn}

\begin{lem}[{\citet[Lemma 1]{carmon2021lower}}] \label{lem:Car-large-grad-convex}
For any $\vx \in \sR^T$ such that $x_{T}= 0$, 
\begin{align*}
    \Vert \nabla \bar f^{\rm c}(\vx) \Vert > 1/ T^{3/2}.
\end{align*}
\end{lem}

Now, if we remove the NC regularizer term $\nu \sum_{i=1}^{T} \Upsilon_r ( {z_i}/{\sqrt{K}} )$ in the upper-level function in Eq. (\ref{eq:our-f}) and keep the lower-level function unchanged, we know that the resulting (unscaled) hyper-objective is $\bar F(\vx) = \bar f^{\rm c}(K^{3/2}\vx)$. 
Then, we can apply similar lower bound arguments
in convex and strongly-convex minimization problems \citep{nesterov2018lectures,carmon2021lower} to prove the following result.

% Now, we sketch our proof as follows.
% Then, using similar lower bound  as , we can easily prove
% the following lower bounds for C-SC and SC-SC bilevel problems.

% \begin{align} \label{eq:scaled-fg-again}
%    f(\vx,\vy,\vz) = \frac{L_1 \beta^2}{\bar \ell_1} \bar f \left( \frac{\vx}{\beta}, \frac{\vy}{\beta}, \frac{\vz}{\beta} \right), \quad  g(\vx,\vy,\vz) = \frac{L_1 \beta^2}{\bar \ell_1} \bar g \left( \frac{\vx}{\beta}, \frac{\vy}{\beta}, \frac{\vz}{\beta} \right),
% \end{align}
% We further set  $\lambda = \mu_x \bar \ell_1/ L_1$ to ensure that $F(\vx)$ is $\mu_x$-strongly-convex.

\begin{thm}[C-SC lower bound] \label{thm:C-SC}
There exists a numerical constants $a_0 \in (0,1)$ such that, 
for any $ L_1>0$, $\Delta >0$ and $\mu_y \in (0, a_0  L_1]$, there exists a quadratic   bilevel problem $(f, g) \in \gF^{\text{cq-scq}}(L_1, \mu_y, \Delta)$, such that, to find an $\epsilon$-stationary point of the  hyper-objective  $F(\vx)$, any deterministic  first-order algorithm $\texttt{A} \in \gA^{\rm fo}$ or HVP-based algorithm $\texttt{A} \in \gA^{\rm hvp}$ requires at least
    \begin{align*}
        \Omega\left( \kappa_y^{5/4} \sqrt{  L_1  D/ \epsilon} \right)
    \end{align*}
calls of first-order oracle when $\epsilon \rightarrow 0$, where $\kappa_y = L_1/\mu_y$.
\end{thm}

\begin{proof}
    See Appendix \ref{apx:proof-thm-CSC} for the complete proof.
\end{proof}

For SC-SC problems, we slightly modify the hard instance for C-SC problems to let $\bar F(\vx) = \bar f^{\rm c} (K^{3/2} \vx) + \mu_x \Vert \vx \Vert^2/2$, which implies the following result by careful calculation. 

\begin{thm}[SC-SC lower bound] \label{thm:SC-SC}
There exists a numerical constants $a_0 \in (0,1)$ such that, 
for any $ L_1>0$, $\Delta >0$, $\mu_x \in (0, L_1]$, $\mu_y \in (0, a_0 L_1]$, there exists a quadratic bilevel problem $( f, g) \in \gF^{\text{scq-scq}}(L_1, \mu_x, \mu_y)$ such that, to find an $\epsilon$-stationary point of the  hyper-objective  $F(\vx)$, any deterministic first-order algorithm $\texttt{A} \in \gA^{\rm fo}$ or HVP-based algorithm $\texttt{A} \in \gA^{\rm hvp}$ requires at least
    \begin{align*}
        \Omega \left( \kappa_y^{5/4} 
 \sqrt{\kappa_x} \ln \left( \frac{\mu_x D^2}{\epsilon}  \right) \right)
    \end{align*}
calls of first-order oracle when $\epsilon \rightarrow 0$, where $\kappa_x = L_1/\mu_x$ and $\kappa_y = L_1/\mu_y$.
%deterministic first-order oracle calls, where and $\kappa_y = L_1/ \mu_y$.
\end{thm}

\begin{proof}
    See Appendix \ref{apx:proof-thm-SCSC} for the complete proof.
\end{proof}

% Our lower bounds of $\Omega(\kappa_y^{5/4} /\sqrt{\epsilon} )  $ and $\tilde \Omega( \kappa_y^{5/4} \sqrt{\kappa_x} )$ for C-SC and SC-SC problems improve the previous result \citep{ji2021bilevel} by a factor of $\kappa_y^{1/4}$. 

As all functions in our construction are quadratic, these lower bounds also apply to objectives with higher-order smoothness. 
Unlike NC-SC problems, higher-order smoothness
does not provide faster convergence for first-order methods under the (S)C-SC setting. % can ot help first-order methods converge faster for C-SC and SC-SC problems.

{ 
\paragraph{Comparison to \citep{ji2023lower}.} 
Theorem \ref{thm:C-SC} and \ref{thm:SC-SC} improve the previous lower bounds of $\tilde \Omega(\min \{ \kappa_y, \epsilon^{-3/2} \} / \sqrt{\epsilon}  )$ and 
$\tilde \Omega (\kappa_y \sqrt{\kappa_x})$
for C-SC and SC-SC  problems \citep{ji2023lower} by a factor of $\kappa_y^{1/4}$. The improvement comes from the two-source characterization of~$\kappa_y$ dependency in our hard instance, where we jointly use the rescaling variable~$\vz$ to capture the $\kappa_y$ dependency in the smoothness of the hyper-objective~$F(\vx)$, and the coupling variable $\vy$ to capture the $\kappa_y$ dependency $\vy^*(\vx)$ in the computation of~$\vy^*(\vx)$. In contrast, the lower-level variable of the hard instance in \citep{ji2023lower} only plays a role as rescaling, making
$\bar F(\vx)$ behave like $\bar f^{\rm c}(\kappa_y \vx)$, but does not couple additional sub-chains to give a better bound.
Next, we turn to discuss the gap between our lower bounds and upper bounds established in \citep{ji2023lower} for (S)C-SC problems.}
% Essentially, ignoring the additional efforts to tackle unbalanced smoothness using ideas borrowed from \citep{zhang2022lower},
% the hard instance in 
% \citep{ji2023lower} plays a similar role as our second construction, which gives $F(\vx) = \bar f^{\rm c}(\kappa_y \vx)$ and can induce a lower bound of $\Omega(\kappa_y / \sqrt{\epsilon})$ for C-SC problems. However, similar to the NC-SC setting, this instance is worse than our first construction since it corresponds to a setting where the length factor $\gamma$ is small.



\begin{remark}[Known C-SC and SC-SC upper bounds]
\citet{ji2023lower} proposed HVP-based methods with the $\tilde \gO(\bar \kappa_y^2 / \sqrt{\epsilon})$ and $\tilde \gO(\bar \kappa_y^2 \sqrt{\kappa_x})$ upper bounds for C-SC and SC-SC bilevel problems, respectively. In addition, for lower-level quadratic problems, their upper bounds can be refined to 
$\tilde \gO(\bar \kappa_y^{3/2} / \sqrt{\epsilon})$ and $\tilde \gO(\bar \kappa_y^{3/2} \sqrt{\kappa_x})$, respectively.
% Their results can be extended to gradient methods using the fully first-order estimator \citep{kwon2023fully,chen2023near} and following the analysis of Section~\ref{sec:ub}.
However, compared to our lower bounds, the gaps of $\kappa_y^{3/4}$ for the general case and $\kappa_y^{1/4}$ for the lower-level quadratic case still exist even when $\bar \kappa_y =  \kappa_y$.
\end{remark}

% Now we formally prove our main results of the lower bounds for deterministic bilevel problems, including Theorem \ref{thm:NC-SC}, \ref{thm:SC-SC}, and \ref{thm:C-SC}.

% Now, we formally prove Theorem \ref{thm:C-SC} and \ref{thm:SC-SC} 

% All the above NC-SC and (S)C-SC deterministic lower bounds apply to the zero-respecting algorithm class (Definition~\ref{dfn:zr-alg}).
All our lower bounds in Section \ref{subsec:lb-det} apply to deterministic algorithms.
In Section~\ref{subsec:lb-stoc}, we will use a different construction to extend our NC-SC lower bound for $p=1$ to randomized algorithms.
It is open whether our other results (SC-SC/C-SC lower bounds and NC-SC lower bounds for $p \ge 2$) can also be extended to the randomized algorithms, which may relate to some open problems in minimization problems~\citep[Section 6.2]{carmon2021lower}.

%The complete proofs can be found in Appendix \ref{cor:apx-scsc} and \ref{cor:apx-csc}.

% \paragraph{Comparison of additional related works.}
% In the literature, there are also many studies on the lower bounds for minimax optimization, which is the special case of bilevel optimization with $g = -f$. 
% Some of our techniques also appear in these prior works. In minimax optimization, either coupling the lower-level zero-chain \citep{li2021complexity,wang2024efficient} or rescaling the smoothness constant \citep{zhang2021complexity,zhang2022lower} could lead to tight lower bounds. However, in bilevel optimization, we need to apply both techniques in a single instance to obtain a good lower bound.
% This also makes our construction more challenging. Thanks to the application of both of these techniques, we can obtain a tighter lower bound than \citep{ji2023lower}, which only used the rescaling technique from  \citet{zhang2022lower}.
% For readers' convenience of comparison, we also present the constructions in prior works in Appendix~\ref{apx:related-lb}.

\subsection{Lower Bounds for Stochastic Problems} \label{subsec:lb-stoc}

In this subsection, we present a different construction for proving lower bounds in first-order smooth NC-SC problems. The advantage of our second construction is that it further encompasses the randomized and stochastic algorithms defined according to \citep{arjevani2023lower}. Here, ``randomness'' means that the algorithm can query random points during its execution, and ``stochasticity'' means that the algorithm has access to noisy first-order oracles instead of the exact gradients. Below, we formally define the stochastic first-order oracle (SFO) and the algorithm class.

\begin{dfn}[Stochastic first-order oracle] \label{dfn:SFO}
Let $f$ and $g$ be two continuous differentiable functions. We define the stochastic first-order oralce (SFO) for the bilevel problem $(f,g)$ as he mapping $\sO^{\rm sfo}(\vx,\vy) = (\hat \nabla f(\vx,\vy), \hat \nabla f(\vx,\vy))$, where $\hat \nabla f(\vx,\vy)$ and $ \hat \nabla f(\vx,\vy)$ are random vectors that satisfy:
\begin{align*}
    \E \hat \nabla f(\vx,\vy) =& \nabla f(\vx,\vy), \quad \E \Vert \hat \nabla f(\vx,\vy) - \nabla  f(\vx,\vy) \Vert^2 \le \sigma^2; \\
    \E \hat \nabla g(\vx,\vy) =& \nabla g(\vx,\vy), \quad \E \Vert \hat \nabla g(\vx,\vy) - \nabla  g(\vx,\vy) \Vert^2 \le \sigma^2,
\end{align*}
and $\sigma^2>0$ is the variance of SFO.
% We call $\sO_h^{{\rm fo}}: \sR^d \rightarrow \sR$ a $\sigma$-SFO for
% a differentiable function $h: \sR^d \rightarrow \sR $ if 
% $\sO_h^{\rm sfo}(\vw) = \hat \nabla h(\vw)$, where $\hat \nabla h(\vw)$ is a random vector 
% satisfying $\E [\hat \nabla h(\vw)] = \nabla h(\vw)$ and $\E  \Vert \hat \nabla h(\vw) - \nabla h(\vw) \Vert^2 \le \sigma^2$, where $\sigma>0$.
% For a bilevel problem $(f,g)$, we define its SFO as the mapping $\sO^{\rm sfo} = (\sO_f^{\rm sfo}, \sO_g^{\rm sfo})$.
\end{dfn}

\begin{dfn}[Randomized stochastic first-order algorithm] \label{dfn:rand-alg-sfo}
% Under both the deterministic and stochastic settings, 
% Let $\sO^{{\rm fo}} = (\sO_f^{{\rm fo}}, \sO_g^{{\rm fo}})$, where 
% where $\sO_f^{{\rm fo}}$ and $\sO_g^{{\rm fo}}$ are the first-order oracles of $f$ and $g$, respectively. 
A stochastic first-order algorithm $\texttt{A}$ consists of a distribution $\sP_r$ over a measurable set $\gR$ and a sequence of measurable mappings $\{ \texttt{A}^{t}\}_{t \in \sN}$ such that $ \texttt{A}^{t}$ takes in the first $t-1$ oracle responses and the random seed $r \in \gR$ to product the $t$-th query $\vw^t$, such that $\vw^0 = \vzero$ and  
\begin{align}  \label{eq:rand-seq-sfo}
    \vw^t = \texttt{A}^t \left(r, \sO^{\rm sfo}(\vw^0), \cdots, \sO^{\rm sfo}(\vw^{t-1}) \right), \quad \forall t \in \sN_+.
\end{align}                      
We denote $\gA^{{\rm sfo}}$ as the set of all randomized algorithms that follow protocol (\ref{eq:rand-seq-sfo}).
% and $\gA^{\text{sfo-det}} \subseteq \gA^{\text{sfo-rand}}$ as the set of all deterministic algorithms that follow protocol~(\ref{eq:rand-seq-sfo}) with a fixed $r$. 
% A first-order algorithm $\texttt{A}$ consists of a sequence of measurable mappings $\{ \texttt{A}^{t}\}_{t \in \sN}$ such that $ \texttt{A}^{t}$ takes in the first $t-1$ oracle responses to product the $t$-th query $\vz^t = (\vx^t,\vy^t)$, such that $\vz^0 = (\vzero,\vzero)$ and  
% \begin{align}  \label{eq:rand-seq}
%     \vz^t = \texttt{A}^t \left(\sO^{\rm fo}(\vz^0), \cdots, \sO^{\rm fo}(\vz^{t-1}) \right), \quad \forall t \in \sN_+.
% \end{align}                                                                            
% We denote $\gA^{{\rm fo}}$ as the set of all deterministic algorithms that follow protocol (\ref{eq:rand-seq}).
\end{dfn}

% {\color{red} randomized algorithm.}

To derive our lower bounds for stochastic bilevel problems, we 
leverage the following stochastic NC zero-chains \citep{arjevani2023lower}.
The first chain in Definition \ref{dfn:Arj-NC-chain-stoc} can give an $\Omega(\epsilon^{-4})$ lower bound for any algorithm following protocol (\ref{eq:rand-seq-sfo}) with a fixed random seed~$r$, and the second one in Definition \ref{dfn:NC-chain-stoc-random} further amplifies the first one with random rotation and a soft projection to give lower bounds for any $r \sim \sP_r$.

% which is used to prove the $\Omega(\epsilon^{-4})$ lower bound for finding $\epsilon$-stationary points in stochastic NC optimization. 

\begin{dfn}[Stochastic NC zero-chain \citep{arjevani2023lower}] \label{dfn:Arj-NC-chain-stoc}
Let $\bar f^{\text{nc-s}}: \sR^T \rightarrow \sR$ and its stochastic first-order oracle $\sO : \vx \mapsto \hat \nabla \bar f^{\text{nc-s}}(\vx)$ be 
\begin{align*}
    \bar f^{\text{nc-s}}(\vx) &=  \left( -\Psi(1) \Phi(\vx_1) + \sum_{i=2}^T [\Psi(-\vx_{i-1}) \Phi(-\vx_i) - \Psi(\vx_{i-1} ) \Phi(\vx_i)] \right), \\
    \hat \nabla \bar f^{\text{nc-s}}(\vx) &= [\nabla \bar f^\text{nc-s}(\vx)]_j  \left( 1 +
    \vone_{j= {\rm prog}(\vx) +1} \left( \frac{\xi}{p}-1\right) \right), \quad \xi \sim {\rm Bernouli}(p), 
\end{align*}
where the component functions $\Psi, \Phi: \sR \rightarrow \sR$ are 
\begin{align*}
    \Psi(x) =  \vone_{x>1/2} \exp \left( 1 -\frac{1}{(2x-1)^2} \right) \quad {\rm and} \quad \Phi(x) = \sqrt{\rm e} \int_{-\infty}^x \exp(- t^2/2 ) {\rm d}t.
\end{align*}
\end{dfn}

% \begin{lem}[{\citet[Lemma 2 and 3]{arjevani2023lower}}] \label{lem:Arj}
% $\bar f^{\text{nc-s}}$ and $\hat \nabla \bar f^{\text{nc-s}}$ satisfy:
% \begin{enumerate}
%     \item We have $\bar f^{\text{nc-s}}(\vzero)-  \inf_{\vx \in \sR^T} \bar f^{\text{nc-s}}(\vx) \le 12 T $;
%     \item $\bar f^{\text{nc-s}}(\vx)$ has $152$-Lipschitz continuous gradients,
%     \item The SFO $\sO : \vx \mapsto \hat \nabla \bar f^{\text{nc-s}}(\vx)$ is unbiased for $\nabla \bar f^{\text{nc-s}}$, has variance bounded by $ 23^2 (1-p)/p $, and $\bar f^{\text{nc-s}}(\vx)$ equipped with $\sO$ is a probability-$p$ zero chain;
%     \item If $\vert \vx_i \vert < 1 $ for any $ i \le T$, then $ \Vert \nabla \bar f^{\text{nc-s}}(\vx) \Vert > 1$. 
% \end{enumerate}
% % \end{lem}
% \citet{arjevani2023lower} gave lower bounds for randomized algorithms by their above construction . We  also recall this construction and its property as follows. 

\begin{dfn}[Randomized stochastic NC zero-chain \citep{arjevani2023lower}] \label{dfn:NC-chain-stoc-random}
Let the matrix $\mU$ be uniformly sampled from ${\rm Orth}(d,T): = \{ \mU \in \sR^{d \times T} \mid \mU^\top \mU  =\mI_T\}$. And
Let $\bar f^{\text{nc-rs}}: \sR^d \rightarrow \sR$ and its stochastic first-order oracle $\sO : \vx \mapsto \hat \nabla \bar f^{\text{nc-rs}}(\vx)$ be 
\begin{align*}
    \bar f^{\text{nc-rs}}(\vx) &=  \bar f^{\text{nc-s}}(\mU^\top \rho(\vx))) + \frac{1}{10} \Vert \vx \Vert^2,   \\
    \hat \nabla \bar f^{\text{nc-rs}}(\vx) &= \mJ(\vx)^\top \mU \hat \nabla \bar f^{\text{nc-s}}(\mU^\top \rho(\vx)) + \frac{1}{5} \vx. 
\end{align*}
where $R=230\sqrt{T}$, $\rho(\vx) = x \big / \sqrt{1+ \Vert \vx \Vert^2/ R^2}$, and $\mJ(\vx) = \left[ \frac{\partial \rho_i(\vx)}{\partial \vx_j} \right]_{i,j}$ is the Jacobian of~$\rho$.
\end{dfn}

Next, we recall the following properties of the hard instance $\bar f^\text{nc-rs}$.

\begin{lem}[{\citet[Lemma 6 and 7]{arjevani2023lower}}] \label{lem:Arj-random}
$\bar f^{\text{nc-rs}}$ and the associated stochastic first-order oracle $\sO : \vx \mapsto \hat \nabla \bar f^{\text{nc-rs}}(\vx)$ satisfy:
\begin{enumerate}
    \item We have $\bar f^{\text{nc-rs}}(\vzero)-  \inf_{\vx \in \sR^T} \bar f^{\text{nc-rs}}(\vx) \le 12 T $;
    \item $\bar f^{\text{nc-rs}}(\vx)$ has $155$-Lipschitz continuous gradients,
    \item The SFO $\sO$ is unbiased for $\nabla \bar f^{\text{nc-rs}}$, has variance bounded by $ 23^2 (1-p)/p $.
    %and $\bar f^{\text{nc-rs}}(\vx)$ equipped with $\sO$ is a probability-$p$ zero chain;
    \item Let
     $\{\vx^t \}_{t \in \sN}$ be the sequence generated by a randomized algorithm $\texttt{A} \in \gA^{{\rm sfo}}$. For all $t \le (T - \log(2/\delta))/ (2p)$ and any deterministic mapping $\gM: \sR^T \rightarrow \sR^T$, with probability $1-\delta$, we have $\Vert \nabla \bar f^{\text{nc-sc}}(\gM(\vx^t)) \Vert \ge 1/2 $.
\end{enumerate}
\end{lem}

% We then extend the lower bound to the stochastic setting. Here we only consider the case $p=1$ and leave the study of high-order smooth problems for future work.
% In this case, we can prove the lower bounds using a simpler construction. 

{Now, we are ready to present our hard instance for stochastic NC-SC bilevel problems. Compared with our construction in the deterministic case, we abandon the coupling variable~$\vy$ in Eq. (\ref{eq:our-f}) since the sub-chain $h^{\rm sc}$ in Eq. (\ref{eq:sc-sub-chain}) has unbounded gradients, which can cause higher variance as in the stochastic setting \citep{li2021complexity}. By only using the rescaling variable $\vz$, our upper- and lower-level functions $f, g: \sR^T \times \sR^T \rightarrow \sR$ and the associated stochastic first-order oracle $\sO: (\vx,\vz) \mapsto ( \hat \nabla f(\vx,\vz),  \hat \nabla g(\vx,\vz))$ are formally defined as:}
\begin{align} \label{eq:our-fg-stoc}
\begin{split}
      f(\vx,\vz)&: = \frac{L_1 \beta^2}{155} \bar f^{\text{nc-rs}}( \vz / \beta ) \quad {\rm and} \quad 
    g(\vx,\vz): = \frac{\mu_y}{2} \Vert \vz \Vert^2 - L_1 \langle \vx, \vz \rangle, \\
    \hat \nabla f(\vx,\vz) &= \left(\vzero_{T}, \frac{L_1 \beta}{155} \hat \nabla \bar f^{\text{nc-rs}}(\vz / \beta) \right)  \quad {\rm and} \quad
   \hat \nabla g(\vx,\vz)  = \nabla g(\vx,\vz),
\end{split}
\end{align}
where both $\bar f^{\text{nc-rs}}$ and $\hat \nabla \bar f^{\text{nc-rs}}$ are from Definition \ref{dfn:Arj-NC-chain-stoc}. 
Let $\epsilon^{-\gamma}$ be the lower bound for nonconvex stochastic optimization \citep{arjevani2023lower}, where $\gamma = 2$ if $\sigma=  0$ and $\gamma=4$ if $\sigma > 0$. Since our construction above ensures that $\vz^*(\vx) = \kappa_y \vx $ and $f(\vx,\vz^*(\vx)) = (L_1 \beta^2/ 155)  \bar f^{\text{nc-rs}}( \kappa_y \vx / \beta)$, we can follow the analysis of \citep{arjevani2023lower}
to obtain the lower bound of $\Omega( (\kappa_y/ \epsilon)^{\gamma})$ by invoking Lemma \ref{lem:Arj-random} and appropriately choosing the parameters~$\beta$ and $T$.  We state the formal theorem below.

\begin{thm}[Stochastic NC-SC lower bound] \label{thm:NC-SC-stoc}
For any $L_1 \ge \mu_y > 0$, $\sigma>0$, and $\Delta >0$, the bilevel problem instance $(f, g) \in \gF^{\text{nc-scq}}(L_1,\mu_y, \Delta)$ equipped with a stochastic first-order oracle (SFO) with variance $\sigma^2$ such that, to find an $\epsilon$-stationary point of the  hyper-objective $F(\vx) = f(\vx,\vz^*(\vx))$, any randomized stochastic first-order algorithm $\texttt{A} \in \gA^{\rm sfo}$ requires at least
\begin{align*}
    \Omega(\kappa_y^2 L_1 \Delta \epsilon^{-2} + \kappa_y^4 L_1 \sigma^2 \Delta \epsilon^{-4})
\end{align*}
SFO calls, where $\kappa_y = L_1/\mu_y$ is the condition number of the lower-level problem.
\end{thm}

\begin{proof}
See Appendix \ref{apx:proof-thm-NCSC-stoc} for the complete proof.
\end{proof}

As the first-order oracle of our lower-level function is deterministic, we can still apply Lemma~\ref{lem:zr-grad-hvp} to obtain the same lower bounds for stochastic HVP-based methods.
Under the noiseless setting $(\sigma=0)$, the result of Theorem \ref{thm:NC-SC-stoc} 
aligns with that of Theorem~\ref{thm:NC-SC} when $p=1$, but allows for randomization in the execution of algorithms. Under the noisy (stochastic) setting $(\sigma>0)$, Theorem \ref{thm:NC-SC-stoc} indicates that the dependencies on $\kappa_y$ and~$1/\epsilon$ both become larger than the noiseless lower bound. Compared with the lower bound of $ \Omega(\sqrt{\kappa_y} \epsilon^{-2} + \kappa_y^{1/3} \sigma^2 \epsilon^{-4} ) $ in NC-SC minimax problems \citep{li2021complexity}, we again find that the condition number dependency is provably larger in NC-SC bilevel problems. Now, we compare our stochastic NC-SC lower bound with the best-known upper bounds.
% Note that our stochastic lower bound only holds for first-order methods, but not for HVP-based methods since 

\begin{remark}[Known stochastic NC-SC upper bounds] The current best upper bound for stochastic NC-SC bilevel problems is $\tilde \gO(\bar \kappa_y^9 \epsilon^{-4})$, achieved by an HVP method  \citep{ji2021bilevel}. 
If the estimator $\hat \nabla g$ is mean-squared smooth, one can turn their upper bound on HVP calls to SFO calls by approximating the HVP oracle via finite difference of gradients \citep{yang2023achieving}. Without the mean-squared smoothness assumption on~$\hat \nabla g$, the best-known upper bound of first-order methods is $\tilde \gO(\bar \kappa_y^{11} \epsilon^{-6})$ \citep{kwon2024complexity,chen2023near}. Notably, significant gaps exist in the condition number dependency compared to our lower bound even when $\bar \kappa_y = \kappa_y$.
\end{remark}

% Now, we prove Theorem \ref{thm:NC-SC-stoc} in the following.
% % Before giving the proof of , we make a few remarks on the theorem. 
% % First, in the noiseless setting $(\sigma=0)$, the lower bound aligns with 
% % Second, we 
% Note that our construction for stochastic bilevel problems does
% not couple the strongly convex sub-chains $h^{\rm sc}$ as in the deterministic setting (Section \ref{subsec:lb-det}). Interestingly,
% we find that directly coupling the sub-chains without modifications can worsen the lower bound in the noisy setting, because the quadratic function $h^{\rm sc}$ is not Lipschitz continuous and results in a larger variance in its stochastic first-order oracle. The same issue also appears in the lower bound of stochastic NC-SC minimax problems \citep{li2021complexity}. 


% \subsection{Old Result}
% We then turn to the lower bound in the stochastic setting. To use Lemma \ref{lem:pro-zero-chain}, we must design a stochastic first-order oracle (SFO) with bounded variance.  A natural candidate of the SFO is the form in Lemma \ref{lem:bernoulli-sg}, but it requires the objective to have a $\gO(1)$-bounded gradient. In order to satisfy this condition, we follow \citet[Section 5.2]{li2021complexity} to restrict the functions on a bounded domain and scale $h^{\rm sc}$ by a factor of $1/K$.
% Formally, define $\gC_R^d \subset \sR^d$ as the hypercube $\gC_R^d = \{ \vx \in \sR^d \mid \Vert \vx \Vert_{\infty} \le R   \}$. We let $\vy = (\vy^{(1)}, \vy^{(2)})$ with $\vy^{(i)} \in \gC_{K R_y}^{2T+1}$ for $i = 1,2$, $\vz = (\vz^{(1)},\cdots, \vz^{(T)})$ with $\vz^{(i)} \in \gC_{K R_y}^K$ for all $i =1,\cdots,T$, and define the (unscaled) upper- and lower-level functions $\bar f^{\text{sg}}, \bar g^{\rm sg} : \gC_{R_x}^{2T+1} \times \gC_{KR_y}^{4T+2} \times \gC_{K R_y}^{KT} \rightarrow \sR $ as % We let 
% % We consider the following bilevel problem on the hypercube:
% % \begin{align*}
% %     \min_{\vx \in \gC^{2T+1}_{R_x}, \vy \in \gC^{4T+2}_{K R_y}, \vz \in \gC^{K T}_{K R_y}} \bar f^{\rm sg}(\vx, \vy,\vz), \quad {\rm s.t.} \quad (\vy,\vz) = \arg \min_{\vy \in \gC^{2T+1}_{K R_y}, \vz \in \gC^{K T}_{K R_y}} \bar g^{\rm sg}(\vx,\vy,\vz),
% % \end{align*}
% % where $\vy = (\vy^{(1)}, \vy^{(2)})$ with $\vy^{(i)} \in \gC^{2T+1}_{K R_y}$ and $\vz = (\vz^{(1)},\cdots, \vz^{(T)})$ with $\vz^{(i)} \in \gC^K_{K R_y}$. The upper- and lower-level functions are defined as
% \begin{align*}
%       \bar f^{\rm sg}(\vx,\vy,\vz) &:= - \Psi(1) \Phi\left(y_{1}^{(1)}\right) 
%     + \sum_{i=1}^T 
%     H(y_{2i}^{(1)}, y_{2i+1}^{(1)}) + \sum_{i=1}^T h^{\rm reg}(y_{2i}^{(2)}, y_{2i-1}^{(2)}/2, \vz^{(i)}), \\
%      \bar g^{\rm sg}(\vx,\vy,\vz) &:= \sum_{i=1}^2 h^{\rm scale}(\vx,\vy^{(i)}) + \frac{1}{K} \sum_{i=1}^T h^{\rm sc}(x_{2i}, x_{2i-1}/2, \vz^{(i)}),
% \end{align*}
% where $H(x_1,x_2) = \Psi(-x_1) \Phi(-x_2) - \Psi(x_1) \Phi(x_2) $ and the component functions $\Psi, \Phi: \sR \rightarrow \sR$ are defined in Definition \ref{dfn:li-NC-chain}; 
% $h^{\rm sc}: \sR \times \sR \times \sR^K \rightarrow \sR$, $h^{\rm reg}:\sR \times \sR \rightarrow \sR$, and $h^{\rm scale}: \sR^{2T+1} \times \sR^{2T+1} \rightarrow \sR$ remain unchanged as in Eq. (\ref{eq:our-components-h}). Note that we can let $K = \Theta(\kappa_y^{1/3})$ to make the lower-level problem $\mu_y$-strongly convex.

% Next, we define the stochastic first-order oracles. We only add stochastic noise to the gradient of $h^{\rm sc}$ with respect to~$\vz$ while keeping the gradients of other components noiseless. Formally, the oracles are
% \begin{align} \label{eq:our-SFO}
% \begin{split}
%  \hat \nabla \bar f^{\rm sg}(\vx,\vy,\vz) &= \nabla \bar f^{\rm sg}(\vx,\vy), \qquad \hat \nabla_x \bar g^{\rm sg}(\vx,\vy,\vz) = \nabla_x \bar g^{\rm sg}(\vx,\vy,\vz), \\
%     \hat \nabla_y \bar g^{\rm sg}(\vx,\vy,\vz) &= \nabla_y \bar g^{\rm sg}(\vx,\vy,\vz), \quad 
%     \hat \nabla_z \bar g^{\rm sg}(\vx,\vy, \vz^{(i)}) =  \frac{1}{K} \hat \nabla_z h^{\rm sc}(x_{2i}, x_{2i+1}/2, \vz^{(i)}),
% \end{split}
% \end{align} 
% Inspired by Lemma \ref{lem:bernoulli-sg}, we set 
% \begin{align} \label{eq:stoc-sg-h}
%     [\hat \nabla_z h^{\rm sc}(x_1,x_2,\vz^{(i)})]_j =  
%     [\nabla_z h^{\rm sc}(x_1,x_2,\vz^{(i)}) ]_j \left( 1 +
%     \vone_{j= {\rm prog}(\vz^{(i)}) +1} \left( \frac{\xi}{p} -1\right) \right),
% \end{align}
% where $\xi \sim {\rm Bernoulli}(p)$. The properties of our hard instance and the associated SFO are summarized in the following lemma.

% \begin{lem} 
% Let $K \ge 1$ and $R_x, R_y = \gO(1)$.
% The hard instance $(\bar f^{\rm sg},\bar g^{\rm sg})$ equipped with the stochastic first-order oracle $\sO: (\vx,\vy,\vz) \mapsto (\hat \nabla \bar f^{\rm sg}(\vx,\vy,\vz), \hat \nabla \bar g^{\rm sg}(\vx,\vy,\vz))$ in Eq. (\ref{eq:our-SFO}) and (\ref{eq:stoc-sg-h}) satisfy the following.
% \begin{enumerate}
%     \item There exists a numerical constant $\bar \ell_1$ such that $\nabla \bar f^{\rm sg}(\vx,\vy,\vz)$ is $\bar \ell_1$-Lipschitz.
%     \item $\nabla \bar g^{\rm sg}(\vx,\vy, \vz)$ is $5$-Lipschitz continuous.
%     \item $\bar g^{\rm sg}(\vx,\vy,\vz)$ is  $(1/K^3)$-strongly convex jointly in $(\vy,\vz)$.
%      \item  There exists a numerical constant $G$ such that $\sO$ is a stochastic first-order oracle with variance  $\sigma^2 \le G^2 (1-p)/p $.
%     \item Let $\vx^0 = \vy^0 = \vz^0 = \vzero$. For any $ t \le (KT - 1-  \log (1/\delta))/ (2p)$, the iterate $(\vx^t,\vy^t\,\vz^t)$ generated by any zero-respecting algorithm $\texttt{A} \in \gA^{\text{zr}}(\sO)$ satisfies $x^t_{2T} = x^t_{2T+1} = 0$ with probability at least $1-\delta$.
% \end{enumerate}
% \end{lem}

% \begin{proof}
% Item 1-3 can be proved in the same way as Lemma \ref{lem:property}. By Lemma \ref{lem:bernoulli-sg}, we know item~3 holds for $G = \Vert \hat \nabla_z \bar g^{\rm sg}(\vx,\vy,\vz) \Vert_{\infty}$. For every $i = 1,\cdots, T$, we have 
% \begin{align*}
%     \nabla_z \bar g^{\rm sg}(\vx,\vy,\vz^{(i)})  =
%     \frac{1}{K}
%     \nabla_z h^{\rm sc}(x_{2i}, x_{2i+1}/2, \vz^{(i)})  = \frac{1}{K} \left( \left( \frac{1}{K^2} \mI_K + \mA_K  \right)  \vy - \vb_{x_{2i},x_{2i+1}/2} \right).
% \end{align*}
% Then it is straightforward to verify that 
% \begin{align*}
%     \Vert \nabla_z \bar g^{\rm sg}(\vx,\vy,\vz^{(i)}) \Vert_{\infty} \le 4 R_y + R_x = \gO(1). 
% \end{align*}
% Finally, item 5 follows from  Lemma \ref{lem:pro-zero-chain} and Lemma \ref{lem:bernoulli-sg}.
% \end{proof}
% Next, we give the analogy to Lemma \ref{lem:hyperfunc} and Corollary \ref{cor:hyperfunc}.

% \begin{lem}\label{lem:hyperfunc-stoc}
% Let $(\vy^*(\vx), \vz^*(\vx)) = \arg \min_{\vy \in \gC_{KR_y}^{2T+1}, \vz \in \gC_{KR_y}^{K T}} \bar g^{\rm sg}(\vx,\vy,\vz)$. If $K \ge 10$ and $R_y \ge 30 R_1$, then there exists some values $c_1,c_2$ (which may depend on $K$) satisfying $\vert c_1 \vert < 4 \times 10^4$ and $\vert c_2 \vert \le 100$ such that when $\vy = \vy^*(\vx)$ and $\vz = \vz^*(\vx)$ we have that
% \begin{align*}
%     h^{\rm reg}(y_{2i}, y_{2i+1}/2, \vz^{(i)}) = 6 K^3 \left( x_{2i} - \frac{1}{2} x_{2i+1} \right)^2.
% \end{align*}
% \end{lem}

% \begin{cor} \label{cor:hyperfunc-stoc}
% In our 
% construction, the hyper-objective $ \barF^{\rm sg}(\vx) = \bar f^{\rm sg}(\vx, \vy^*(\vx), \vz^*(\vx))$ is 
% \begin{align*}
%     \bar F^{\rm sg}(\vx)  = - \Psi(1) \Phi(\hat x_1)
%     + \sum_{i=1}^T  H(\hat x_{2i}, \hat x_{2i-1}) + 6 \sum_{i=1}^T \left( \hat x_{2i} - \frac{1}{2} \hat x_{2i+1} \right)^2 := \bar f^{\text{nc-2}}(\hat \vx).
% \end{align*}
% where $\hat \vx = K^{3/2} \vx$ and  $\bar f^{\text{nc-2}}:\sR^{2T+1} \rightarrow \sR$ the nonconvex zero-chain in Definition \ref{dfn:li-NC-chain}.
% \end{cor}

% \begin{proof}[Proof of Lemma \ref{lem:hyperfunc-stoc}]
% By Lemma \ref{lem:li-y-star}, the minimizer to $\bar g^{\rm sg}$ remains the same as if there are no constraints. Then we can prove this lemma in the same way as we prove Lemma \ref{lem:hyperfunc}.
% \end{proof}

% Finally, we prove a lower bound of 
% $\Omega( \kappa_y^{7/3} \epsilon^{-4})$ for stochastic first-order bilevel optimization, as stated in the following theorem. Our result significantly improves the prior lower bounds, $\Omega(\epsilon^{-4})$ for NC optimization \citep{arjevani2023lower} and $\Omega(\kappa_y^{1/3} \epsilon^{-4})$ for NC-SC minimax optimization \citep{li2021complexity}, in the $\kappa_y$ dependency.

% \begin{thm}[Stochastic NC-SC lower bound]
% There exists a numerical constants $a_0 \in (0,1)$ such that, 
% for any $L_1 >0$, $\Delta >0$ and $\mu_y \in (0, a_0 L_1]$, there exists a bilevel problem instance $(f, g) \in \gF^{\text{nc-scq}}(L_1,\mu_y, \Delta)$ equipped with a stochastic first-order oracle (SFO) with variance $\sigma^2$ such that, to find an $\epsilon$-stationary point of the  hyper-objective $F(\vx)$, any first-order zero-respecting algorithm requires at least
% \begin{align*}
%     \Omega ( \kappa_y^{7/3} L_1 \sigma^2 \Delta \epsilon^{-4})
% \end{align*}
% SFO calls, 
%     where $\kappa_y = L_1/\mu_y$ is the condition number of the lower-level problem.
% \end{thm}

% \begin{proof}
% Similar to the proof of deterministic lower bound in Theorem \ref{thm:NC-SC}, we rescale the instance as well as their domains and define $f,g : \gC^{2T+1}_{\beta R_x} \times \gC^{4T+2}_{\beta K R_y} \times \gC^{K T}_{\beta K R_y} \rightarrow \sR$ as
% \begin{align*}
%      f(\vx,\vy,\vz) = \frac{L_1 \beta^2}{\bar \ell_1} \bar f^{\rm sg} \left( \frac{\vx}{\beta}, \frac{\vy}{\beta}, \frac{\vz}{\beta} \right), \quad g(\vx,\vy,\vz) = \frac{L_1 \beta^2}{\bar \ell_1} \bar g^{\rm sg} \left( \frac{\vx}{\beta}, \frac{\vy}{\beta}, \frac{\vz}{\beta} \right),
% \end{align*}
% where $R_x= 2$, $R_y = 30R_x$ and $\bar \ell_1 = 80204$. The above rescaling ensures that both $f$ and $g$ are $L_1$-Lipschitz continuous.
% We can choose  $ K = \left \lfloor \left({{L_1}/{(\bar \ell_1 \mu_y)}} \right)^{1/3} \right \rfloor$ 
% to ensure that $g(\vx,\vy,\vz)$ is $\mu_y$-strongly convex jointly in $(\vy,\vz)$. If 
% $\mu_y \lesssim L_1$ then we have $K \ge 10$, which satisfies the condition of Lemma \ref{lem:li}.
% The hyper-objective after scaling is
% \begin{align*}
%     F(\vx) &=  \frac{ L_1}{\bar \ell_1} \left( \beta^2 \Psi(1) \Phi(\hat x_1/\beta)
%     + \sum_{i=1}^T \beta^2 H(\hat x_{2i}/\beta, \hat x_{2i-1}/\beta) + 6 \sum_{i=1}^T \left( \hat x_{2i} - \frac{1}{2} \hat x_{2i+1} \right)^2 \right) \\
%     &= \frac{L_1 \beta^2}{\bar \ell_1} \bar f^{\text{nc-2}}(\hat \vx/ \beta),
% \end{align*}
% where the rescaled variable is $\hat \vx = K^{3/2} \vx$. 
% By Lemma \ref{lem:property-stoc} (item 4), the variance of SFO for $(f,g)$ is upper bounded by $(G L_1 \beta / \bar \ell)^2 (1-p)/p$, which means we can 
% let the variance be bounded by $\sigma^2$ if we 
% choose
% \begin{align} \label{eq:para-sigma}
% p = \min \left\{1, \left( \frac{G L_1 \beta}{\sigma \bar \ell_1} \right)^2  \right\}.
% \end{align}
% By Lemma \ref{lem:property-stoc} (item 5), for any $ t \le ((KT-1-  \log (1/\delta))/ (2p)$ we have $\vx_{2T}^t =  0$. Furthermore, by Lemma \ref{lem:li} and our rescaling, 
% $\vx^t$ is not a $2\epsilon$-stationary point of $F(\vx)$ if we 
% set
% \begin{align} \label{eq:hyper-beta-stoc}
%     \beta = \frac{6 \bar \ell_1 \epsilon}{L_1 K^{3/2}}.
% \end{align}
% Then it implies an $\Omega(K T)$ lower bound for any first-order zero-respecting algorithm to find an $\epsilon$-stationary point of $F(\vx)$. To make the lower bound as large as possible, we choose $T$ to be largest possible value that satisfies the constraint $F(\vzero) - \inf_{\vx \in \sR^{2T+1}} F(\vx) \le \Delta$. By item 1 of Lemma \ref{lem:Car-Psi-Phi}  we know that $\Psi(0) = \Phi(0) =0 $, and item 4 we know that $ 0 < \inf_{x,y \in \sR} \Psi(x) \Phi(y) < \sup_{x,y \in \sR} \Psi(x) \Phi(y) < 12 $. Therefore, we have that
% \begin{align*}
%     F(\vzero) - \inf_{\vx \in \sR^{2T+1}} F(\vx) \le \frac{12 T L_1 \beta^2}{\bar \ell_1} \le \Delta,
% \end{align*}
% where the last inequality holds if we set
% \begin{align} \label{eq:para-T-stoc}
%     T = \left \lfloor 
%      \frac{\bar \ell_1 \Delta}{12L_1 \beta^2} 
%     \right \rfloor.
% \end{align}
% Let $\delta = 1/2$.
% Combining all the above analysis, we know that with probability at least $1/2$, any zero-respecting algorithm can not find a $2\epsilon$-stationary point of $F(\vx)$ in SFO calls less than 
% \begin{align*}
%     \Omega(K T / p) &= \Omega \left( 
%     \left( \frac{\bar \ell_1}{L_1} \right)^3 \frac{K \sigma^2 \Delta}{G^2 \beta^4}
%     \right) = \Omega \left( \frac{L_1 \Delta \sigma^2 K^7}{\bar \ell_1 G^2 \epsilon^4} \right) = \Omega\left( \frac{L_1 \Delta \sigma^2 \kappa_y^{7/3}}{\epsilon^4} \right),
% \end{align*}
% where the last step uses $K = \Theta(\kappa_y^{1/3})$.
% Finally, by Markov inequality, we know that the above result also implies the same lower bound for finding $\epsilon$-stationary point in expectation.
% \end{proof}

% % \subsection{Finite-Sum Setting}

% % To prove the lower bound for finite-sum NC-SC bilevel problems, we separately consider the case $n \ge m$ and $n< m$. For the case $n \ge m $, we let $\{ \mQ_j \}_{j=1}^n \in {\rm Orth}(T,nT,n)$ be the orthogonal matrix sequence defined according to Definition \ref{dfn:orth-matrix} and define the unscaled upper- and lower-level functions $\bar f^{\text{nc-sc-n}} , \bar g^{\text{sc-n}} : \sR^{n T} \rightarrow \sR$ as
% % \begin{align} \label{eq:our-hard-ins-n}
% %     \bar f^{\text{nc-sc-n}}(\vx,\vy) = \frac{1}{n} \sum_{j=1}^n \sqrt{n} \bar f^{\rm nc}_{1,1}(\mQ_j \vy)
% %     \quad \bar g^{\text{sc-n}}(\vx,\vy) = \frac{1}{2 \kappa_y} \Vert \vy \Vert^2 - \langle \vx, \vy \rangle.
% % \end{align}
% % For the case $n <m $, the construction is slightly more complicated.  
% % We
% % let $\vy = (\vy^{(0)}, \vy^{(1)}, \vy^{(2)})$ with $ \vy^{(0)} \in \sR^m$, $ \vy^{(1)}, \vy^{(2)} \in \sR^{m(2T+1)}$, $\vz = (\vz^{(1)},\cdots, \vz^{(T)})$ with $\vz^{(i)} \in \sR^{nK}$ for all $i = 1,\cdots,T$, and define the (unscaled) upper- and lower-level functions $\bar f^{\text{nc-sc-m}}, \bar g^{\text{sc-m}} : \sR^{m(2T+1)} \times \sR^{m(4T+3)} \times \sR^{mKT} \rightarrow \sR $ as
% % \begin{align} \label{eq:our-fg-fs}
% % \begin{split}
% %      \bar f^{\text{nc-sc-m}}(\vx,\vy,\vz)&:= \frac{1}{m} \sum_{j=1}^m \bar f^{\text{nc-sc}}(\mP_j \vx), \\
% % \bar g^{\text{sc-m}}(\vx,\vy,\vz) &:= \sum_{i=1}^2 h^{\rm scale}(\vx,\vy^{(i)}) +  \sum_{i=1}^T h^{\text{sc-m}}(x_{2i}, x_{2i+1}, \vz^{(i)}),
% % \end{split}
% % \end{align}
% % where $\Upsilon_r: \sR \rightarrow \sR$ is from Definition \ref{dfn:nc-zo-chain}; $h^{\rm sc}: \sR \times \sR \times \sR^K \rightarrow \sR$ is from Definition \ref{dfn:li-SC-subchain}, 
% % $h^{\rm chain}: \sR^{2T+1} \rightarrow \sR$, $h^{\rm reg}:\sR \times \sR \rightarrow \sR$, and $h^{\rm scale}: \sR^{2T+1} \times \sR^{2T+1} \rightarrow \sR$ are
% % \begin{align} \label{eq:our-components-h}
% % \begin{split}
% %       h^{\rm chain}(\vy) &:=\frac{\sqrt{\nu}}{2} (y_1 - 1)^2 + \frac{1}{2} \sum_{i=1}^T  (y_{2i-1}- y_{2i})^2; \\
% %        h^{\rm scale}(\vx,\vy) &:= \frac{1}{2 K^{3/2}} \Vert \vy \Vert^2 -  \langle \vx, \vy \rangle; \\
% %      h^{\rm reg}(y_1,y_2, \vz)&: = c_1 y_1^2 + c_1 y_2^2 + c_2 \Vert \vz \Vert^2
% % \end{split}
% % \end{align}
% % with $c_1,c_2\in \sR$ being the parameters to be chosen later. 
% % We let $K = \Omega(\sqrt{\kappa_y})$ to ensure the lower-level function is $\mu_y$-strongly convex.

% % We modify our hard instance as follows.
% % \begin{align*}
% %     h^{\text{sc-m}}(x_{2i}, x_{2i+1}, \vz^{(i)}) = \frac{1}{n} \sum_{j=1}^{n} \frac{1}{2} \left(\vz^{(j)}\right)^\top \left( \frac{1}{(nK)^2} \mI_{nK} + \mQ_j^\top \mA_K \mQ_j \right) \vz^{(j)} - \left \langle \vb_j, \mQ_j \vz^{(j)} \right \rangle,
% % \end{align*}
% % where $\vb_j = (\mQ_j \vx)_{2i} 
% % \ve_1 - (\mQ_j x)_{2i+1} \ve_K$.

% % We let $s_i$ be the last index $j$ such that the corresponding sub-chain is fully activated. Then we define the upper-level function as
% % \begin{align*}
% %    h^{\rm chain}(\vy) = \frac{\sqrt{\nu}}{2} \left( 
% %    \vy_{s_1} - \vy_{s_0} 
% %    \right) ^2 + \frac{1}{2} \sum_{i=1}^T \left( \vy_{2s_i-1} - \vy_{2s_i} \right)^2.
% % \end{align*}

% % \begin{thm}[NC-SC lower bound, finite-sum setting] \label{thm:NC-SC-fs}
% % There exists a numerical constants $a_0 \in (0,1)$ such that, 
% % for any $L_1 >0$, $\Delta >0$ and $\mu_y \in (0, a_0 L_1]$, there exists a bilevel problem instance $(f,g) \in \gF^{\text{nc-scq-as}}(L_1,\mu_y, \Delta)$ with the finite-sum structure $f = \sum_{j=1}^n f_j / n$,  and an incremental first-order oracle (IFO) such that, to find an $\epsilon$-stationary point of the  hyper-objective $F(\vx)$, any first-order zero-respecting algorithm requires at least
% % \begin{align*}
% %     \Omega ( n+ m + \sqrt{n+m} \kappa_y^2  L_1  \Delta \epsilon^{-2})
% % \end{align*}
% % IFO calls, where $\kappa_y = L_1/\mu_y$ is the condition number of the lower-level problem.
% % \end{thm}

% % \paragraph{Instance I.}
% % An $\Omega(m)$ lower bound is oblivious. We also present the instance for completeness. xxx

% % We let

% % Then $\{ f_j\}_{j=1}^m$ is $\gO(1)$-averaged smooth, and the resulting hyper-objective is
% % \begin{align*}
% %     F(\vx) = \frac{1}{m} \sum_{j=1}^m \sqrt{m} f^{\rm nc}(\mQ_j \bar \vx),
% % \end{align*}
% % where $\bar \vx = \kappa_y \vx$. Then we can following the same analysis as Fang et al. to prove an $\Omega(\sqrt{m} \kappa^2 \epsilon^{-2})$ lower bound.

% % \paragraph{Instance II.}


% % Setting $K = \Theta(\sqrt{\kappa_y/n})$, we finally obtain an $\Omega( \sqrt{n} \kappa^2 \epsilon^{-2} )$ lower bound.

% % Note that we can also get an $\Omega(n)$ lower bound by using the same resisting strategy to define $s_0$.  We let
% % \begin{align*}
% %     h^{\rm start}(\vz) = \frac{1}{n}\sum_{j=1}^n \frac{1}{2} \Vert \vz \Vert^2 -  \langle \ve_j, \vz \rangle.
% % \end{align*}
% % Then it requires at least $\Omega(n)$ IFO calls to activate the coordinate $\vy_{s_0}$.

% % \paragraph{The final lower bound.}
% % Combining them together, the final lower bound is $\Omega \left( m+n + \sqrt{m+n} \kappa^2 \epsilon^{-2} \right)$, which significantly improves the previous result of 
% % $\Omega(m+ \sqrt{m} \epsilon^{-2}) $ lower bound that immediately comes from single-level lower bound.

\section{Upper Bounds} \label{sec:ub}

% {\color{red} need to rewrite this section.}

We observe a significant gap between 
the current best upper bounds and the lower bounds we established in the previous section in the condition number dependency even for deterministic smooth NC-SC bilevel problems: Specifically, our lower bound is $\Omega(\kappa_y^2 \epsilon^{-2})$ , while the state-of-the-art upper bound \citep{ji2021bilevel,chen2023near} is $\tilde \gO(\bar \kappa_y^4 \epsilon^{-2})$. In this section,
we show that the upper bound achieved by fully first-order method in \citep{chen2023near} can be improved to $\tilde \gO(\bar \kappa_y^{7/2} \epsilon^{-2})$ 
by solving the lower-level problem with AGD \citep{nesterov1983method}. 
% Although it is possible to extend the HVP methods in \citep{ji2021bilevel} to achieve the same rates, we do not go in that direction.
Furthermore, we also obtain a $\tilde \gO(\bar \kappa_y^{13/4} \epsilon^{-7/4})$ upper bound for second-order smooth NC-SC problems, using fully first-order oracles to achieve the same complexity given by the accelerated HVP-based method in \citep{yang2023accelerating}.

% For second-order smooth NC-SC problems, our lower bounds is $\Omega(\kappa_y^{25/14} \epsilon^{-12/7})$, while the upper bound by \citet{chen2023near} is $\tilde \gO(\bar \kappa_y^{15/4} \epsilon^{-7/4})$.
% %For stochastic smooth NC-SC problems, our lower bound is $\Omega(\kappa_y^4 \epsilon^{-4})$, while the upper bound by \citet{chen2023near} is $\tilde \gO(\kappa_y^{12} \epsilon^{-6})$.
% In this section, we provide improved upper bounds to reduce the gap in this deterministic case, and the study of stochastic cases is left to future work.

\subsection{Algorithm Overview}

The main challenge of designing fully first-order method for bilevel optimization is that the algorithm only has access to first-order oracles of $f$ and $g$, while the calculation of hyper-gradient $\nabla F$ in Eq. (\ref{eq:hyper-grad}) requires second-order information of $g$. 
To address this issue, we follow \citet{kwon2023fully} to use the following fully first-order estimator:
% The main idea of fully first-order bilevel approximation \citep{kwon2023fully,chen2023near} is to approximate the hyper-gradient in Eq. (\ref{eq:hyper-grad}), which involves second-order information according to the following formula that only involves first-order information:
\begin{align} \label{eq:f2ba-grad}
    \nabla F_{\lambda} (\vx) = \nabla_x f(\vx, \vy_{\lambda}^*(\vx)) + \lambda ( \nabla_x g(\vx, \vy_{\lambda}^*(\vx)) - \nabla_x g(\vx,\vy^*(\vx)),
\end{align}
where $ \vy_{\lambda}^*(\vx) =  \arg \min_{\vy \in \sR^{d_y}} \{f(\vx,\vy) + \lambda g(\vx,\vy')\}$. \citet{kwon2023fully} showed that $\nabla F_{\lambda}$ is a good approximation to $\nabla F$ for a large penalty $\lambda$, as we recall below.
\begin{lem}[{\citet[Lemma 3.1]{kwon2023fully}}] Let $\lambda \ge 2 L_1/ \mu_y$. For first-order smooth NC-SC bilevel problem $(f,g) \in$ $\gF^{\text{nc-sc}}(L_0,L_1,L_2,\mu_y,\Delta)$, 
$\nabla F_{\lambda}$ is well-defined and satisfies
\begin{align*}
    \Vert \nabla F_{\lambda}(\vx) - \nabla F(\vx) \Vert \le 8 \bar \kappa_y^3 L_1 / \lambda.
\end{align*}
\end{lem}
Based on the above lemma, \citet{kwon2023fully} proposed a fully first-order method that can find an $\epsilon$-stationary point in $\tilde \gO({\rm poly}(\bar \kappa_y) \epsilon^{-3})$ first-order oracle calls. 
Later on, \citet{chen2023near} proposed the Fully First-order Bilevel Approximation (F${}^2$BA) method that improved the complexity to 
$\tilde \gO(\bar \kappa_y^4 \epsilon^{-2})$. The F${}^2$BA method consists of two loops:
(1) The outer loop performs gradient descent (GD) on $\vx$ using the hyper-gradient estimator~$\nabla F_\lambda$, under the setting $\lambda = \Omega(\epsilon^{-1})$ that ensures  $\Vert \nabla F_\lambda(\vx) - \nabla F(\vx) \Vert \le \epsilon$. (2) The inner loop computes the lower-level solutions $\vy_{\lambda}^*(\vx)$ and $\vy^*(\vx)$ by GD and uses them to construct the estimator $\nabla F_\lambda(\vx)$. 

These loops introduce two sources of condition number dependency in the complexity: (1) The outer-loop complexity depends on $L_F$, the smoothness of the hyper-objective $F(\vx)$, which further depends on~$\bar \kappa_y$. (2) Both the computation of $\vy_{\lambda}^*(\vx)$ and $\vy^*(\vx)$ requires solving a strongly-convex sub-problem, whose complexity depends on $\kappa_y$. However, the lower-level GD sub-solver used in \citep{chen2023near} has a suboptimal $\tilde \gO(\kappa_y)$ complexity.

In this work, we obtain better upper bounds by replacing GD with AGD \citep{nesterov1983method}, which only requires  $\tilde \gO(\sqrt{\kappa_y})$ gradient steps to solve the strongly-convex lower-level problems. This naturally results in an improvement of a $\sqrt{\kappa_y}$ factor in the total complexity. The procedure of AGD is presented in Algorithm~\ref{alg:AGD}, with its guarantee stated in Lemma~\ref{lem:AGD}.

% Recently, \citet{chen2023near} improved the outer loop complexity from $\gO(\lambda L_F \epsilon^{-2})$ in the original analysis by \citet{kwon2023fully} to $\gO( \log(\lambda) L_F \epsilon^{-2})$, leading to a near-optimal dependency on $\epsilon$.  
% both  use naive gradient descent (GD) in the lower level that 


\begin{algorithm*}[htbp]  
\caption{AGD $(h,K,  \vz_0)$} \label{alg:AGD}
\begin{algorithmic}[1] 
\STATE $ \tilde \vz_0 = \vz_0$ \\[1mm]
\STATE \textbf{for} $ k =0,1,\cdots,K-1 $ \\[1mm]
\STATE \quad $\vz_{k+1} = \tilde \vz_k - \frac{1}{L} \nabla h(\tilde \vz_k)$ \\[1mm]
\STATE \quad $\tilde \vz_{k+1} = \vz_{k+1} + \frac{\sqrt{\kappa_h}-1}{\sqrt{\kappa_h} +1} (\vz_{k+1} - \vz_k) $ \\[1mm]
% \STATE \quad \textbf{Option II: } $x_{t+1}  = x_t - \eta \hat \nabla \fL_{\lambda}^*(x_t) / \Vert \hat \nabla \fL_{\lambda}^*(x_t) \Vert $ \quad // Normalized Gradient Descent 
\STATE \textbf{end for} \\[1mm]
\STATE \textbf{return} $\vz_K$ 
\end{algorithmic}
\end{algorithm*}
% of. However, GD is suboptimal compared to the accelerated gradient descent (AGD) \citep{nesterov1983method} which can achieve the optimal $\gO(\sqrt{\kappa_y})$ complexity. 
% The upper bound in \citep{chen2023near} has achieved near-optimal dependency on $\epsilon$. However, the lower-level solver they used is suboptimal for strongly convex optimization. 
% We can get an improved upper bound if we replace it by the , which only requires  gradient steps to solve the lower-level problem. This results in an improvement of a  factor in the total complexity, as stated in the following theorem.
% The above lemma motivates us to set  and use the fully first-order hyper-gradient estimator $\nabla F_{\lambda}(\vx)$ to perform gradient descent on variable $\vx$ to find an $\epsilon$-stationary point of $F(\vx)$. 
% To obtain the estimator $\nabla F_{\lambda}(\vx)$, the algorithm requires computing the 
% optimal lower-level solutions $\vy_{\lambda}^*(\vx)$ and $\vy^*(\vx)$. 
% These two steps form the basis of fully first-order bilevel approximation \citep{kwon2023fully} and they also correspond to the following two sources of the condition number dependency in the complexity upper bounds:
% \begin{enumerate}
%     \item The complexity of gradient descent on $F(\vx)$ depends on its smoothness $L_F$, which further depends on the global condition number $\bar \kappa_y$ \citep{ghadimi2018approximation}.
%     \item The expression of estimator $\nabla F_{\lambda}(\vx)$ in Eq. (\ref{eq:f2ba-grad}) involves the optimal solution to the lower-level function $\vy^*(\vx)$. This means that each time an iterative algorithm uses $\nabla F_\lambda(\vx)$, it needs to solve a strongly convex lower-level problem that requires at least $\Omega(\sqrt{\kappa_y})$ complexity in the worst case \citep[Section 2.1.4]{nesterov2018lectures}.
% \end{enumerate}
% In the follow-up work of, \citet{chen2023near} analyze the source 1 complexity in detail and 

% Obtaining $\mF(\vx)$ requires accessing the , which are approximated via vanilla gradient descent (GD) in prior works \citep{ji2021bilevel,kwon2023fully,chen2023near}. The upper-level requires $\gO(L_{F} \epsilon^{-2})$ gradient steps, while the lower-level requires $\tilde \gO(\kappa_y )$ gradient steps, resulting in the total first-order complexity of $\tilde \gO(\bar \kappa_y^4 \epsilon^{-2})$ \citep{chen2023near}.

 
\begin{lem}[{\citet[Theorem 2.2.4]{nesterov2018lectures}}] \label{lem:AGD}
Running Algorithm \ref{alg:AGD} on an $L_h$-gradient Lipschitz and $\mu_h$-strongly convex objective $h: \sR^d \rightarrow \sR$ outputs $\vz_K$ satisfying
\begin{align*}
    \Vert \vz_K - \vz^* \Vert^2 \le (1 + \kappa_h) \left( 1 - \frac{1}{\sqrt{\kappa_h}} \right)^K \Vert \vz_0 - \vz^* \Vert^2,
\end{align*}
where $\vz^* = \arg \min_{\vz \in \sR^d} h(\vz)$ and $\kappa_h = L_h/ \mu_h$ is the condition number of $h$.
\end{lem}

\subsection{Upper Bounds for First-order Smooth Problems} \label{subsec:F2BA}


% \begin{algorithm*}[t]  
% \caption{MS-AC-SA $(h,K,\gV_0, \vz_0)$} \label{alg:AC-SA}
% \begin{algorithmic}[1] 
% \STATE $ \tilde \vz_0 = \vz_0$ \\[1mm]
% \STATE \textbf{for} $ k =0,1,\cdots,K-1 $ \\[1mm]
% \STATE \quad $N_k = \left \lceil \max \left \{ 4 \sqrt{\frac{2L_h}{}} \right\} \right \rceil $
% \STATE \quad $\vz_{k+1} = \tilde \vz_k - \frac{1}{L} \nabla h(\tilde \vz_k)$ \\[1mm]
% \STATE \quad $\tilde \vz_{k+1} = \vz_{k+1} + \frac{\sqrt{\kappa_h}-1}{\sqrt{\kappa_h} +1} (\vz_{k+1} - \vz_k) $ \\[1mm]
% % \STATE \quad \textbf{Option II: } $x_{t+1}  = x_t - \eta \hat \nabla \fL_{\lambda}^*(x_t) / \Vert \hat \nabla \fL_{\lambda}^*(x_t) \Vert $ \quad // Normalized Gradient Descent 
% \STATE \textbf{end for} \\[1mm]
% \end{algorithmic}
% \end{algorithm*}

% For proving upper bounds, we require the following additional assumption, which always holds on a compact set if $\nabla_y f(\vx,\vy)$ is $L_1$-Lipschitz continuous.
% \begin{asm} \label{asm:Lip-y}
% There exists $L_0 >0$ such that $\Vert \nabla_y f(\vx,\vy) \Vert \le L_0 $ for every $\vx \in \sR^{d_x}$.
% \end{asm}

% In this subsection, we present upper complexity bounds for NC-SC first-order smooth problems. 
We present the F${}^2$BA${}^+$ method in Algorithm \ref{alg:F2BA} for NC-SC bilevel problems, which uses AGD to solve the lower-level problem in the F${}^2$BA method \citep{chen2023near}.
The analysis mainly follows that of \citep{chen2023near} and relies on the following two lemmas. The first lemma gives the smoothness of $F(\vx)$.


%$\kappa_y$ dependency of bilevel problems, which comes from the following two sources: 
% First, the smoothness of  hyper-objective $F(\vx) = f(\vx,\vy^*(\vx))$ depends on the smoothness of $\vy^*(\vx)$, which further depends on $\kappa_y$ in general \citep[Lemma 2.2]{ghadimi2018approximation}.
% Second, 
\begin{lem} \label{lem:smooth-phi}
For first-order smooth NC-SC bilevel problem $(f,g) \in$ $\gF^{\text{nc-sc}}(L_0,L_1,L_2,\mu_y,\Delta)$, the hyper-objective $F(\vx) = f(\vx,\vy^*(\vx))$ satisfies
%[] 
% Consider NC-SC bilevel function class in Definition \ref{dfn:NC-SC-func}. For every $\vx, \vx' \in \sR^{d_x}$, we have 
\begin{align} \label{eq:hyper-smooth}
    \Vert \nabla F(\vx) - \nabla F(\vx') \Vert \le L_F \Vert \vx-  \vx' \Vert, \quad \forall \vx,\vx' \in \sR^{d_x},
\end{align}
where $L_F = \gO( \bar \kappa_y^3 L_1 )$ and $\bar \kappa_y = \max\{L_0,L_1,L_2 \} / \mu_y$. In particular, if the lower-level function is quadratic and $(f,g) \in \gF^{\text{nc-scq}}(L_1,\mu_y, \Delta)$, we have $L_F = \gO(\kappa_y^2 L_1)$, where $\kappa_y = L_1/ \mu_y$.
\end{lem}

\begin{proof}
In the general case, $L_F = \gO( \bar \kappa_y^3 L_1)$ can be verified by Eq. (\ref{eq:hyper-grad}). See, for example, {\citep[Lemma 2.2]{ghadimi2018approximation}}. For quadratic lower-level problems defined in Eq.~(\ref{eq:quadratic-g}), we have $F(\vx) = \nabla_x f(\vx,\vy^*(\vx)) - \mJ \mH^{-1} \nabla_y f(\vx,\vy^*(\vx))$. Then, from the facts $\Vert \mJ \Vert \le L_1$, $\Vert \mH^{-1} \Vert \le \1/\mu_y$, and $\Vert \nabla \vy^*(\vx) \Vert \le \kappa_y$, we can easily prove that $L_F = \gO(\kappa_y^2 L_1)$.
\end{proof}

The second lemma from \citep{chen2023near}, states that if the lower-level solver achieves sufficient accuracy, then the outer loop of Algorithm \ref{alg:F2BA} requires $\gO(L_F \Delta \epsilon^{-2})$ iterations to find an $\epsilon$-stationary point.

\begin{algorithm*}[t]  
\caption{F${}^2$BA${}^+$ $(f,g, \eta_x, T, K,\vx_0,\vz_0)$} \label{alg:F2BA}
\begin{algorithmic}[1] 
\STATE $ \vz_0 = \vy_0$ \\[1mm]
\STATE \textbf{for} $ t =0,1,\cdots,T-1 $ \\[1mm]
\STATE \quad $\vz_{t+1} = \text{AGD}( g(\vx_t,\,\cdot\,), K, \vz_t)$ \\[1mm]
\STATE \quad $\vy_{t+1} = \text{AGD} (f(\vx_t,\,\cdot\,) + \lambda g(\vx_t, \,\cdot\,) , K, \vy_t    )$ \\[1mm]
% \STATE \quad $ \vy_t^0 =  \vy_{t}, ~ \vz_t^0 = \vz_{t}$ \\[1mm]
% \STATE \quad \textbf{for} $ k =0,1,\cdots,K-1$ \\[1mm]
% \STATE \quad \quad $ \vz_t^{k+1} = \vz_t^{k}- \frac{1}{L_1} \nabla_y g(\vx_t, \vz_t^k)   $ \\[1mm]
% \STATE \quad \quad $ \vy_t^{k+1} = \vy_t^k - \frac{1}{2 \lambda L_1} \left( \nabla_y f(\vx_t,\vy_t^k) +  \lambda  \nabla_y g(\vx_t,\vy_t^k) \right)$ \\[1mm]
% \STATE \quad \quad $\tilde \vz^{k+1}_t = \vz^{k+1}_t + \frac{\sqrt{\kappa_y}-1}{\sqrt{\kappa_y} +1} (\vz^{k+1}_t - \vz^k_t) $ \\[1mm]
% \STATE \quad \quad $\tilde \vy^{k+1}_t = \vy^{k+1}_t + \frac{2\sqrt{\kappa_y}-1}{2\sqrt{\kappa_y} +1} (\vy^{k+1}_t - \vy^k_t) $  \\[1mm]
% \STATE \quad \textbf{end for} \\[1mm]
\STATE \quad $ \hat \nabla F_\lambda(\vx_t)= \nabla_x f(\vx_t,\vy_{t+1}) + \lambda ( \nabla_x g(\vx_t,\vy_{t+1}) - \nabla_x g(\vx_t,\vz_{t+1}) )$ \\[1mm]
\STATE \quad  $\vx_{t+1} = \vx_t -  \eta_x  \hat \nabla F_{\lambda}(\vx_t)$ \\[1mm]
% \STATE \quad \textbf{Option II: } $x_{t+1}  = x_t - \eta \hat \nabla \fL_{\lambda}^*(x_t) / \Vert \hat \nabla \fL_{\lambda}^*(x_t) \Vert $ \quad // Normalized Gradient Descent 
\STATE \textbf{end for} \\[1mm]
\STATE \textbf{return} $\vx_{\rm out} = \frac{1}{T} \sum_{t=0}^{T-1} \vx_{t}$ \\[1mm]
\end{algorithmic}
\end{algorithm*}

\begin{lem}[{\citet[Theorem 4.1]{chen2023near}}]\label{lem:F2BA}
For first-order smooth NC-SC bilevel problem $(f,g) \in$ $\gF^{\text{nc-sc}}(L_0,L_1,L_2,\mu_y,\Delta)$, 
if in each iteration $t$ of Algorithm \ref{alg:F2BA} the lower-level solver outputs 
$\vz_t^K$ and $\vy_t^K$ such that 
\begin{align*}
    \Vert \vz_{t+1} - \vy^*(\vx_t) \Vert^2+ \Vert \vy_{t+1} - \vy_{\lambda}^*(\vx_t) \Vert^2 \le \gamma \Vert \vz_t - \vy^*(\vx_t) \Vert^2+ \Vert \vy_t - \vy_{\lambda}^*(\vx_t) \Vert^2
\end{align*}
with $\gamma \lesssim \min\{ 1/2, \epsilon^2/ (\bar \kappa_y^3 L_1^3) \} / \lambda^2$, then setting $\lambda \asymp L_1 \bar \kappa_y^3 \epsilon^{-1}$ and $\eta_x \asymp 1/L_F$ can ensure the algorithm to find an $\epsilon$-stationary point of $F(\vx)$ in $T = \gO( L_{F} \Delta \epsilon^{-2})$ outer-loop iterations.
\end{lem}

% The third lemma recalls the theoretical guarantee of AGD, which implies that the inner loop of our algorithm only requires $\tilde \gO(\sqrt{\kappa_y})$ iterations.

% gives the inner complexity of our algorithm.
% The convergence rate depends on the smoothness constant of $F(\vx)$ given by Lemma \ref{lem:smooth-phi}. 
% Define the lower-level distances. 
% We observe a gap between 
% the existing upper bounds \citep{chen2023near,ji2023lower} and our lower bounds. 
% The second contribution of our paper is to provide improved upper bounds to reduce the gap. Below, we first review state-of-the-art methods \citep{chen2023near} and show that using accelerated gradient descent \citep{nesterov1983method} to solve the lower-level leads to an improvement for solving NC-SC bilevel problems (Definition \ref{dfn:NC-SC-func}). As a warm up, we start with the case $p=1$, and the improvement in the case $p=2$ can be obtained for the same reason.
% \begin{cor} \label{cor:smooth-phi}
% \end{cor}
% The upper bound in \citep{chen2023near} has achieved near-optimal dependency on $\epsilon$. However, the lower-level solver they used is suboptimal for strongly convex optimization. 
% We can get an improved upper bound if we replace it by the accelerated gradient descent (AGD) \citep{nesterov1983method}, which only requires $\tilde \gO(\sqrt{\kappa_y})$ gradient steps to solve the lower-level problem. 

Combining Lemma \ref{lem:AGD}, \ref{lem:smooth-phi}, and \ref{lem:F2BA}, we can easily prove that the total complexity of F${}^2$BA${}^+$ is $\tilde \gO(\sqrt{\kappa_y} L_F \Delta \epsilon^{-2})$,
where $L_F = \gO(\bar \kappa_y^3 L_1)$ for general lower-level functions and $L_F = \gO(\kappa_y^2 L_1) $ for quadratic lower-level functions. We summarize the complexity of F${}^2$BA${}^+$ in the following, which improves the $\tilde \gO(\kappa_y L_F \Delta \epsilon^{-2})$
complexity of F${}^2$BA \citep{chen2023near}
 by a $\sqrt{\kappa_y}$ factor due to the use of lower-level AGD. 

%Finally, we obtain the following result by plugging the value of $L_F$ by Lemma \ref{lem:smooth-phi}.

% This results in an improvement of a $\sqrt{\kappa_y}$ factor in the total complexity, as stated in the following theorem.

\begin{restatable}{thm}{thmubp}  \label{thm:ub-p1}
Under the setting of Lemma \ref{lem:F2BA}, Algorithm \ref{alg:F2BA} can find an
$\epsilon$-stationary point of $F(\vx)$ in 
$\tilde \gO ( \bar \kappa_y^{7/2} L_1 \Delta \epsilon^{-2}  )$
first-order oracle complexity. In particular, if the lower-level function  $g(\vx,\vy)$ is quadratic as in Eq.  (\ref{eq:quadratic-g}), the complexity can be refined to $\tilde \gO ( \kappa_y^{5/2} L_1 \Delta \epsilon^{-2}  )$.
\end{restatable}

\begin{proof}
Recall that $\vy_\lambda^*(\vx) = \arg \min_{\vy \in \sR^{d_y}} f(\vx,\vy) + \lambda g(\vx,\vy)$. For any $\lambda \ge 2 L_1/ \mu_y$, we know that sub-problem for variable $\vy$ is $2\lambda L_1$-smooth and $\lambda \mu_y/2$-strongly convex. Therefore, we can apply Lemma \ref{lem:AGD} with condition number $4 L_1 / \mu_y$, in conjunction with \ref{lem:F2BA}, to prove the upper complexity bound of $\tilde \gO(\kappa_y L_F \Delta \epsilon^{-2})$, where $L_F$ is given by Lemma \ref{lem:smooth-phi}.
\end{proof}

Compared with the $\Omega(\kappa_y^2 \epsilon^{-2})$ lower bound in 
Theorem \ref{thm:NC-SC} when $p=1$, our upper bounds in Theorem \ref{thm:ub-p1} still have a gap of 
$\kappa_y^{3/2}$ even when $\bar \kappa_y = \kappa_y$ for general lower-level problems, and a gap of $\sqrt{\kappa_y}$ for quadratic lower-level problems.

\subsection{Upper Bounds for Second-order Smooth Problems}

Now, we consider second-order smooth NC-SC bilevel problems, where the hyper-objective $F(\vx)$ is proved to have Lipschitz continuous Hessians.

% The following lemma gives the Hessian smoothness constants for the problem.

\begin{lem}[{\citet[Lemma 2.4]{huang2025efficiently}}] \label{lem:smooth-phi-second}
For second-order smooth NC-SC bilevel problem $(f,g) \in \gF^{\text{nc-sc}}(L_0,\cdots,L_3,\mu_y,\Delta)$, the hyper-objective $F(\vx) = f(\vx,\vy^*(\vx))$ satisfies
%[] 
% Consider NC-SC bilevel function class in Definition \ref{dfn:NC-SC-func}. For every $\vx, \vx' \in \sR^{d_x}$, we have 
\begin{align} \label{eq:hyper-smooth-second}
    \Vert \nabla^2 F(\vx) - \nabla F(\vx') \Vert \le \rho_F \Vert \vx-  \vx' \Vert, \quad \forall \vx,\vx' \in \sR^{d_x},
\end{align}
where $\rho_F = \gO( \bar \kappa_y^5 L_2 )$ and $\bar \kappa_y = \max\{L_0,L_1,L_2,L_3 \} / \mu_y$. In particular, if the lower-level function is quadratic and $(f,g) \in \gF^{\text{nc-scq}}(L_1,L_2,\mu_y, \Delta)$, we have $\rho_F = \gO(\bar \kappa_y^3 L_2)$.
\end{lem}

Following \citep{nesterov2006cubic}, we target to find stronger solutions known as second-order stationary points for Hessian smooth nonconvex functions. Compared with first-order stationary points in Definition \ref{dfn:fo-sta}, the second-order stationary points in the following Definition \ref{dfn:so-sta} exclude saddle points and can be equivalent to global minimizers in certain applications \citep{ge2015escaping}.

\begin{dfn} \label{dfn:so-sta}
We call 
a point $\vx \in \sR^d$ an $\epsilon$-second-order stationary point of $F$ if 
\begin{align*}
    \Vert \nabla F(\vx) \Vert \le \epsilon \quad \text{and} \quad \nabla^2 F(\vx) \succeq -\sqrt{\rho_{F} \epsilon} \mI_{d}
\end{align*}
\end{dfn}

For second-order smooth problems, \citet{chen2023near} proposed AccF${}^2$BA to achieve a $\tilde \gO( \kappa_y \sqrt{L_F} \rho_F^{1/4} \epsilon^{-7/4} )$ upper bound by applying nonconvex AGD \citep{carmon2017convex,jin2018accelerated,li2023restarted} in the upper level.
But again, their lower-level GD sub-solver is suboptimal. Based on this observation, we propose 
AccF${}^2$BA${}^+$ in Algorithm \ref{alg:AccF2BA}, which replaces the GD sub-solver in AccF${}^2$BA with AGD.
By keeping the outer loop unchanged, we can directly apply the results from \citep{chen2023near} to obtain the following guarantee of the outer loop, provided that the lower-level sub-solver achieves sufficient accuracy.

% Since we keep the outer loop unchanged, we can directly invoke the following result of \citep{chen2023near} to obtain the guarantee for the outer loop 
% once the lower-level sub-solver can achieve sufficient accuracy.
% \citet{chen2023near} also provided upper bounds for the case $p=2$, where the hyper-objective has Lipschitz continuous Hessians and the outer-level complexity can be further improved \citep{carmon2017convex,li2023restarted}. But again, they used a suboptimal lower-level solver. Therefore, we can also obtain an improved upper bound in this setting by using AGD in the lower level. 


\begin{algorithm*}[t]  
\caption{AccF${}^2$BA${}^+$} \label{alg:AccF2BA}
\begin{algorithmic}[1] 
\STATE $ \vz_0 = \vy_0$, $\vx_{-1} = \vx_0$ \\[1mm]
\STATE \textbf{while} $ t < T $ \\[1mm]
\STATE \quad $\tilde \vx_{t} = \vx_t + (1-\theta) (\vx_t -\vx_{t-1})$ \\[1mm]
\STATE \quad $\vz_{t+1} = \text{AGD}( g(\tilde \vx_t,\,\cdot\,), K_t, \vz_t)$ \\[1mm]
\STATE \quad $\vy_{t+1} = \text{AGD} (f(\tilde \vx_t,\,\cdot\,) + \lambda g(\vx_t, \,\cdot\,) , K_t, \vy_t    )$ \\[1mm]
% \STATE \quad $ \vy_t^0 =  \vy_{t}, ~ \vz_t^0 = \vz_{t}$ \\[1mm]
% \STATE \quad \textbf{for} $ k =0,1,\cdots,K_t-1$ \\[1mm]
% % \STATE \quad \quad $ \vz_t^{k+1} = \vz_t^{k}- \eta_z \lambda \nabla_y g(\vx_{t+1/2}, \vz_t^k)   $ \\[1mm]
% % \STATE \quad \quad $ \vy_t^{k+1} = \vy_t^k - \eta_y \left( \nabla_y f(\vx_{t+1/2},\vy_t^k) +  \lambda  \nabla_y g(\vx_{t+1/2},\vy_t^k) \right)$ \\[1mm]
% % \STATE \quad \textbf{end for} \\[1mm]
% \STATE \quad \quad $ \vz_t^{k+1} = \vz_t^{k}- \frac{1}{L_1} \nabla_y g(\vx_t, \vz_t^k)   $ \\[1mm]
% \STATE \quad \quad $ \vy_t^{k+1} = \vy_t^k - \frac{1}{2 \lambda L_1} \left( \nabla_y f(\vx_t,\vy_t^k) +  \lambda  \nabla_y g(\vx_t,\vy_t^k) \right)$ \\[1mm]
% \STATE \quad \quad $\tilde \vz^{k+1}_t = \vz^{k+1}_t + \frac{\sqrt{\kappa_y}-1}{\sqrt{\kappa_y} +1} (\vz^{k+1}_t - \vz^k_t) $ \\[1mm]
% \STATE \quad \quad $\tilde \vy^{k+1}_t = \vy^{k+1}_t + \frac{2\sqrt{\kappa_y}-1}{2\sqrt{\kappa_y} +1} (\vy^{k+1}_t - \vy^k_t) $  \\[1mm]
% \STATE \quad $\vz_{t+1} = \vz_t^{K_t},~ \vy_{t+1} = \vy_t^{K_t} $ \\[1mm]
\STATE \quad $ \hat \nabla F_\lambda(\tilde \vx_{t})= \nabla_x f(\tilde \vx_{t},\vy_{t+1}) + \lambda ( \nabla_x g(\tilde \vx_{t},\vy_{t+1}) - \nabla_x g(\tilde \vx_{t},\vz_{t+1}) )$ \\[1mm]
% \STATE \quad $x_{t+1/2}^- = x_{t+1/2}$ \\[1mm]
\STATE\quad  $\vx_{t+1} = \tilde \vx_{t} - \eta_x \hat \nabla F_\lambda(\tilde \vx_{t}) $ \\[1mm]
\STATE \quad $ t = t+1 $ \\ [1mm]
\STATE \quad \textbf{if} $ t \sum_{j=0}^{t-1} \Vert \vx_{t+1} - \vx_t \Vert^2 \ge B^2$ \\[1mm]
\STATE \quad \quad $t=0$, $\vx_{-1} = \vx_0 = \vx_t + \xi_t \vone_{ \Vert \hat \nabla F_\lambda(\tilde \vx_{t}) \Vert \le \frac{B}{2 \eta_x} }$, {\rm where} $\xi_t \sim \sB(r)$  \\[1mm] 
\STATE \quad \textbf{end if} \\[1mm]
% \STATE \quad \textbf{if} $ \Vert \hat \nabla \fL_{\lambda}^*(x_t) \Vert \le \frac{4}{5} \epsilon$ \textbf{and} no perturbation added in the last $\fT  $ steps \\[1mm]
% \STATE \quad \quad $x_{t} = x_t - \eta_x \xi_t $, where $\xi_t \sim \sB(r)$  \\[1mm]
% \STATE \quad \textbf{end if} \\[1mm]
% \STATE \quad $x_{t+1} = x_t -  \eta_x  \hat \nabla \fL_{\lambda}^*(x_t)$ \\[1mm]
\STATE \textbf{end while} \\[1mm]
\STATE $T_0 = \arg \min_{\lfloor \frac{T}{2} \rfloor \le t \le T-1} \Vert \vx_{t+1} - \vx_t \Vert$ \\[1mm]
\STATE \textbf{return} $\vx_{\rm out} = \frac{1}{T_0+1} \sum_{t=0}^{T_0} \tilde \vx_{t}$ \\[1mm]
\end{algorithmic}
\end{algorithm*}

\begin{lem}{\citep[Theorem 5.2]{chen2023near}} \label{lem:AccF2BA}
Let $\sqrt{\epsilon \rho_F} \le L_F$.
For second-order smooth NC-SC bilevel problem $(f,g) \in \gF^{\text{nc-sc}}(L_0,L_1,L_2,L_3,\mu_y,\Delta)$, 
if in each iteration $t$ of Algorithm~\ref{alg:AccF2BA} the lower-level solver outputs 
$\vz_{t+1}$ and $\vy_{t+1}$ such that 
\begin{align*}
    \Vert \vz_{t+1} - \vy^*(\vx_t) \Vert^2+ \Vert \vy_{t+1} - \vy_{\lambda}^*(\vx_t) \Vert^2 \le \gamma \Vert \vz_t - \vy^*(\vx_t) \Vert^2+ \Vert \vy_t - \vy_{\lambda}^*(\vx_t) \Vert^2
\end{align*}
with $\gamma \lesssim {\rm poly}(\epsilon) / {\rm poly}(\log d_x, \bar \kappa_y, L_1) $, then setting the hyper-parameters as
\begin{align*}
    \lambda \asymp \max\{ L_1 \bar \kappa_y^3 /\epsilon, ~\bar \kappa_y^{7/2} \sqrt{L_2/\epsilon} \}, ~~ \eta_x \asymp 1/L_F, ~~ \chi \asymp \log(\nicefrac{d_x}{\delta \epsilon}),  \\
    \theta \asymp \sqrt{\eta_x} (L_\rho \epsilon)^{1/4}, ~~ T \asymp \chi / \theta,  ~~ B \asymp \sqrt{\epsilon/\rho_F} / \chi^2, ~~ {\text{and}} ~~ r\asymp \epsilon 
\end{align*}
ensures the algorithm to output an $\epsilon$-second-order stationary point with probability $1-\delta$
in no more than  $\gO( \sqrt{L_{F}} \rho_{F}^{1/4} \Delta \epsilon^{-7/4} \log^6(\nicefrac{d_x}{\delta \epsilon}))$ steps of update in the upper-level variable $\vx$.
\end{lem}

Now, by combining Lemma \ref{lem:AGD}, \ref{lem:smooth-phi-second}, and \ref{lem:AccF2BA}, we naturally arrive at the following theorem for the total complexity of AccF${}^2$BA${}^+$. 

\begin{restatable}{thm}{thmubpacc}  \label{thm:ub-p2}
Under the same setting of Lemma \ref{lem:AccF2BA}, Algorithm \ref{alg:AccF2BA} can find an
$\epsilon$-stationary point of $F(\vx)$ in 
$\tilde \gO ( \bar \kappa_y^{13/4} \sqrt{L_1} L_2^{1/4} \Delta \epsilon^{-7/4}  )$
first-order oracle complexity. In particular, if the lower-level function $g(\vx,\vy)$ is quadratic as in Eq.  (\ref{eq:quadratic-g}), the complexity can be refined to $\tilde \gO ( \bar \kappa_y^{9/4} \sqrt{L_1} L_2^{1/4} \Delta \epsilon^{-7/4} )$.
\end{restatable}

\begin{proof}
By combining Lemma \ref{lem:AGD} and \ref{lem:AccF2BA}, we know that the total complexity of Algorithm \ref{alg:AccF2BA} is no more than $\gO( \sqrt{\kappa_y L_F} \rho_F^{1/4} \Delta \epsilon^{-7/4} \log^6 (\nicefrac{d_x}{\delta \epsilon})  )$. Then we can obtain the claimed complexity by substituting the value of $L_F = \gO(\bar \kappa_y^3 L_1)$, $\rho_F = \gO(\bar \kappa_y^5 L_2)$ for general lower-level problems and $L_F = \gO( \kappa_y^2 L_1 )$, $\rho_F = \gO(\bar \kappa_y^3 L_2)$ according to Lemma \ref{lem:smooth-phi} and \ref{lem:smooth-phi-second}.
\end{proof}

Our upper bounds in Theorem \ref{thm:ub-p2} still have gaps 
compared to the lower bound of $\Omega(\kappa_y^{25/14} \epsilon^{-12/7})$ in Theorem \ref{thm:NC-SC} when $p=2$. Even assuming $\bar \kappa_y = \kappa_y$, there remains a gap of $\kappa_y^{41/28} \epsilon^{-1/28} \approx \kappa_y^{1.464} \epsilon^{-0.036}$ for general lower-level problems, and a gap of $\kappa_y^{13/28} \epsilon^{-1/28} \approx \kappa_y^{0.464} \epsilon^{-0.036}$ for quadratic lower-level problems. 

%It is important for future work to close the gap.

% \begin{remark}[Gap between our lower bound]
% The lower bound in  is , which still has a gap of 
% \end{remark}

% \begin{cor} \label{cor:ub-p1}
% When the quadratic function, the complexity in Theorem \ref{thm:ub-p1} can be refined to 
% \end{cor}

%Closing the gap is definitely an important future work.
 
% \begin{remark}[Gap between our lower bound]
% % The lower bound in Theorem \ref{thm:NC-SC} if $p=2$ is $. 

% % which still has a gap of $\kappa_y^{}$ even when $\bar \kappa_y = \kappa_y$ for general lower-level problems, and a gap of $\sqrt{\kappa_y}$ for quadratic lower-level problems.    
% \end{remark}

% \begin{cor} \label{cor:ub-p2}
% When $g(\vx,\vy)$ is the quadratic function (\ref{eq:quadratic-g}), the complexity in Theorem \ref{thm:ub-p2} can be refined to $\tilde \gO \left( \bar \kappa_y^{2.25} L_1^{0.5} L_2^{0.25} \Delta \epsilon^{-1.75}  \right)$.
% \end{cor}


%
% \begin{proof}
% By Lemma \ref{lem:AGD}, the lower-level problem requires $\tilde \gO(\sqrt{\kappa_y})$ first-order oracle calls every iteration of $\vx$, where $\kappa_y = L_1/ \mu_y$. And the complexity of outer loop directly follows Lemma \ref{lem:AccF2BA}, where we have $L_{\varphi} = \bar \kappa_y^3 L_1 $  by Lemma \ref{lem:smooth-phi} and $\rho_\varphi = \bar \kappa_y^5 L_2$ by \citep[Lemma 5.1]{chen2023near}. Therefore, the total complexity of the algorithm is 
% \begin{align*}
%     \tilde \gO( \kappa_y^{0.5} L_{\varphi}^{0.5} \rho_{\varphi}^{0.25} \Delta \epsilon^{-1.75} ) = \tilde \gO(\bar \kappa_y^{3.25} L_1^{0.5} L_2^{0.25} \Delta \epsilon^{-1.75}).
% \end{align*}
% For the case when the lower-level problem is the quadratic function (\ref{eq:quadratic-g}), also from \citep[Lemma 5.1]{chen2023near} we have the refined smoothness constants $L_\varphi = \bar \kappa_y^2 L_1 $ and $\rho_{\varphi} = \bar \kappa_y^3 L_2$, leading to the refined complexity of $\tilde \gO(\bar \kappa_y^{2.25} L_1^{0.5} L_2^{0.25} \Delta \epsilon^{-1.75})$.
% \end{proof}

\begin{figure}[htbp]
    \centering
    \includegraphics[scale=0.3]{fig/l2reg_det.pdf} 
\caption{Performances of different algorithms when learning the optimal regularization.}
\label{fig:l2reg}
\end{figure}

\subsection{Numerical Experiments}

Our improved upper bounds are simply obtained by plugging the guarantee of AGD into existing analysis of fully first-order methods \citep{chen2023near}. % It is also possible that similar improvements can be obtained by replacing GD with AGD in state-of-the-art HVP-based methods \citep{ji2021bilevel,arbel2022amortized}.
In this section, we conduct experiments to verify the benefit of using AGD to solve the lower-level problems.

% However, we are a bit surprised that most works in the bilevel literature \citep{ghadimi2018approximation,ji2021bilevel,shen2023penalty,chen2023near} only use GD instead of AGD to solve the lower-level problem even in the deterministic case. 
% We also conduct experiments to show the effectiveness of using AGD.



% \newpage

Following \citet{grazzi2020iteration,ji2021bilevel,liu2022bome,chen2023near}, we consider the ``learn-to-regularize'' problem on the ``20 Newsgroup'' dataset, which solves the following bilevel optimization problem:
\begin{align*}
    \min_{\vx \in \sR^{p}} \frac{1}{m} \sum_{i=1}^m \ell_i^{\rm val}(\mY), \quad {\rm s.t.} \quad \mY \in \arg \min_{\mY \in \sR^{p \times q}} \frac{1}{n} \sum_{i=1}^n \ell_i^{\rm tr}(\mY) + \Vert \mW_x \mY  \Vert^2_F ,
\end{align*}
where $\mW_\vx = {\rm diag} \left(\sqrt{\exp(\vx)}\right)$, 
$\vy \in \sR^{p \times q}$ is the parameter matrix of a multi-class logistic model for $p$ features and $q$ classes, $\{ \ell_i^{\rm val}\}_{i=1}^m$ and $\{ \ell_i^{\rm tr}\}_{i=1}^n$ are the validation and training loss on every sample each sample with $m,n$ being the size of the datasets.
In our experiments, we have $p= 130, 107$, $q = 20$, and $m = n = 5,657$. We compare six algorithms, including a baseline without learning the regularization ``w/o Reg'', the HVP-based method AID \citep{ji2021bilevel}, two gradient-based methods F${}^2$BA and AccF${}^2$BA \citep{chen2023near}, and our improved versions with lower-level AGD denoted by F${}^2$BA${}^+$ and AccF${}^2$BA${}^+$.
The experiment result is presented in Figure \ref{fig:l2reg}. From the figure, we can observe a significant gain by solving the lower-level problem with AGD, confirming our theoretical claims.

\section{Conclusion and Future Work}

In this work, we provide lower bounds and accelerated fully first-order algorithms in NC-SC bilevel optimization. For deterministic first-order smooth problems, we prove an $\Omega(\kappa_y^2 \epsilon^{-2})$ lower bound and a $\tilde \gO(\bar \kappa_y^{7/2} \epsilon^{-2})$ upper bound. For deterministic second-order smooth problems, we prove an $\Omega(\kappa_y^{25/14} \epsilon^{-12/7})$ lower bound and a $\tilde \gO(\bar \kappa_y^{13/4} \epsilon^{-7/4})$ upper bound. We also extend our lower bound for many other setups, including deterministic arbitrarily smooth NC-SC problems, deterministic smooth (S)C-SC problems, and stochastic smooth NC-SC problems, demonstrating the gaps between existing upper bounds. 

Definitely, closing the gaps between our lower and upper bounds is an important future work and may require new techniques. It would also be interesting to study other setups, such as 
mean-squared smooth problems \citep{khanduri2021near,yang2021provably, dagreou2022framework,yang2023achieving,chu2024spaba}, 
structured NC-NC problems \citep{kwon2024penalty,chen2024finding,jiang2025beyond}, and  zeroth-order bilevel optimization \citep{aghasi2025fully,aghasi2025optimal} in the future.


%Additionally, our algorithms consist of two loops, and it would also be important to consider single-loop methods for practical use \citep{chen2021closing,chen2022single,dagreou2022framework,hong2023two,chen2024optimal,shen2025single}. 

\section*{Acknowledgments}

We greatly thank Jeongyeol Kwon for his helpful comments on the preliminary version of this manuscript.


\bibliographystyle{plainnat}
\bibliography{sample}

\appendix
\newpage 
% \section{Missing Proofs in Section \ref{sec:lb}}

% % \subsection{Proof of Lemma \ref{lem:property}}

% % \lempro*

% \subsection{Proof of Lemma \ref{lem:hyperfunc}} \label{apx:proof-hyperfunc}

% \lemhyperfunc*

% \begin{proof}

% % Note that $g(\vx,\vy,\vz)$ is separable in each block variable $\vy^{(i)}$ and $\vz^{(i)}$. Therefore, we can calculate the optimal lower-level solution $(\vy^*(\vx),\vz^*(\vx))$ by setting
% % \begin{align*}
% %     \vy^{(i)}&= \arg \min_{\vy \in \sR^{2T+1}} h^{\rm scale}(\vx,\vy) , \quad i = 1,2; \\
% %     \vz^{(i)} &= , \quad i = 1,\cdots,K.
% % \end{align*}
% From 
% The solution to the above equations $\vz = K^2 \vx$ and
% \begin{align} \label{eq:close-yz-star}
% \begin{split}
%     \vy^{(i)} &= \arg \min_{\vz \in \sR^K} h^{\rm sc}(x_{2i},x_{2i+1}, \vz) =  \left( \frac{1}{K^2} \mI_K + \mA_K \right)^{-1} \vb^{(i)}, \quad \forall i = 1,\cdots,K.
% \end{split}
% \end{align}
% where $\vb^{(i)} = x_{2i} \ve_1 - x_{2i+1} \ve_K$. For $\vy^{(i)}$ satisfying Eq. (\ref{eq:close-yz-star}), we have
% \begin{align*}
%     \Vert \vz^{(i)} \Vert^2 =  \Vert \mB \vb^{(i)} \Vert^2,
% \end{align*}
% where $\mB = (\alpha \mI_K + \mA_K)^{-1}$ and $\alpha = 1/K^2$. Let $\mD = \mB^2$, then
% \begin{align*}
%     \Vert \vz^{(i)} \Vert^2  = \left \langle \vb^{(i)}, \mD \vb^{(i)} \right \rangle = D_{1,1} x_{2i}^2 - 2 D_{1,K} x_{2i} x_{2i+1} + D_{K,K} x_{2i+1}^2.
% \end{align*}
% Therefore, plugging in $\vy^*(\vx)$ and $\vz^*(\vx)$ given by Eq. (\ref{eq:close-yz-star}) to $h^{\rm reg}$ in Eq. (\ref{eq:our-f}) yields 
% \begin{align*}
%        % \varphi(\vx) &= \frac{\sqrt{\nu}}{2} ( \bar x_1 - 1)^2 +  \frac{1}{2}\sum_{i=1}^T  (\bar x_{2i-1} - \bar x_{2i})^2 + \nu \sum_{i=1}^{2T} \Upsilon_r (\bar x_i) \\
% h^{\rm reg}(y_{2i},y_{2i+1}, \vz^{(i)}) = (c_1 K^3 + c_2 D_{1,1} x_{2i}^2) - 2 c_2 D_{1,K} x_{2i} x_{2i+1} + (c_1 K^3 + c_2 D_{K,K}) x_{2i+1}^2,
% \end{align*}
% From Lemma \ref{lem:li} we know 
% \begin{align*}
%     D_{K,K} = D_{1,1} = \sum_{i=1}^K B_{1,i} B_{i,1} = \sum_{i=1}^K B_{i,1}^2 \in [0.01 K^3, 400 K^3] 
% \end{align*}
% and
% \begin{align*}
%     D_{1,K} = D_{K,1} = \sum_{i=1}^K B_{K,i} B_{i,1} =  \sum_{i=1}^K B_{1,K+-i} B_{i,1} \in [0.01K^3, 400K^3],
% \end{align*}
% where the fact $B_{K,i} = B_{1,K+1-i}$ can be verified by $\mB = \mM^\top / \det ( \alpha \mI_K + \mA_K )$ with $\mM$ is the cofactor matrix of $  \alpha \mI_K + \mA_K $. Now, Lemma \ref{lem:hyperfunc} is proved by letting 
% \begin{align} \label{eq:setting-c}
%     c_2 = \frac{K^3}{D_{1,K}}, \quad c_1 = 1 - \frac{D_{1,1}}{D_{1,K}}.
% \end{align}
% Finally,
% from the range of $D_{1,1}$ and $D_{1,K}$ it is easy to check that
% \begin{align*}
%     |c_2| \le 100, \quad \vert c_1 \vert < 4 \times 10^4,
% \end{align*}
% which means both $c_1$ and $c_2$ are (almost) numerical constants. This completes the proof.
% \end{proof}


\section{Examples with Lower-Level Quadratics} \label{apx:example-Q}

This section provides a collection of (S)C-SC and NC-SC bilevel optimization applications in which the lower-level function is quadratic. 

\subsection{Convex-Strongly-Convex Examples} \label{apx:exmple-C-SC}
The (S)C-SC setting is often assumed without any examples in the literature \citep{ghadimi2018approximation,hong2023two}. To the best of our knowledge, the only example is a fair resource allocation problem~\citep{srikant2014communication} given by \citet{ji2023lower}. In the following, we provide several machine learning problems that provably satisfy the assumptions, which verify the rationality of assuming the (strong) convexity of the hyper-objective in previous work \citep{ghadimi2018approximation,hong2023two}.

Let $(\mA,\vb)$ denote a dataset with $n$ samples, where each row of $\mA \in \sR^{n \times d}$ represents a $d$-dimensional feature of a sample, and the corresponding element in $\vb \in \sR^n$ represents the corresponding label. We consider the following examples.

\begin{exmp}
Given $T$ different tasks, with each task $i$ corresponds to a training set $(\mA_i^{\rm tr}, \vb_i^{\rm tr})$ and  a testing set $(\mA_i^{\rm test}, \vb_i^{\rm  test})$. Define the covariance matrices of each task as $\mG_i^{\rm test}  =\left(\mA_i^{\rm test}  \right)^\top \mA_i^{\rm test} $ and $\mG_i^{\rm tr}  =\left(\mA_i^{\rm tr}  \right)^\top \mA_i^{\rm tr}$. 
We consider the meta-learning for linear regression~\citep{rajeswaran2019meta}:
    \begin{align*}
        &\quad \min_{\vx \in \sR^d} \frac{1}{T} \sum_{i=1}^T \frac{1}{2} \Vert \mA_i^{\rm tr} \vy_i^*(\vx) - \vb_i^{\rm tr} \Vert^2, \\
        &{\rm s.t.} \quad \vy_i^*(\vx) = \arg \min_{\vy' \in \sR^d} \frac{1}{2} \Vert \mA_i^{\rm test} \vy' - \vb_i^{\rm test} \Vert^2 + \frac{\lambda}{2} \Vert \vy' - \vx \Vert^2.
    \end{align*}
In this example, the hyper-objective $\varphi(\vx) = f(\vx,\vy^*(\vx))$ is (strongly) convex.
\end{exmp}

\begin{proof}
The lower-level solution has a closed form
\begin{align*}
    &\quad \vy_i^*(\vx) = (\mG_i^{\rm test} + \lambda \mI_d)^{-1} \left(\lambda \vx + \left(\mA_i^{\rm test} \right)^\top \vb_i^{\rm test} \right).
\end{align*}
Plugging the above equation into the upper-level function gives the closed form of $\varphi(\vx)$, which is a convex quadratic function because since its Hessian matrix is positive (semi)definite:
\begin{align*}
    &\quad \mH = \frac{1}{T} \sum_{i=1}^T \lambda^2 (\mG_i^{\rm test} + \lambda \mI_d)^{-1} \mG_i^{\rm tr} (\mG_i^{\rm test} + \lambda \mI)^{-1} \succeq \mO_d.
\end{align*}.
\end{proof}

\begin{exmp}
Given a dataset $(\mA,\vb)$ with $n$ samples, we  
consider the Stackelberg game \citep{bruckner2011stackelberg} for robust linear regression. In this model, the learner's goal is to learn to good linear predictor $\vx$ with small error, but the label $\vy$ is modified by an adversarial data provider, whose goal is to make the loss (per sample weighted by $w_i^2$) high, with a cost to modify the original label $\vb$. 
Let $\mD = {\rm diag}(w_1,\cdots, w_n)$.
It can be formulated as the following bilevel problem
\begin{align*}
    &\quad \min_{\vx \in \sR^d} \frac{1}{2} \Vert \mA \vx - \vy^*(\vx) \Vert^2 \\
    &{\rm s.t.} \quad \vy^*(\vx) = \arg \min_{\vy' \in \sR^n} - \frac{1}{2} \Vert \mD (\mA \vx - \vy') \Vert^2 + \frac{\lambda}{2} \Vert \mD(\vy'  -\vb) \Vert^2.
\end{align*}
In this example, the hyper-objective $\varphi(\vx) = f(\vx,\vy^*(\vx))$ is (strongly) convex if $\lambda >1$.
\end{exmp}

\begin{proof}
When $\lambda >1$, the lower-level solution has a closed form
\begin{align*}
    \vy^*(\vx) = \frac{\mD  (\lambda \vb - \mA  \vx)}{\lambda -1}.
\end{align*}
Plugging the above equation into the upper-level function gives the closed form of $\varphi(\vx)$, which is a convex quadratic function because since its Hessian matrix is positive (semi)definite:
\begin{align*}
    \mH = \frac{1}{(\lambda-1)^2}  \mA^\top ( (\lambda -1)^2 \mI_n + \mD_n^2 ) \mA \succeq \mO_d.
\end{align*}
\end{proof}

\begin{exmp}
    \citet{yang2021graph} reformulates graph convolution network (GCN) training~\citep{kipf2016semi} into a bilevel optimization problem, where the lower-level problem is to find the optimal node embeddings that minimize a graph energy function, and the upper-level problem is the loss of the embedding. For the node regression problem, we have the dataset $(\mA,\vb)$ with $n$ nodes. The bilevel problem is:
    \begin{align*}
        &\quad \min_{\vx \in \sR^d} \frac{1}{2} \Vert \vy^*(\vx) - \vb \Vert^2 \\
        &{\rm s.t.} \quad \vy^*(\vx) = \arg \min_{\vy \in \sR^{n}} \frac{1}{2} \Vert \mA \vx - \vy \Vert^2 + \frac{\lambda}{2} \vy^\top \mL \vy,
    \end{align*}
where $ \mL$ is the graph Laplacian matrix. In this example, the hyper-objective $\varphi(\vx) = f(\vx,\vy^*(\vx))$ is (strongly) convex if $\lambda>0$.
% When $A$ is non-singular, by the inequality $\sigma_i(AB) \ge \sigma_i(A) \sigma_n(B)$, we know that $H \succ O$, which implies that $\varphi(\phi)$ is strongly convex.
\end{exmp}

\begin{proof}
When $\lambda >0$, the lower-level solution has a closed form 
\begin{align*}
    \vy^*(\vx) = (\mI_n + \lambda \mL)^{-1} \mA \vx.
\end{align*}
Plugging the above equation into the upper-level function gives the closed form of $\varphi(\vx)$, which is a convex quadratic function because since its Hessian matrix is positive (semi)definite:
\begin{align*}
    \mH = \mA^\top (\mI_n + \lambda \mL)^{-1} (\mI_n + \lambda \mL)^{-1} \mA \succeq \mO_d.
\end{align*}
\end{proof}

\subsection{Nonconvex-Strongly-Convex Examples} \label{apx:exmple-NC-SC}

The NC-SC bilevel problems have many applications in the actor-critic method \citep{konda1999actor}, which is a class of important algorithms in reinforcement learning.  For this algorithm class, there is an actor in the upper level that optimizes a (typically nonconvex) cost function to select a good policy, and a critic in the lower level that evaluates by estimating the Q-function. The Q-function can be calculated based on the Bellman equation, which corresponds to solving a strongly convex function in the basic linear Markov decision process (MDP) setup.

Formally, we define the MDP as a tuple with five elements $(S,A, \gamma, P, R)$, where $S$ is the state space, $A$ is the action space, $\gamma \in [0,1)$ is the discount factor of the rewards, $P$ defines the transition kernel such that $P(s' \mid s,a)$ denotes the probability of transitioning to state $s'$ after taking action $a$
in state $s$, and $R$ defines the reward function such that $R(s,a) $ denotes the immediate reward received after taking action $a$ in state $s$. 
Let $Q^\pi$ denote the Q-function for policy $\pi$. 
Under the linear setup, $Q^\pi$ is assumed to have a linear representation $Q^\pi(s,a) = \langle \theta^\pi, \phi(s,a) \rangle$, where $ \theta^\pi \in \sR^d$ is the weight vector of parameters and $\phi(s,a) \in \sR^d$ is the 
feature vector for the state-action pair $(s,a)$. Then we can show that the induced bilevel optimization problem is nonconvex-strongly-convex with a quadratic lower-level problem.
%For a policy $\pi$ such that $\pi(a \mid s)$ denotes the probability of choosing action $a$ in state $s$, we further define $P^\pi$ as the transition kernel if the agent choose actions following policy $\pi$, \textit{i.e.}, $P^\pi(s' \mid s) = \sum_{a \in A} \pi(a \mid s) P(s' \mid s, a)$.

\begin{exmp}[Actor-critic for linear MDP]
Let the initial state $s_0$ be drawn from a fixed distribution $\rho_0$, and let $\Pi \subset \sR^{\vert S \vert \times \vert A \vert}$ be the policy space. We consider the following problem:
\begin{align*}
    &\quad \max_{\pi \in \Pi} \sum_{a \in A} Q^{\pi}(s_0,a) \pi(a \mid s_0) \rho_0(s_0). \\
    &{\rm s.t.} \quad Q^\pi(s,a) = R(s,a) + \gamma \sum_{s' \in S}  \sum_{a \in A} Q^\pi(s',a) P(s' \mid s ,a) \pi(a \mid s).
\end{align*}
Using linear function approximation to $Q^\pi$, the lower-level function is a linear system in $\theta^\pi$, which can also be reformulated as a equivalent (strongly) convex quadratic optimization.
\end{exmp}

There are many similar examples in reinforcement learning, 
where the lower-level problem is exactly a linear system using linear function approximation.
Interested readers can refer to references \citep{hong2023two,zeng2024two,zeng2024fast}.

\section{Proof of Lemma \ref{lem:hyperfunc}} \label{apx:proof-lem-hyperfunc}

To prove the result, we first recall the following fact about the matrix $\mA_K$.
\begin{lem}[{\citet[Lemma 7]{li2021complexity}}] \label{lem:li}
Let $\mB = ( \mI_K / K^2 + \mA_K )^{-1}$. If $K \ge 10$, we have for all $i = 1,\cdots,K$ that
\begin{align*}
     0.1K \le B_{1,i} \le 20K.
\end{align*}
\end{lem}

Now, we give the proof of Lemma \ref{lem:hyperfunc}.

\begin{proof}[Proof of Lemma \ref{lem:hyperfunc}]
%[Proof of Lemma \ref{lem:hyperfunc}]
% Note that $g(\vx,\vy,\vz)$ is separable in each block variable $\vy^{(i)}$ and $\vz^{(i)}$. Therefore, we can calculate the optimal lower-level solution $(\vy^*(\vx),\vz^*(\vx))$ by setting
% \begin{align*}
%     \vy^{(i)}&= \arg \min_{\vy \in \sR^{2T+1}} h^{\rm scale}(\vx,\vy) , \quad i = 1,2; \\
%     \vz^{(i)} &= , \quad i = 1,\cdots,K.
% \end{align*}
Because the lower-level function $\bar g^{\rm sc}$ is separable in $\vy$ and $\vz$, it is easy to see that $\vz^*(\vx) = K^2 \vx$, and $\vy^*(\vx) = (\vy^{(1)}_*(\vx),\cdots, \vy^{(T-1)}_*(\vx))$ satisfies 
\begin{align*} 
    \vy^{(i)}_*(\vx) =
    %&=  \left( \frac{1}{K^2} \mI_K + \mA_K \right)^{-1} (x_{i} \ve_1 - x_{i+1} \ve_K) := 
    \mB (x_{i} \ve_1 - x_{i+1} \ve_K), \quad \forall i = 1,\cdots,T-1.
\end{align*}
Note that $B_{1,K} = B_{K,1}$ and $B_{1,1} = B_{K,K}$. We have
\begin{align} \label{eq:close-Hx}
    H(\vx) = \sum_{i=1}^{T-1} \frac{a_K}{2} \cdot K^4 (x_i^2 +x_{i+1}^2) + \frac{b_K}{2} \cdot K^2 ( B_{1,1} x_i^2 - 2 B_{1,n} x_i x_{i+1} + B_{1,1} x_{i+1}^2).
\end{align}
According to the above equation, we can let $H(\vx)$ be equivalent to Eq. (\ref{eq:Hx-goal}) by setting
\begin{align} \label{eq:setting-alpha-beta}
    a_K = \frac{1}{K} \left( 1 - \frac{B_{1,1}}{B_{1,K}} \right) \quad {\rm and} \quad b_K = \frac{K}{B_{1,K}}.
\end{align}
Finally, by Lemma \ref{lem:li} it is easy to see that the ranges of $a_K$ and $b_K$ satisfy that
\begin{align*}
    \vert a_K \vert \le 20 \quad {\rm and }\quad \vert b_K \vert \le 10.
\end{align*}
% We can let 
% where $\vb^{(i)} = x_{2i} \ve_1 - x_{2i+1} \ve_K$. For $\vy^{(i)}$ satisfying Eq. (\ref{eq:close-yz-star}), we have
% \begin{align*}
%     \Vert \vz^{(i)} \Vert^2 =  \Vert \mB \vb^{(i)} \Vert^2,
% \end{align*}
% where $\mB = (\alpha \mI_K + \mA_K)^{-1}$ and $\alpha = 1/K^2$. Let $\mD = \mB^2$, then
% \begin{align*}
%     \Vert \vz^{(i)} \Vert^2  = \left \langle \vb^{(i)}, \mD \vb^{(i)} \right \rangle = D_{1,1} x_{2i}^2 - 2 D_{1,K} x_{2i} x_{2i+1} + D_{K,K} x_{2i+1}^2.
% \end{align*}
% Therefore, plugging in $\vy^*(\vx)$ and $\vz^*(\vx)$ given by Eq. (\ref{eq:close-yz-star}) to $h^{\rm reg}$ in Eq. (\ref{eq:our-f}) yields 
% \begin{align*}
%        % \varphi(\vx) &= \frac{\sqrt{\nu}}{2} ( \bar x_1 - 1)^2 +  \frac{1}{2}\sum_{i=1}^T  (\bar x_{2i-1} - \bar x_{2i})^2 + \nu \sum_{i=1}^{2T} \Upsilon_r (\bar x_i) \\
% h^{\rm reg}(y_{2i},y_{2i+1}, \vz^{(i)}) = (c_1 K^3 + c_2 D_{1,1} x_{2i}^2) - 2 c_2 D_{1,K} x_{2i} x_{2i+1} + (c_1 K^3 + c_2 D_{K,K}) x_{2i+1}^2,
% \end{align*}
% From Lemma \ref{lem:li} we know 
% \begin{align*}
%     D_{K,K} = D_{1,1} = \sum_{i=1}^K B_{1,i} B_{i,1} = \sum_{i=1}^K B_{i,1}^2 \in [0.01 K^3, 400 K^3] 
% \end{align*}
% and
% \begin{align*}
%     D_{1,K} = D_{K,1} = \sum_{i=1}^K B_{K,i} B_{i,1} =  \sum_{i=1}^K B_{1,K+-i} B_{i,1} \in [0.01K^3, 400K^3],
% \end{align*}
% where the fact $B_{K,i} = B_{1,K+1-i}$ can be verified by $\mB = \mM^\top / \det ( \alpha \mI_K + \mA_K )$ with $\mM$ is the cofactor matrix of $  \alpha \mI_K + \mA_K $. Now, Lemma \ref{lem:hyperfunc} is proved by letting 
% \begin{align} \label{eq:setting-c}
%     c_2 = \frac{K^3}{D_{1,K}}, \quad c_1 = 1 - \frac{D_{1,1}}{D_{1,K}}.
% \end{align}
% Finally,
% from the range of $D_{1,1}$ and $D_{1,K}$ it is easy to check that
% \begin{align*}
%     |c_2| \le 100, \quad \vert c_1 \vert < 4 \times 10^4,
% \end{align*}
% which means both $c_1$ and $c_2$ are (almost) numerical constants. This completes the proof.
\end{proof}

\section{Proof of Theorem \ref{thm:NC-SC}} \label{apx:proof-thm-SCNC}

\begin{proof}
From Lemma \ref{lem:property}, we know that the hard instance we construct has $\bar \ell_1$-Lipschitz continuous gradients if we set $\nu r^{3-p} \le 1$, where $\bar \ell_1 :=124 + \nu r^{3-p} \ell_1$. When $T \ge 2$, we let
$m =  2T+ K(T-1)$ and $n = T(K+1)-K$ be the dimensions of $\bar f^{\text{nc-sc}}, \bar g^{\rm sc}$ and effective lower bound dimension, respectively. We have $n \ge m$. Let $\vw = (\vx,\vy,\vz)$.
For any given $L_1 \ge \mu_y >0$, we define the following rescaled and rotated functions $f,g: \sR^{6n} \rightarrow \sR$ as
\begin{align} \label{eq:scaled-fg}
   f(\vw) = \frac{L_1 \beta^2}{\bar \ell_1} \bar f^{\text{nc-sc}} \left(\frac{\mP^\top \vw}{\beta} \right) \quad {\rm and} \quad g(\vw) = \frac{L_1 \beta^2}{\bar \ell_1} \bar g^{\rm sc} \left( \frac{\mP^\top \vw}{\beta} \right)
\end{align}
to ensure that both $f$ and $g$ have $L_1$-Lipschitz continuous gradients for any $\beta>0$, where the matrix $\mP \in \sR^{6 n} \times \sR^{m}$ satisfies $\mP^\top \mP = \mI_{m}$ and is defined as $\mP = \mU \mE$, where
%composited by three matrices $\mP_x \in \sR^{d \times T}$, $\mP_z \in \sR^{d \times T} $, and $\mP_z \in \sR^{d \times K(T-1)}$ via
\begin{align} \label{eq:dfn-matrix-P}
\mP = 
  \begin{bmatrix}
        \mP_x & & \\
        & \mP_z & \\
        & & \mP_y
    \end{bmatrix}
    =
    \begin{bmatrix}
        \mU_x \mE_x & & \\
        & \mU_z \mE_x & \\
        & & \mU_y \mE_y
    \end{bmatrix}
\end{align}
In the above, we let $\mE_x \in \sR^{n\times T}, \mE_y \in \sR^{n \times K(T-1)}$ be
\begin{align*}
 &\mE_x = (\ve_1, \ve_{K+2}, \ve_{2K+3}, \cdots, \ve_{n-(K+1)}, \ve_{n}) \\
 {\rm and}~~ & \mE_y = (\ve_2, \ve_3, \cdots, \ve_{K+1}, \ve_{K+3}, \ve_{K+4}, \cdots, \ve_{2(K+1)}, \ve_{2K+4}, \ve_{2K+5} \cdots, \ve_{T(K+1)}),
\end{align*}
respectively, and adversarially select $\mU_x, \mU_z, \mU_y \in \sR^{2n\times n}$ as three matrices with orthogonal columns $(\vu_x^0,\cdots,\vu_{x}^n)$, $(\vu_z^0,\cdots,\vu_{z}^n)$, and $(\vu_y^0,\cdots,\vu_{y}^n)$ according to the procedure below:
\begin{enumerate}
    \item Pick $(\vu_x^0,\vu_z^0,\vu_y^0)$ as three arbitrary unit vectors in $\sR^{2n}$. At the initialization $\vw^0 = (\vx^0,\vz^0,\vy^0) = (\vzero, \vzero, \vzero)$, we know from the zero-chain structure of Eq. (\ref{eq:our-f}) that the oracle response $\sO(\vw^0)$ only depends on $(\vu_x^0,\vu_z^0,\vu_y^0)$, which is already fixed. Since the algorithm is deterministic, we know the next query $\vw^1 = (\vx^1, \vy^1, \vz^1)$ is also fixed.
    \item For any $t =1, \cdots,n $, we repeat the following process: We compute $\vw^t = (\vx^t, \vz^t, \vy^t)$ from the algorithm, and pick $\vu_x^t, \vu_z^t$ and $ \vu_y^t$ as the unit vectors satisfying 
    \begin{align*}
    &\langle \vu_x^t, \vu_x^i \rangle = \langle \vu_z^t, \vu_z^i \rangle = \langle \vu_y^t, \vu_y^i \rangle = 0,\quad \forall i = 1,\cdots,t-1; \\
    {\rm and} ~~ &\langle \vu_x^t, \vx^i \rangle = \langle \vu_z^t, \vz^i \rangle =  \langle \vu_y^t, \vy^i \rangle = 0,\quad \forall i = 1,\cdots,t; 
    \end{align*}
    We know such vectors $\vu_x^t, \vu_z^t,$ and $\vu_y^t$ always exist in the space $\sR^{2n}$ since $ 2t-1 < 2n$.
    Assuming that the subsequent construction will always ensure that all the vectors $(\vu^{t+1}_x, \cdots, \vu^n_x)$, $(\vu^{t+1}_z, \cdots, \vu^n_z)$, and $(\vu^{t+1}_y, \cdots, \vu^n_y)$  are orthogonal to $\vx^t$, $\vz^t$, and $\vy^t$, respectively. we know from the zero-chain structure of Eq. (\ref{eq:our-f}) that the oracle response $\sO(\vw^t)$ only depends on $(\vu^0_x,\cdots, \vu^t_x)$, $(\vu^0_y,\cdots, \vu^t_y)$, and $(\vu^0_z,\cdots, \vu^t_z)$, which are already fixed. Since the algorithm is deterministic, it means that the next query $\vw^{t+1}$ is also fixed.
    % \item At the end of the procedure, we pick the remaining $3T$ vectors 
    % % \begin{align*}
    % % \end{align*}
\end{enumerate}
Let $n' = n - (K+1) = (T-1)(K+1) - K = \Omega(KT )$.
From the above procedure, any first-order deterministic algorithm with $t \le n'$ must satisfy $\langle \vu_x^n, \vx^t \rangle = \langle \vu_x^{n'}, \vx^{t} \rangle = 0$.
We also choose  $ K =  \left \lfloor \sqrt{{L_1}/{(\bar \ell_1 \mu_y)}} \right \rfloor$ 
to ensure that $g(\vx,\vz,\vy)$ is $\mu_y$-strongly convex jointly in $(\vz,\vy)$. If 
$\mu_y \lesssim L_1$ then we have $K \ge 10$, which satisfies the condition of Lemma \ref{lem:li}. Combining Lemma \ref{lem:hyperfunc}, we know the hyper-objective after our rescaling and rotation is
\begin{align} \label{eq:varphi-ncsc}
     F(\vx) &= \frac{L_1 \beta^2}{\bar \ell_1} \bar f^{\rm nc}_{\nu, r}(K^{3/2} \mP_x^\top \vx/ \beta), \quad {\rm where} ~~ \mP_x = \mU_x \mE_x.   %&=  \frac{ L_1 }{\bar \ell_1} \left( \frac{\sqrt{\nu}}{2} ( \hat x_1 - \beta)^2 + \frac{1}{2} \sum_{i=2}^{2T+1} (\hat x_{i} - \hat x_{i-1})^2 +  \nu \beta^2 \sum_{i=1}^{2T} \Upsilon_r (\hat x_i/ \beta) \right) \\
\end{align}
% The zero-chain property in Lemma \ref{lem:property} (item~4) together with 
Since $\langle \vu_x^n, \vx^t \rangle = \langle \vu_x^{n'}, \vx^{t} \rangle = 0$, Lemma~\ref{lem:Car-Upsion} (item 4) means that for any $t \le n'$ we have 
\begin{align} \label{eq:varphi-grad}
    \Vert \nabla F(\vx^t) \Vert = \frac{L_1 \beta K^{3/2}}{\bar \ell_1} \Vert \nabla f^{\rm nc}_{\nu, r}(K^{3/2} \mP_x^\top \vx^t / \beta) \Vert >  \frac{L_1 \beta K^{3/2} \nu^{3/4}}{4 \bar \ell_1} \ge \epsilon,
\end{align}
where the last inequality holds if we set 
\begin{align} \label{eq:para-beta}
     \beta = \frac{4 \bar \ell_1 \epsilon}{K^{3/2} \nu^{3/4} L_1}.
\end{align}
It implies an $\Omega(K T)$ lower bound for any first-order zero-respecting algorithm to find an $\epsilon$-stationary point of $F(\vx)$. To make the lower bound as large as possible, we choose $T$ to be the largest possible value that satisfies the constraint $F(\vzero) - \inf_{\vx \in \sR^{2n}} F(\vx) \le \Delta$. By Lemma~\ref{lem:Car-Upsion}, we know that $\Upsilon_r(0) \le 10$ and  $\Upsilon_r(x) > \Upsilon_r(1) =0$. Therefore, the minimizer of $\bar f^{\rm nc}_{\nu, r}(\vx)$ is obtained at $\vx  =\vone$ with $\bar f_{\nu, r}^{\rm nc}(\vone) = 0$, which further leads to
\begin{align} \label{eq:Delta}
    F(\vzero) - \inf_{\vx \in \sR^{2n}} F(\vx) = F(\vzero) \le \frac{L_1 \beta^2}{\bar \ell_1} \left( \frac{\sqrt{\nu}}{2} + 10 \nu (T-1) \right).
\end{align} 
It means that we can fulfill the constraint $F(\vzero) - \inf_{\vx \in \sR^{T}} F(\vx) \le \Delta$ by setting
\begin{align} \label{eq:para-T}
    T = \left \lfloor 
    \left( \frac{\bar \ell_1 \Delta}{L_1 \beta^2} - \frac{\sqrt{\nu}}{2} \right) \big / (10\nu)
    \right \rfloor + 1.
\end{align}
Now, it remains to determine the setting
g of $\nu$ and $r$, which we choose to satisfy the smoothness conditions under the following three different cases.
\begin{enumerate}
    \item \textbf{The case $p=1$.} We simply set $\nu = r = 1$ to obtain the lower bound 
\begin{align*}
    \Omega(K T) =  \Omega \left( \kappa_y^{2} L_1 \Delta \epsilon^{-2} \right).
\end{align*}
\item \textbf{The case $p=2$.} We let $r = 1$, then the condition $\nu r^{3-p} \le 1$ still holds if $\nu \le 1$, which means both $f$ and $g$ have $L_1$-Lipschitz continuous gradients. 
Next, we ensure that  $f$ has $L_2$-Lipschitz continuous Hessians by setting 
\begin{align} \label{eq:para-nu-p2}
    \nu = \frac{\beta \bar \ell_1 L_2}{L_1 \ell_2}.
\end{align}
Combining with the setting of $\beta$ in Eq. (\ref{eq:para-beta}), we can solve that
\begin{align*}
    \nu &= \left( {4 \epsilon}/{K^{3/2}}\right)^{4/7} \left( {\bar \ell_1}/{L_1} \right)^{8/7} \left({L_2}/{\ell_2} \right)^{4/7}, \\
    \beta &= \left( {4 \epsilon}/{K^{3/2}} \right)^{4/7} \left( {\bar \ell_1}/{L_1}  \right)^{1/7} \left( {\ell_2}/{L_2}  \right)^{3/7}.
\end{align*}
Under our parameter settings, and $\nu \le 1$  holds if
\begin{align*}
    \epsilon \lesssim \kappa_y^{3/4}\left( {L_1}/{\bar \ell_1} \right)^2 \left({\ell_2}/{L_2} \right).
\end{align*}
Finally, we substitute $\beta$ and $\nu$ into Eq. (\ref{eq:para-T}) to calculate $T$. We assume that
\begin{align} \label{eq:assume}
     \frac{\bar \ell_1 \Delta}{L_1 \beta^2} \gtrsim  \sqrt{\nu},
\end{align}
which implies that $T \ge 2$ and is true under the condition
\begin{align*}
    \epsilon \lesssim \kappa_y^{3/4} \Delta^{7/10} (\bar \ell_1/ L_1)^{1/10} (L_2/ \ell_2)^{2/5}. 
\end{align*}
Then we have
\begin{align*}
    T+1 \ge  
    \frac{\bar \ell_1 \Delta}{40 \nu L_1 \beta^2} = \frac{\ell_2 \Delta}{40 \beta^3 L_2  }  = \frac{\Delta K^{18/7}}{40 \cdot 4^{12/7}} ( L_1 / \bar \ell_1 )^{3/7} (L_2/ \ell_2)^{2/7} \epsilon^{-12/7},
\end{align*}
which indicates the lower bound of
\begin{align*}
    \Omega(K T) = \Omega ( \kappa_y^{25/14}  L_1^{3/7} L_2^{2/7} \Delta \epsilon^{-12/7}    ).
\end{align*}
\item \textbf{The case $p \ge 3$.} Let $r = \sqrt{1/\nu}$ then the condition $\nu r^{3-p} \le 1$ still holds if $\nu \le 1$, which means both $f$ and $g$ have $L_1$-Lipschitz continuous gradients.
For $q = 2,\cdots,p$, the $q$th-order derivative of $g$ is zero and 
the $q$th-order derivatives of $f$ is $ (L_1/\bar \ell_1) \beta^{1-q} \nu^{(q-1)/2} \ell_q$-Lipschitz continuous, where $\ell_q$ is the numerical constant defined in Lemma \ref{lem:Car-Upsion}.
Hence, we can ensure that the $q$th-order derivative of $f$ is $L_q$-Lipschitz continuous for all $q = 2,\cdots,p$ by setting 
\begin{align} \label{eq:para-nu-p3}
    \nu =  \beta^2  \min_{q=2, \cdots, p}  \left( \frac{\bar \ell_1 L_q}{L_1 \ell_q} \right)^{2/(q-1)} := \beta^2 L_*.
\end{align}
Combining with the setting of $\beta$ in Eq. (\ref{eq:para-beta}), we can solve that
\begin{align*}
    \beta = L_*^{-3/10}\left( \frac{4 \bar \ell_1 \epsilon}{K^{3/2} L_1
    } \right)^{2/5} \quad {\text{and}} \quad \nu = L_*^{2/5} \left( \frac{4 \bar \ell_1 \epsilon}{K^{3/2} L_1 } \right)^{4/5}.
\end{align*}
Under our parameter settings, $\nu \le 1$  holds if
\begin{align*}
    \epsilon \lesssim 
    L_*^{-1/2}\kappa_y^{3/4} ( L_1/{\bar \ell_1}).
\end{align*}
Finally, we substitute $\beta$ and $\nu$ into Eq. (\ref{eq:para-nu-p3}) to calculate $T$. As the case $p=2$, we also assume that Eq. (\ref{eq:assume}) holds, 
which is now true under the condition
\begin{align*}
    \epsilon \lesssim \kappa_y^{3/4} \Delta^{5/6} (L_1 / \bar \ell_1)^{1/6} L_*^{1/3}.
\end{align*}
Then we have
\begin{align*}
    T+1 \ge  
    \frac{\bar \ell_1 \Delta}{40 \nu L_1 \beta^2} = \frac{\bar \ell_1 \Delta}{40 L_1 \beta^4 L_*} = \frac{\Delta K^{12/5}}{40 \cdot 4^{8/5}} (L_1 / \bar \ell_1 )^{3/5} L_*^{1/5} \epsilon^{-8/5},
\end{align*}
which indicates the lower bound of
\begin{align*}
    \Omega(K T) = \Omega ( \kappa_y^{17/10} L_1^{3/5} L_*^{1/5} \Delta \epsilon^{-8/5} ).
\end{align*}
\end{enumerate}

\end{proof}

\section{Proof of Theorem \ref{thm:C-SC}} \label{apx:proof-thm-CSC}


% We set $\lambda = 0 $ in Eq. (\ref{eq:our-f-convex}), then the hyper-objective in Eq. (\ref{eq:varphi-sc}) becomes
% \begin{align*}
% F(\vx) =  \frac{L_1 \beta^2}{\bar \ell_1} \bar f^{\rm c}(K^{3/2} \vx / \beta),
% \end{align*}
% where $\bar \ell_1 = 124$ and $\bar f^{\rm c}: \sR^T \rightarrow \sR$ is
% the convex zero-chain in Definition \ref{dfn:finite-Nes-func}.

\begin{proof}
To give lower bounds for C-SC problems, let us consider the following instance, which modifies the rescaled and rotated NC upper-level function in Eq. (\ref{eq:scaled-fg}) to
\begin{align} \label{eq:our-f-convex}
    f(\vx,\vz,\vx) :=  \frac{L_1 \beta^2}{2\bar \ell_1} \left(
    \left( \frac{\langle \vp_z^1, \vz\rangle}{\beta \sqrt{K}} - 1 \right)^2 + h\left(\frac{\mP_z^\top \vz}{\beta},\frac{\mP_y^\top \vy}{\beta} \right) \right),
\end{align}
where $\mP = {\rm diag}(\mP_x, \mP_z, \mP_y)$ be the same column-orthogonal matrix defined in Eq. (\ref{eq:dfn-matrix-P}) and $\vp_z^1$ be the first column of $\mP_z$. Both the component $h: \sR^T \times \sR^{K(T-1)} \rightarrow \sR$ in Eq.~(\ref{eq:our-hard-h}) and the SC lower-level function $g(\vx,\vy,\vz) $ in Eq. (\ref{eq:scaled-fg}) remain unchanged as in the NC-SC case:
\begin{align*}
    g(\vx,\vy,\vz) := \frac{L_1 \beta^2}{\bar \ell_1} \bar g^{\rm sc} \left( \frac{\mP_x^\top \vx}{\beta}, \frac{\mP_z^\top \vz}{\beta}, \frac{\mP_y^\top \vy}{\beta} \right).
\end{align*}
Note that both
Lemma \ref{lem:property} and \ref{lem:hyperfunc} still apply to the above instance, indicating that both 
$f$ and~$g$ have $L_1$-Lipschitz continuous gradients for any $\beta>0$, and the hyper-objective $F$ is 
% % $\bar f^{\text{sc-sc}}$ has $\bar \ell_1 = 124$-Lipschitz continuous gradients and $\bar g^{\rm sc}$ has $5$-Lipschitz continuous gradients. Therefore, we can follow the proof of NC-SC setting to define the scaled functions $f$ and $g$ according to Eq. (\ref{eq:scaled-fg}), which ensures both $
% % Since Lemma \ref{}
% % Since , we can 
% % . We can still apply Lemma \ref{lem:hyperfunc} to obtain 
% \begin{align*}
%     \bar F^{\text{sc}}(\vx) := \bar f^{\text{nc-sc}}(\vx,\vz^*(\vx), \vy^*(\vx)) = 
% \end{align*}
\begin{align} \label{eq:varphi-sc}
      F(\vx) 
      %&= \frac{L_1}{\bar \ell_1} \left(\frac{\sqrt{\nu}}{2} ( \hat x_1 - \beta)^2 + \frac{1}{2} \sum_{i=2}^{2T+1} (\hat x_{i} - \hat x_{i-1})^2 \right) + \frac{L_1 \sigma }{2 \bar \ell_1} \Vert \vx \Vert^2 \\
   = \frac{L_1 \beta^2}{\bar \ell_1} \bar f^{\rm c} \left(\frac{K^{3/2} \mP_x^\top \vx}{\beta}\right),
\end{align}
where $\bar f^{\rm c}: \sR^T \rightarrow \sR$ is the convex zero-chain in Definition \ref{dfn:finite-Nes-func}. We can observe that
the minimizer of 
the hyper-objective $F(\vx): \sR^{2n} \rightarrow \sR$ is achieved at $\mP_x^\top \vx^* = \beta \vone / K^{3/2}$. Then, setting $\beta = {K^{3/2} D}/{\sqrt{T}}$ satisfies the constraint that
$\min\{ \Vert \vx \Vert : \vx \in  \arg \min_{\vx \in \sR^{2n}} F(\vx)  \} \le D$.
As in the proof of NC-SC lower bound, our construction of $\mP = {\rm diag}(\mP_x, \mP_z, \mP_y)$ ensures that 
any first-order deterministic algorithm with $t \le n$ must satisfy $\langle \vu_x^n, \vx^t \rangle  = 0$, where $n = T (K+1) - K = \Omega(TK)$. Hence, by Lemma \ref{lem:Car-large-grad-convex}, we have 
%From Lemma \ref{lem:Car-large-grad-convex} and item 4 of Lemma~\ref{lem:property}, we know that for any $t \le K(T-1)$, we have
\begin{align*}
    \Vert \nabla F(\vx^t) \Vert = \frac{L_1 \beta K^{3/2} }{\bar \ell_1} \Vert \nabla \bar f^{\rm c}_1(K^{3/2} \mP_x^\top \vx_t/\beta) \Vert > \frac{K^{3}L_1 D}{\bar \ell_1 T^2}.
\end{align*}
Setting $K = \Theta(\sqrt{\kappa_y})$ and $T = \Theta( \sqrt{K^{3} L_1 D / \epsilon} )$ proves the lower bound of
\begin{align*}
    \Omega (K T) = \Omega\left( \kappa_y^{5/4} \sqrt{ L_1  D/ \epsilon} \right).
\end{align*}
\end{proof}

\section{Proof of Theorem \ref{thm:SC-SC}}  \label{apx:proof-thm-SCSC}
\begin{proof}
To give lower bounds for SC-SC problems, we add the quadratic regularizer $\mu_x \Vert \vx \Vert^2/2$ to the upper-level function $f$ of C-SC problems in Eq. (\ref{eq:our-f-convex}), that is,
\begin{align*} 
    f(\vx,\vz,\vx) :=  \frac{L_1 \beta^2}{2\bar \ell_1} \left(
    \left( \frac{\langle \vp_z^1, \vz\rangle}{\beta \sqrt{K}} - 1 \right)^2 + h\left(\frac{\mP_z^\top \vz}{\beta},\frac{\mP_y^\top \vy}{\beta} \right) \right)
    + \frac{\mu_x}{2} \Vert \vx \Vert^2,
\end{align*}
and keep the lower-level function $g$ unchanged. Then, the hyper-objective in Eq. (\ref{eq:varphi-sc}) is
\begin{align*} 
      F(\vx) 
      %&= \frac{L_1}{\bar \ell_1} \left(\frac{\sqrt{\nu}}{2} ( \hat x_1 - \beta)^2 + \frac{1}{2} \sum_{i=2}^{2T+1} (\hat x_{i} - \hat x_{i-1})^2 \right) + \frac{L_1 \sigma }{2 \bar \ell_1} \Vert \vx \Vert^2 \\
   = \frac{L_1 \beta^2}{\bar \ell_1} \bar f^{\rm c} \left(\frac{K^{3/2} \mP_x^\top \vx}{\beta}\right) + \frac{\mu_x}{2} \Vert \vx \Vert^2.
\end{align*}
Recall $\mA_{T}$ is the matrix defined in Eq. (\ref{eq:finite-A}) with $K = T$. We solve the minimizer  $\vx^* = \arg \min_{\vx \in \sR^{2n}} F(\vx)$ by taking gradient in Eq. (\ref{eq:varphi-sc}) and set it to be zero, yielding that
\begin{align} \label{eq:first-order-equation}
    \left(\mP_x \left( \mA_{T} + \ve_1 \ve_1^\top \right) K^{3/2} \mP_x^\top + \alpha \mI_{2n} \right) \vx^* = \beta \ve_1, \quad {\rm where} \quad \alpha = \frac{\bar \ell_1 \mu_x}{L_1 K^3}.
\end{align} 
Recall $\mP_x \in \sR^{2 n \times T}$ is a matrix with $T$ orthogonal columns $(\vp_x^1, \cdots, \vp_x^T)$. Since $T \le 2 n$, according to the extension theorem, it can be extended to a basic $(\vp_x^1, \cdots, \vp_x^T, \vp_x^{T+1}, \cdots, \vp_x^{2n})$. Let $\tilde \mP_x \in \sR^{2n \times 2n} $ be the matrix whose columns consist of the basis.
% As in \citep{nesterov2018lectures}, we set  $n = T = \infty$. 
Define the rotated minimizer $\tilde \vx^* = \tilde \mP_x^\top \vx^*$
Then, expanding Eq. (\ref{eq:first-order-equation}) gives 
\begin{align} \label{eq:three-term}
    \begin{cases}
        -\beta / K^{3/2} + (2 + \alpha)  \tilde x_1^* - \tilde x_2^* = 0 \\
        - \tilde x_1^* + (2+ \alpha) \tilde x_2^* - \tilde x_3^* = 0 \\
        - \tilde x_2^* + (2+ \alpha) \tilde x_3^* - \tilde x_4^* = 0  \\
        \quad \quad \quad \quad \vdots \\
        % % -\bar x_{2T-1}^* + (2+ \mu_x) \bar x_{2T}^* - \bar x_{2T+1}^* = 0 \\
        - \tilde x_{T-1}^* + (2+ \alpha) \tilde x_{T}^* = 0,
    \end{cases}
\end{align}
and $\tilde x_{i}^* =0 $ for all the remaining $i = T+1,\cdots, 2n$.

Let $q \in (0,1)$ be the smallest root of the quadratic equation $-1+(2+\alpha) q -q^2 = 0$, \textit{i.e.},
\begin{align*}
    q = \frac{2 + \alpha - \sqrt{\alpha^2 + 4\alpha }}{2}.
\end{align*}
We can observe that the variable $\hat \vx^* \in \sR^{T}$ defined by 
$\hat x_i^* = \beta q^i / K^{3/2}$ satisfies Eq. (\ref{eq:three-term}) except the last line with a residual of $\beta q^{T+1} / K^{3/2} $, which means 
\begin{align*}
    \left( \mA_T + \ve_1 \ve_1^\top + \alpha \mI_{T} \right) K^{3/2} \hat \vx^* = \beta q^{T+1} \ve_{T} / K^{3/2} 
\end{align*}
Compared with Eq. (\ref{eq:first-order-equation}), we know that
\begin{align*}
    \hat \vx^* - \tilde \vx^* = \left( \mA_T + \ve_1 \ve_1^\top + \alpha \mI_{T} \right)^{-1} \beta q^{T+1} \ve_{T} / K^{3}.
\end{align*}
By the inequality $\left( \mA_T + \ve_1 \ve_1^\top + \alpha \mI_{T} \right)^{-1} \succeq \alpha^{-1}\mI_T  $, taking norm on both sides above gives
\begin{align} \label{eq:error-x-star}
    \Vert \hat \vx^* - \tilde \vx^* \Vert \le \beta q^{T+1} / (K^{3} \alpha).
\end{align}
If $\alpha \le 1/2$ and $T \ge 4$. We let $T$ also satisfies  
\begin{align} \label{eq:cond-T}
    T \ge  2  \ln \left( 4 \left( 1 + \frac{1-q}{\alpha K^{3/2}} \right) \right)     \Big / \ln (q^{-1}).
\end{align}
Then we have
\begin{align*}
    &\Vert \vx^* \Vert = \Vert \tilde \vx^* \Vert \le \Vert \hat \vx^* \Vert + \Vert \hat \vx^* - \tilde \vx^* \Vert \\
    \le& \frac{\beta q (1-q^T)}{(1-q) K^{3/2} } + \frac{\beta q^{T+1}}{\alpha K^{3}} \le \frac{2 \beta q}{(1-q) K^{3/2}},
\end{align*}
where the second inequality uses $\hat x_i^* = \beta q^i / K^{3/2}$ and Eq. (\ref{eq:error-x-star}), and the last one uses Eq.~(\ref{eq:cond-T}) that $q^T \le \alpha/ ((1-q) K^{3/2})$. From the above inequality, we know that
% Then $\tilde \vx^*$ satisfies $\tilde x_i^* = \beta q^i / K^{3/2}$ for all $i=1,2,3,\cdots$, which means that
letting $ \beta = K^{3/2} (1-q) D / (2q) $ satisfies the constraint 
$\min\{ \Vert \vx \Vert : \vx \in  \arg \min_{\vx \in \sR^{2n}} F(\vx)  \} \le D$.
% Similarly, we have
% \begin{align*}
%     &\Vert \vx^* \Vert = \Vert \tilde \vx^* \Vert \ge \Vert \hat \vx^* \Vert + \Vert \hat \vx^* - \tilde \vx^* \Vert  \\
%     \ge & \frac{\beta q (1-q^T)}{(1-q) K^{3/2} } + \frac{\beta q^{T+1}}{\alpha K^{3/2}} \ge \frac{\beta q}{2(1-q) K^{3/2}}
% \end{align*}
As in the proof of NC-SC lower bound, our construction of $\mP = {\rm diag}(\mP_x, \mP_z, \mP_y)$ ensures that 
any first-order deterministic algorithm with $t \le n'$ must satisfy
\begin{align*}
    \tilde \vx^t_i  = 0, \quad \forall i = \lfloor T/2 \rfloor +1, \cdots, T,
\end{align*}
where $n' = \lfloor T/ 2 \rfloor(K+1) -K = \Omega(K T )$. This further gives the following lower bound of the hyper-gradient norm:
\begin{align*}
    & \Vert \nabla F(\mP_x^\top \vx^t) \Vert \ge \mu_x \Vert \mP_x^\top \vx^t -\vx^* \Vert \ge \mu_x \Vert \mP_x^\top \vx^t -\tilde \vx^* \Vert \\
    \ge & \mu_x \left( \Vert \mP_x^\top \vx^t -\tilde \vx^* \Vert - \Vert \tilde \vx^* - \hat \vx^* \Vert \right) \\
    \ge & \mu_x \left(\frac{\beta q^{\lfloor T/2 \rfloor +1} (1 - q^{T - \lfloor T/2 \rfloor})}{(1-q) K^{3/2}} -  \frac{\beta q^{T+1}}{\alpha K^{3}} \right) \\
    = & \frac{\mu_x D q^{\lfloor T/2 \rfloor }}{2}   \left(1 -  \left( 1 + \frac{1-q}{\alpha K^{3/2}} \right) q^{T - \lfloor T/2 \rfloor}  \right) \ge  \frac{\mu_x D q^{\lfloor T/2 \rfloor }}{8}.
\end{align*}
In the above, the first inequality uses the $\mu_x$-strong convexity of $F(\vx)$, the second last inequality uses $\hat x_i^* = \beta q^i / K^{3/2}$ and Eq. (\ref{eq:error-x-star}), and the last one uses the condition of $T$ in Eq. (\ref{eq:cond-T}).
% % $\langle \vu_x^n, \vx^t \rangle  = 0$, where $n = T (K+1) - K = \Omega(TK)$. Therefore, we have 
% % Since item 4 of Lemma \ref{lem:property} still holds by the same chain structure, for any $t \le K (T-1)$, we have $x_j^t = 0 $ for all $j \ge T$ and
% \begin{align*}
%     \Vert \nabla F(\mP_x^\top \vx^t) \Vert \ge  \mu_x \Vert \vx^t -\vx^* \Vert \ge   \mu_x \sqrt{\sum_{i=T}^{\infty} (x_i^*)^2} \ge  \mu_x q^T \Vert \vx^0 - \vx^* \Vert,
% \end{align*}
To let $\Vert \nabla F(\vx^t) \Vert > \epsilon$ we require that 
\begin{align} \label{eq:cond-T-eps}
    T \le 2\ln \left( \frac{\mu_x D}{8 \epsilon}  \right) \Big / \ln (q^{-1}).    
    % = \Omega \left( \kappa_y^{3/4} \sqrt{\kappa_x}  \ln \left( \frac{\mu_x \Vert \vx^0 - \vx^* \Vert^2}{\epsilon}  \right) \right),
\end{align}
Recall that $K = \Theta(\sqrt{\kappa_y})$ and $\alpha^{-1} = \Theta\left(\kappa_y^{3/2} \kappa_x \right)$. We can derive that
\begin{align} \label{eq:lb-q}
    \ln (q^{-1}) \ge \frac{q}{1-q} = \frac{2 + \alpha - \sqrt{\alpha^2 + 4 \alpha}}{\sqrt{\alpha^2 + 4 \alpha} - \alpha} = \frac{2}{\sqrt{\alpha^2 + 4\alpha} - \alpha} -1 \ge \sqrt{\frac{1}{\alpha}} - 1,
\end{align}
where the last inequality uses $\sqrt{a+b} \le \sqrt{a} + \sqrt{b}$.
If the accuracy $\epsilon$ satisfies that $\epsilon \lesssim \alpha \mu_x D$, or equivalent to $\epsilon \lesssim \mu_x D / (\kappa_x \kappa_y)$, then we know there exists $T$ that jointly fulfills 
 Eq. (\ref{eq:cond-T}) and~(\ref{eq:cond-T-eps}). Now,
we let $T$ be the smallest value that fulfills Eq. (\ref{eq:cond-T-eps}), and by Eq. (\ref{eq:lb-q}), we have 
% \begin{align*}
%      \asymp 
$T = \Omega( \kappa_y^{3/4} \sqrt{\kappa_x}  \ln (\mu_x D/\epsilon))$, which implies the final lower bound of 
\begin{align*}
    \Omega(KT)  = \Omega \left( \kappa_y^{5/4} 
\sqrt{\kappa_x}  \ln (\mu_x D/\epsilon)  \right).
\end{align*}

\end{proof}

\section{Proof of Theorem \ref{thm:NC-SC-stoc}} \label{apx:proof-thm-NCSC-stoc}

\begin{proof}
From Eq. (\ref{eq:our-fg-stoc}) it is clear that $(f,g) \in  \gF^{\text{nc-scq}}(L_1,\mu_y, \Delta)$.
By  Definition~\ref{dfn:rand-alg}, a randomized algorithm $\texttt{A}$ consists of the mappings $\{\texttt{A}^t \}_{t \in \sN}$ recursively defined~as
\begin{align*}
    (\vx^t,\vz^t) = \texttt{A}^t(r, \hat \nabla \bar f^{\text{nc-rs}}(\vz^0), \nabla g(\vx^0,\vz^0), \cdots, \hat \nabla \bar f^{\text{nc-rs}}(\vz^{t-1}), \nabla g(\vx^{t-1},\vz^{t-1})   )
\end{align*}
Note that $\nabla g(\vx,\vz)$ is a deterministic mapping of $(\vx,\vz)$. We know that there exists mappings $\{\texttt{A}_1^t \}_{t \in \sN}$ that satisfy the recursion
\begin{align*}
    (\vx^t,\vz^t) = \texttt{A}_1^t(r, \hat \nabla \bar f^{\text{nc-rs}}(\vz^0), \vx_0,\vz_0, \cdots, \hat \nabla \bar f^{\text{nc-rs}}(\vz^{t-1}), \vx^{t-1}, \vz^{t-1}   )
\end{align*}
Expanding the recursion for $(\vx_t,\vz_t)$, we can observe that the above fact implies some mappings  $\{\texttt{A}_2^t \}_{t \in \sN}$ satisfying 
\begin{align*}
     (\vx^t,\vz^t) = \texttt{A}_2^t(r, \hat \nabla \bar f^{\text{nc-rs}}(\vz^0), \cdots, \hat \nabla \bar f^{\text{nc-rs}}(\vz^{t-1}) ).
\end{align*}
Then we can invoke Lemma \ref{lem:Arj-random} (item 4) to conclude that for all $t \le (T - \log 4 ) /(2p)$ and any deterministic mapping $\gM : \sR^T \rightarrow \sR^T$, with probability at least $1/2$, we have $\Vert \nabla \bar f^{\text{nc-rs}}(\gM(\vx^t)) \Vert \ge 1/2$.  By observation, we have $F(\vx) = (L_1 \beta^2 / 155) \bar f^{\text{nc-rs}}(\kappa_y \vx / \beta) $. Thus,
\begin{align*}
    \Vert \nabla F(\vx^t) \Vert = \frac{L_1 \beta \kappa_y}{155} \Vert \nabla \bar f^{\text{nc-rs}}( \kappa_y \vx^t / \beta)\Vert \ge \frac{L_1 \beta \kappa_y}{310},
\end{align*}
which means we can set
\begin{align} \label{eq:para-beta-stoc}
    \beta = \frac{310 \epsilon}{L_1 \kappa_y}
\end{align}
to obtain an $\Omega( T/p )$ lower bound for finding an $\epsilon$-stationary point of $F(\vx)$. Next, by Lemma~\ref{lem:Arj-random} (item 1), we have that
\begin{align*}
    F(\vzero) - \inf_{\vx \in \sR^T} F(\vx) \le \frac{12 L_1 \beta^2 T}{155},
\end{align*}
which means that we can fulfill the constraint $F(\vzero) - \inf_{\vx \in \sR^T} F(\vx) \le \Delta$ by setting
\begin{align} \label{eq:para-T-stoc}
      T = \left \lfloor 
     \frac{155 \Delta}{12 L_1 \beta^2} 
    \right \rfloor.
\end{align}
By Lemma \ref{lem:Arj-random} (item 3), the variance of the SFO for $(f,g)$ is upper bounded by $ (23 L_1 \beta / 155)^2$ $(1-p)/p $, which means that we can let the variance be bounded by $\sigma^2$ by setting
\begin{align} \label{eq:para-p-stoc}
    p = \min \left\{  1, \left( \frac{23 L_1 \beta}{\sigma 155} \right)^2 \right\}.
\end{align}
Finally, we can conclude the lower bound of 
\begin{align*}
    \Omega(T/p) = \Omega \left( \frac{\Delta}{L_1 \beta^2} \left(1+ \frac{\sigma^2}{L_1^2 \beta^2} \right)  \right) = \Omega \left( \frac{\kappa_y^2 L_1 \Delta }{\epsilon^2} \left( 1+ \frac{\sigma^2 \kappa_y^2}{\epsilon^2} \right)\right).
\end{align*}
\end{proof}


\end{document}